\documentclass[11pt,a4paper]{report}
\usepackage[utf8]{inputenc}
\usepackage[T1]{fontenc}
\usepackage[french]{babel}

%-----DIAGRAMMES (TikZ)-------------
\usepackage{tikz}
\usetikzlibrary{arrows.meta, positioning, shapes.geometric, calc,
                decorations.pathreplacing, patterns, decorations.markings}

% --- PACKAGES DE MATHÉMATIQUES ---
\usepackage{amsmath, amssymb, amsthm, mathrsfs}
\usepackage{bm}      % vecteurs gras italiques : \bm{U}

% --- CONFIGURATION DES LIENS CLIQUABLES ---
\usepackage{hyperref}
\hypersetup{
    colorlinks=true,
    linkcolor=blue,
    citecolor=green,
    urlcolor=cyan
}

% --- ENCADRÉS ARRONDIS PERSONNALISÉS (tcolorbox) ---
\usepackage[most]{tcolorbox}

\tcbset{
    commonbox/.style={
        enhanced,
        arc=4mm,
        boxrule=1.5pt,
        colback=white,
        fonttitle=\bfseries\sffamily,
        attach boxed title to top left={yshift=-2mm, xshift=5mm},
        boxed title style={sharp corners=downhill, arc=1.5mm},
        top=4mm, bottom=3mm, left=3mm, right=3mm
    }
}

% Encadré bleu --- relations importantes
\newtcolorbox{importantbox}[1]{
    commonbox,
    colframe=blue!70!black,
    colbacktitle=blue!70!black,
    coltitle=white,
    title=#1
}

% Encadré rouge --- définitions
\newtcolorbox{defbox}[1]{
    commonbox,
    colframe=red!75!black,
    colbacktitle=red!75!black,
    coltitle=white,
    title=#1
}

% --- MISE EN PAGE ---
\usepackage{geometry}
\geometry{hmargin=2.5cm,vmargin=3cm}

% --- TABLEAUX ET GRAPHIQUES ---
\usepackage{pgfplots}
\pgfplotsset{compat=1.18}
\usepackage{booktabs}

% --- COMMANDES PERSONNALISÉES ---
\newcommand{\DD}[1]{\dfrac{D #1}{Dt}}
\newcommand{\pd}[2]{\dfrac{\partial #1}{\partial #2}}
\newcommand{\ten}[1]{\overline{\overline{#1}}}   % tenseur d'ordre 2

% ==========================================
\begin{document}
% ==========================================

\title{\Huge\bfseries Mécanique des Fluides}
\author{\Large Nael LEPETRE \\ \small Centrale Lyon --- LMFA}
\date{\the\year}
\maketitle

\tableofcontents
\newpage

% ==========================================
\chapter{Cinématique, lois fondamentales et modèle non visqueux}
% ==========================================

Un Rafale Marine décolle d'un porte-avions à 180~km/h, 9~tonnes sous les ailes.
La question naïve---mais fondamentale---est : \textit{pourquoi cet engin ne s'écrase
pas dans l'océan~?}

La réponse tient en un mot : l'air. Plus précisément, les forces que l'air exerce
sur l'aile. Et pour calculer ces forces, il faut d'abord disposer d'un modèle
mathématique décrivant comment l'air s'écoule. C'est exactement l'objet de ce
premier chapitre.

On va construire brique par brique les outils de la cinématique des fluides, puis
on posera les lois de conservation (masse, quantité de mouvement), et on finira
par dériver l'équation de Bernoulli---qui explique pourquoi le Rafale vole---et son
application directe au tube de Pitot qui mesure la vitesse de l'avion.

\begin{defbox}{Fil conducteur}
Le Rafale en vol subsonique ($M_a \approx 0{,}2$ à l'approche, $M_a \leq 1{,}8$
en pointe) est notre système de référence tout au long de ce chapitre. On y reviendra
à chaque étape pour ancrer les concepts abstraits dans une réalité physique concrète.
\end{defbox}

% ------------------------------------------
\section{Décrire le mouvement d'un fluide}
% ------------------------------------------

\subsection{L'hypothèse du milieu continu}

L'air est un gaz composé de molécules qui se baladent en tous sens, se cognent
les unes aux autres, et font leur vie à l'échelle nanométrique. Si on prend au
sérieux cette description microscopique, le champ de vitesse de l'air autour du
Rafale est un chaos discontinu, inutilisable pour calculer quoi que ce soit de
macroscopique.

La bonne nouvelle : on n'a pas besoin de ce niveau de détail. Dès que l'échelle
caractéristique de l'écoulement $L$ est très grande devant le libre parcours moyen
des molécules $\lambda$, les comportements moléculaires se moyennent et le fluide
se comporte comme un milieu \textit{continu}---avec des champs de pression, de
densité et de vitesse bien définis en tout point.

\begin{defbox}{Hypothèse du milieu continu --- Nombre de Knudsen}
On caractérise l'échelle relative de la structure moléculaire par le
\textbf{nombre de Knudsen} :
\[
    Kn \equiv \frac{\lambda}{L},
\]
où $\lambda$ est le libre parcours moyen moléculaire et $L$ l'échelle caractéristique
de l'écoulement. L'hypothèse du milieu continu est valide dès lors que
\[
    Kn \ll 1.
\]
\end{defbox}

Pour le Rafale à basse altitude : $\lambda \sim 70~\text{nm}$ et $L \sim 1~\text{m}$
(envergure), soit $Kn \sim 7 \times 10^{-8}$. On est parfaitement dans le régime
continu. L'hypothèse ne lâche qu'à très haute altitude (re-entrée atmosphérique,
vol hypersonique) ou pour des objets microscopiques (MEMS, microfluidique)---ce n'est
pas notre problème ici.

Une seconde hypothèse de travail sera souvent faite dans ce chapitre : celle
d'\textbf{incompressibilité}. Un gaz est compressible par nature, mais si les vitesses
en jeu sont petites devant la vitesse du son, les variations de densité induites par
les variations de pression sont négligeables.

\begin{defbox}{Nombre de Mach et condition d'incompressibilité}
Le \textbf{nombre de Mach} compare la vitesse de l'écoulement à la vitesse du son $c$
dans le milieu :
\[
    M_a \equiv \frac{U}{c}.
\]
Un gaz peut être traité comme \textbf{incompressible} ($\rho \approx \text{cst}$)
dès lors que
\[
    M_a \lesssim 0{,}3.
\]
\end{defbox}

Le Rafale en phase d'approche vole à $\sim 250~\text{km/h} \approx 70~\text{m/s}$.
La vitesse du son à $20\,°C$ est $c \approx 340~\text{m/s}$, donc $M_a \approx 0{,}2$ :
l'hypothèse d'incompressibilité est valide. Ce premier chapitre se place donc dans ce
cadre. On y reviendra au chapitre~7 pour le vol supersonique (Mach~1{,}8).

\subsection{Quantités macroscopiques : les champs de l'écoulement}

Dans le cadre du milieu continu, on remplace la description moléculaire par des
\textbf{champs continus}. Les deux quantités fondamentales qui nous intéressent sont :

\begin{defbox}{Champs macroscopiques}
\begin{itemize}
    \item La \textbf{masse volumique} $\rho(\bm{x},t)$ (en kg/m$^3$) : masse par unité
    de volume au point $\bm{x}$ à l'instant $t$.
    \item Le \textbf{champ de vitesse} $\bm{U}(\bm{x},t)$ (en m/s) : vitesse
    macroscopique du fluide au point $\bm{x}$ à l'instant $t$.
\end{itemize}
Ces champs sont définis en \textit{tout} point de l'espace occupé par le fluide, à
tout instant.
\end{defbox}

Ce point de vue---fixer un point dans l'espace et observer ce qui s'y passe au fil
du temps---est la description \textbf{eulérienne}. C'est la description naturelle en
expérience : un tube de Pitot est installé en un point fixe sur le Rafale et mesure
$\bm{U}$ en ce point.

L'autre description, dite \textbf{lagrangienne}, consiste à suivre chaque particule
fluide individuellement dans son mouvement. Les deux approches sont équivalentes ; on
passera de l'une à l'autre selon la commodité mathématique.

\subsection{Domaine matériel et particule fluide}

Imaginons qu'on colorie en rouge un petit volume d'air juste en amont de l'aile du
Rafale. Ce volume se déplace et se déforme en suivant l'écoulement : il forme ce
qu'on appelle un \textbf{domaine matériel}.

\begin{defbox}{Domaine matériel}
Un \textbf{domaine matériel} $\mathcal{D}(t)$, délimité par une surface de contrôle
$\mathcal{S}(t)$, est un volume se déplaçant et se déformant avec le fluide. Sa
surface se déplace à la vitesse locale $\bm{U}_\mathcal{S} = \bm{U}$.

Un domaine matériel est toujours constitué des mêmes particules fluides : sa masse
est conservée.
\end{defbox}

La \textbf{particule fluide} est la limite infinitésimale de ce concept : un volume
élémentaire $d\mathcal{V}$ qu'on suit dans son mouvement. Elle a une masse
$dm = \rho\,d\mathcal{V}$ et une vitesse $\bm{U}$. C'est sur elle qu'on appliquera
le principe fondamental de la dynamique dans la suite.

\subsection{Lignes de courant et trajectoires}

Pour visualiser un écoulement, on dispose de deux outils complémentaires :

\begin{defbox}{Lignes de courant}
À un instant $t$ fixé, les \textbf{lignes de courant} sont les courbes partout
tangentes au champ de vitesse $\bm{U}(\bm{x},t)$. Ce sont des photographies
instantanées de la structure de l'écoulement.
\end{defbox}

\begin{defbox}{Trajectoires}
La \textbf{trajectoire} d'une particule fluide est le chemin parcouru au cours du
temps. Si $\bm{x}(t)$ désigne la position de la particule, elle vérifie :
\[
    \frac{d\bm{x}}{dt} = \bm{U}\!\left(\bm{x}(t),\, t\right).
\]
\end{defbox}

Pour un \textbf{écoulement stationnaire} ($\bm{U}$ indépendant de $t$), lignes de
courant et trajectoires coïncident. Pour un écoulement instationnaire (comme le
sillage tourbillonnaire derrière le Rafale en virage serré), les deux divergent.

\subsection{La dérivée matérielle : voir le monde depuis la particule}

On veut quantifier comment une grandeur physique $f(\bm{x},t)$ (température,
pression, norme de la vitesse\ldots) évolue \textit{vue depuis une particule fluide}
en mouvement.

La particule est en $\bm{x}(t)$ à l'instant $t$. Un instant $dt$ plus tard, elle
est en $\bm{x}(t+dt) = \bm{x}(t) + \bm{U}\,dt$. La variation de $f$ le long de
sa trajectoire vaut donc :
\[
    df = f\!\left(\bm{x} + \bm{U}\,dt,\; t + dt\right) - f(\bm{x},t)
       = \underbrace{\pd{f}{t}\,dt}_{\text{variation temporelle locale}}
       + \underbrace{(\bm{U}\cdot\boldsymbol{\nabla} f)\,dt}_{\text{variation convective}}.
\]

En divisant par $dt$, on obtient l'opérateur qui traduit ce ``regard en suivant
la particule'' :

\begin{importantbox}{Dérivée matérielle}
\[
    \frac{Df}{Dt} \;\equiv\; \pd{f}{t} + \bm{U}\cdot\boldsymbol{\nabla}f
    \;=\; \pd{f}{t} + U_j\,\pd{f}{x_j}.
\]
La dérivée matérielle $D/Dt$ donne le taux de variation d'une grandeur en suivant
une particule fluide. Elle se décompose en :
\begin{itemize}
    \item un terme \textbf{instationnaire} $\partial f/\partial t$ : variation due
    à l'évolution temporelle du champ en un point fixe,
    \item un terme \textbf{convectif} $\bm{U}\cdot\boldsymbol{\nabla}f$ : variation
    due au déplacement de la particule dans un champ spatialement non uniforme.
\end{itemize}
\end{importantbox}

Appliquée au champ de vitesse lui-même, la dérivée matérielle donne
l'\textbf{accélération} d'une particule fluide, qui vaut composante par composante :
\[
    \left.\frac{D\bm{U}}{Dt}\right|_i = \pd{U_i}{t} + U_j\,\pd{U_i}{x_j}.
\]
C'est ce terme qui joue le rôle du ``$\bm{a}$'' dans le principe fondamental de la
dynamique.

\subsection{Le théorème de Reynolds : différentier une intégrale sur domaine mobile}

On voudra souvent écrire la dérivée temporelle d'une grandeur intégrée sur un
domaine matériel en mouvement---la masse totale, la quantité de mouvement, etc. Le
problème est que les bornes d'intégration se déplacent avec le fluide. Le théorème
de Reynolds gère exactement cette difficulté.

La dérivée temporelle d'une intégrale de volume se décompose en deux contributions :
la variation de la grandeur à l'intérieur du domaine, et le flux de la grandeur à
travers la frontière qui se déplace à la vitesse $\bm{U}_\mathcal{S}$ :

\begin{importantbox}{Théorème de Reynolds}
Pour tout domaine $\mathcal{D}(t)$ de surface $\mathcal{S}(t)$ se déplaçant à la
vitesse de surface $\bm{U}_\mathcal{S}$, et pour toute grandeur scalaire
$\varphi(\bm{x},t)$ :
\[
    \frac{d}{dt}\int_{\mathcal{D}} \varphi\,d\mathcal{V}
    = \int_{\mathcal{D}} \pd{\varphi}{t}\,d\mathcal{V}
    + \int_{\mathcal{S}} \varphi\,\bm{U}_\mathcal{S}\cdot\bm{n}\,ds.
\]
Pour une grandeur \emph{massique} $\chi$ (i.e.\ $\varphi = \rho\chi$), on montre
que :
\[
    \frac{d}{dt}\int_{\mathcal{D}} \rho\chi\,d\mathcal{V}
    = \int_{\mathcal{D}} \rho\,\frac{D\chi}{Dt}\,d\mathcal{V}
    + \int_{\mathcal{S}} \rho\chi(\bm{U}_\mathcal{S}-\bm{U})\cdot\bm{n}\,ds.
\]
\end{importantbox}

Deux cas particuliers reviennent constamment :

\begin{itemize}
    \item \textbf{Domaine matériel} ($\bm{U}_\mathcal{S} = \bm{U}$) : le flux
    surfacique s'annule et $\dfrac{d}{dt}\displaystyle\int_{\mathcal{D}} \rho\chi\,d\mathcal{V}
    = \displaystyle\int_{\mathcal{D}} \rho\,\dfrac{D\chi}{Dt}\,d\mathcal{V}$.

    \item \textbf{Domaine fixe} ($\bm{U}_\mathcal{S} = \bm{0}$) :
    $\dfrac{d}{dt}\displaystyle\int_{\mathcal{D}} \rho\chi\,d\mathcal{V}
    = \displaystyle\int_{\mathcal{D}} \rho\,\dfrac{D\chi}{Dt}\,d\mathcal{V}
    - \displaystyle\int_{\mathcal{S}} \rho\chi\,\bm{U}\cdot\bm{n}\,ds$.
\end{itemize}

% ------------------------------------------
\section{Conservation de la masse}
% ------------------------------------------

\subsection{L'équation de continuité}

La masse d'un domaine matériel est constante---les mêmes particules sont toujours
là. Dit mathématiquement :
\[
    \frac{d}{dt}\int_{\mathcal{D}} \rho\,d\mathcal{V} = 0.
\]

On applique le théorème de Reynolds avec $\chi = 1$ sur un domaine matériel
($\bm{U}_\mathcal{S} = \bm{U}$, terme surfacique nul) :
\[
    \int_{\mathcal{D}} \frac{D\rho}{Dt}\,d\mathcal{V} = 0.
\]

On développe la dérivée matérielle :
$D\rho/Dt = \partial_t\rho + \bm{U}\cdot\boldsymbol{\nabla}\rho$, et on utilise
l'identité vectorielle
$\boldsymbol{\nabla}\cdot(\rho\bm{U}) = \rho\,\boldsymbol{\nabla}\cdot\bm{U}
+ \bm{U}\cdot\boldsymbol{\nabla}\rho$
pour regrouper :
\[
    \int_{\mathcal{D}} \left(\pd{\rho}{t} + \boldsymbol{\nabla}\cdot(\rho\bm{U})\right)d\mathcal{V} = 0.
\]

Comme le domaine $\mathcal{D}$ est arbitraire, l'intégrand doit être nul en tout
point :

\begin{importantbox}{Équation de continuité}
\[
    \pd{\rho}{t} + \boldsymbol{\nabla}\cdot(\rho\bm{U}) = 0.
\]
Cette relation, valable pour tout fluide (compressible ou non), est la transcription
locale de la conservation de la masse.
\end{importantbox}

\subsection{L'incompressibilité : condition de divergence nulle}

Pour un fluide \textbf{incompressible}, chaque particule conserve sa densité :
$D\rho/Dt = 0$. On peut alors développer l'équation de continuité :
\[
    \underbrace{\pd{\rho}{t} + \bm{U}\cdot\boldsymbol{\nabla}\rho}_{= D\rho/Dt = 0}
    + \rho\,\boldsymbol{\nabla}\cdot\bm{U} = 0
    \quad\Longrightarrow\quad
    \rho\,\boldsymbol{\nabla}\cdot\bm{U} = 0.
\]

Comme $\rho \neq 0$, on conclut :

\begin{importantbox}{Incompressibilite : divergence nulle}
\[
    \boldsymbol{\nabla}\cdot\bm{U} = 0.
\]
Cette condition est vérifiée pour tout écoulement vérifiant $M_a \lesssim 0{,}3$.
Elle signifie que le champ de vitesse est \textbf{à divergence nulle} : tout volume
de fluide se conserve (ni compression, ni dilatation).
\end{importantbox}

\subsection{Application : conservation du débit volumique en conduite}

L'une des conséquences directes de $\boldsymbol{\nabla}\cdot\bm{U} = 0$ est la
conservation du débit dans une conduite. Considérons un conduit avec deux sections
$S_1$ et $S_2$ et une paroi imperméable $S_w$ ($\bm{U}\cdot\bm{n}=0$). En intégrant
la condition d'incompressibilité sur le volume du conduit et en appliquant le théorème
de la divergence, le flux net à travers la surface fermée $\mathcal{S} = S_1 \cup S_2
\cup S_w$ est nul :
\[
    \int_{S_1} \bm{U}\cdot\bm{n}\,ds + \int_{S_2} \bm{U}\cdot\bm{n}\,ds
    + \underbrace{\int_{S_w} \bm{U}\cdot\bm{n}\,ds}_{=\,0} = 0.
\]

\begin{importantbox}{Conservation du débit volumique}
Le débit volumique $Q_v = \int_S \bm{U}\cdot\bm{n}\,ds$ est conservé le long d'une
conduite en écoulement incompressible. La \textbf{vitesse de débit}
$U_d = Q_v/S$ est donc inversement proportionnelle à la section.
\end{importantbox}

Léonard de Vinci avait empiriquement constaté ce résultat sur les cours d'eau dès
1502---avant même que les mathématiques de la mécanique des fluides n'existent.

% ------------------------------------------
\section{Les forces sur un domaine fluide}
% ------------------------------------------

Pour écrire le principe fondamental de la dynamique, il faut recenser toutes les
forces qui agissent sur une portion de fluide. Ces forces sont de deux types.

\subsection{Forces volumiques et forces surfaciques}

La force totale $\bm{F}$ sur un domaine $\mathcal{D}$ se décompose en :
\[
    \bm{F} = \bm{F}_s + \bm{F}_v.
\]

Les \textbf{forces volumiques} $\bm{F}_v$ agissent à distance sur l'ensemble du
volume (pesanteur, forces électromagnétiques pour les fluides conducteurs). Pour
la gravité :
\[
    \bm{F}_v = \int_{\mathcal{D}} \rho\bm{g}\,d\mathcal{V}.
\]

Les \textbf{forces surfaciques} $\bm{F}_s$ résultent des interactions (essentiellement
des chocs moléculaires) entre le domaine $\mathcal{D}$ et le fluide environnant, au
travers de la surface de contrôle $\mathcal{S}$. On introduit le \textbf{vecteur
contrainte} $\bm{T}(\bm{x},t,\bm{n})$ (en Pa), force par unité de surface, tel
que :
\[
    \bm{F}_s = \int_{\mathcal{S}} \bm{T}\,ds.
\]

Le vecteur $\bm{n}$ dans l'argument de $\bm{T}$ est la normale extérieure à $\mathcal{S}$ :
la contrainte dépend de l'orientation de la surface sur laquelle elle s'exerce. La
troisième loi de Newton impose immédiatement :
\[
    \bm{T}(\bm{x},t,-\bm{n}) = -\bm{T}(\bm{x},t,\bm{n}).
\]

\subsection{Le tenseur des contraintes de Cauchy}

La dépendance de $\bm{T}$ en $\bm{n}$ pourrait a priori être quelconque. Augustin-Louis
Cauchy (1789--1857) a démontré qu'elle est en réalité \textbf{linéaire} : il existe
un tenseur d'ordre~2, le \textbf{tenseur des contraintes} $\ten{\sigma}$, tel que
$\bm{T}$ s'exprime simplement.

L'idée de la preuve est de calculer $\bm{T}(\bm{e}_i)$ sur chaque face du trièdre
de base. Par linéarité de la décomposition, $\bm{T}(\bm{n})$ pour une normale
quelconque s'obtient en combinant ces trois vecteurs contrainte via le tenseur :

\begin{defbox}{Tenseur des contraintes de Cauchy}
Le tenseur des contraintes $\ten{\sigma}$ est défini par ses composantes :
\[
    \sigma_{ij} \equiv T_i(\bm{e}_j)
    \quad \text{(composante }i\text{ de la contrainte sur une face de normale }\bm{e}_j\text{)}.
\]
Le vecteur contrainte sur une surface de normale $\bm{n}$ quelconque vaut :
\[
    \bm{T}(\bm{x},t,\bm{n}) = \ten{\sigma}\cdot\bm{n}.
\]
Le tenseur $\ten{\sigma}$ est \textbf{symétrique} : $\sigma_{ij} = \sigma_{ji}$
(résultat de l'équilibre des moments cinétiques internes).
\end{defbox}

On peut alors exprimer la force surfacique totale. On applique le théorème de la
divergence à $\bm{F}_s = \int_\mathcal{S}\ten{\sigma}\cdot\bm{n}\,ds$, ce qui
convertit l'intégrale de surface en intégrale de volume :
\[
    \bm{F}_s = \int_{\mathcal{S}} \ten{\sigma}\cdot\bm{n}\,ds
             = \int_{\mathcal{D}} \boldsymbol{\nabla}\cdot\ten{\sigma}\,d\mathcal{V}.
\]

\textbf{Cas du fluide non visqueux ou au repos.} Quand seule la pression contribue
aux contraintes surfaciques, la contrainte est purement normale et isotrope :
$\bm{T} = -p\bm{n}$ (le signe moins vient du fait que la pression comprime le
volume). Le tenseur se réduit alors à :
\[
    \ten{\sigma} = -p\,\ten{I}
    = \begin{pmatrix} -p & 0 & 0 \\ 0 & -p & 0 \\ 0 & 0 & -p \end{pmatrix},
    \qquad
    \bm{F}_s = -\int_\mathcal{S} p\bm{n}\,ds.
\]

% ------------------------------------------
\section{Conservation de la quantité de mouvement}
% ------------------------------------------

On a maintenant tous les ingrédients. On applique le principe fondamental de la
dynamique ($\bm{F} = m\bm{a}$) à une particule fluide de masse $\rho\,d\mathcal{V}$ :
\[
    \rho\,d\mathcal{V}\cdot\frac{D\bm{U}}{Dt}
    = d\bm{F}_s + d\bm{F}_v.
\]

En exprimant $\bm{F}_s$ et $\bm{F}_v$ sous forme volumique et en identifiant les
intégrands (valable pour tout domaine arbitraire), on obtient la forme locale :

\begin{importantbox}{Conservation de la quantité de mouvement}
\[
    \rho\,\frac{D\bm{U}}{Dt} = \boldsymbol{\nabla}\cdot\ten{\sigma} + \rho\bm{g}.
\]
C'est le ``$m\bm{a} = \bm{F}$'' de la mécanique des fluides : l'accélération
matérielle de la particule est égale à la résultante des forces surfaciques
($\boldsymbol{\nabla}\cdot\ten{\sigma}$) et volumiques ($\rho\bm{g}$).
\end{importantbox}

Via le théorème de Reynolds, cette même conservation s'écrit sous forme intégrale
pour un domaine quelconque (fixe ou mobile) :
\[
    \frac{d}{dt}\int_{\mathcal{D}} \rho\bm{U}\,d\mathcal{V}
    = \int_{\mathcal{S}} \ten{\sigma}\cdot\bm{n}\,ds
    + \int_{\mathcal{D}} \rho\bm{g}\,d\mathcal{V}
    + \int_{\mathcal{S}} \rho\bm{U}(\bm{U}_\mathcal{S}-\bm{U})\cdot\bm{n}\,ds.
\]

% ------------------------------------------
\section{Hydrostatique : quand le fluide ne bouge pas}
% ------------------------------------------

Avant d'aborder l'écoulement non visqueux, traitons le cas le plus simple qui soit :
le fluide au repos, $\bm{U} = \bm{0}$. C'est l'\textbf{hydrostatique}.

\subsection{Équilibre hydrostatique et profil de pression}

Quand $\bm{U} = \bm{0}$, l'accélération est nulle et $\ten{\sigma} = -p\ten{I}$
(pas de viscosité, le fluide ne cisaille rien). L'équation de conservation de
la quantité de mouvement se réduit à :
\[
    \bm{0} = -\boldsymbol{\nabla}p + \rho\bm{g}.
\]

On choisit $x_3$ comme axe vertical ascendant, $\bm{g} = -g\bm{e}_3$. Les
composantes horizontales donnent $\partial p/\partial x_1 = \partial p/\partial x_2 = 0$,
et la composante verticale donne $\partial p/\partial x_3 = -\rho g$. Par intégration :

\begin{importantbox}{Profil de pression hydrostatique}
\[
    p(x_3) = p_0 - \rho g x_3,
\]
où $p_0$ est la pression à $x_3 = 0$ (surface libre). La pression augmente avec
la profondeur. Pour l'eau ($\rho \simeq 10^3~\text{kg/m}^3$), elle augmente d'environ
une atmosphère ($10^5~\text{Pa}$) tous les 10~m.
\end{importantbox}

\subsection{La force d'Archimède}

Considérons un corps immergé dans le fluide au repos. La force exercée par le fluide
sur ce corps est une force surfacique associée à la pression. Elle vaut :
\[
    \bm{F}_A = \int_\mathcal{S} (-p\ten{I})\cdot\bm{n}\,ds
             = -\int_\mathcal{S} p\bm{n}\,ds.
\]

On applique le théorème de la divergence (en se rappelant que
$\boldsymbol{\nabla}p = \rho\bm{g}$ à l'équilibre) :
\[
    \bm{F}_A = -\int_\mathcal{D} \boldsymbol{\nabla}p\,d\mathcal{V}
             = -\int_\mathcal{D} \rho\bm{g}\,d\mathcal{V}
             = -\bm{g}\int_\mathcal{D} \rho\,d\mathcal{V}
             = -\bm{g}\,M_f.
\]

\begin{importantbox}{Principe d'Archimède}
La force de flottabilité exercée par un fluide au repos sur un corps immergé est :
\[
    \bm{F}_A = -M_f\bm{g},
\]
où $M_f$ est la masse du fluide déplacé. Cette force est verticale, ascendante,
et s'applique au \textbf{centre de masse du fluide déplacé}
$\bm{x}_A$, défini par $M_f\bm{x}_A \equiv \int_\mathcal{D}\rho\bm{x}\,d\mathcal{V}$.
\end{importantbox}

Si le pilote du Rafale s'éjecte en mer, il flotte parce que sa densité moyenne
(humain + combinaison) est inférieure à celle de l'eau de mer
($\rho_{mer} \approx 1025~\text{kg/m}^3$). Archimède (287--212 av.~J.-C.) aurait
compris le problème de survie en mission.

\textbf{Stabilité d'un corps immergé.} Pour qu'un corps complètement immergé (un
sous-marin, par exemple) soit stable, son centre de gravité $\bm{x}_G$ doit se
trouver \textit{en dessous} du centre $\bm{x}_A$ du fluide déplacé. Si c'est
l'inverse, un léger roulis crée un couple déstabilisant qui chavire le corps.

% ------------------------------------------
\section{Le modèle non visqueux}
% ------------------------------------------

\subsection{Motivations}

Tout fluide réel possède une viscosité : c'est sa résistance interne au cisaillement,
la propriété qui fait qu'un filet d'eau ne glisse pas librement sur une surface. Dans
l'expression générale du tenseur des contraintes, la viscosité introduit des termes
supplémentaires au-delà de la pression (nous les étudierons au chapitre~2).

Cependant, pour des fluides peu visqueux comme l'air, et loin des parois où les
frottements se concentrent, ces termes visqueux sont souvent négligeables devant les
effets de pression. On peut alors adopter le \textbf{modèle non visqueux} (ou fluide
parfait) :

\begin{defbox}{Modèle non visqueux}
Dans le modèle non visqueux, le tenseur des contraintes se réduit à la pression :
\[
    \ten{\sigma} = -p\,\ten{I}, \qquad \bm{T} = -p\bm{n}.
\]
Il n'y a ni frottement ni cisaillement : la seule force surfacique est due à la
pression, qui agit toujours perpendiculairement à la surface.
\end{defbox}

\subsection{Les équations d'Euler}

En substituant $\ten{\sigma} = -p\ten{I}$ dans l'équation de conservation de la
quantité de mouvement, et en utilisant
$\boldsymbol{\nabla}\cdot(-p\ten{I}) = -\boldsymbol{\nabla}p$, on obtient
les équations d'Euler (1757) :

\begin{importantbox}{Équations d'Euler --- écoulement incompressible non visqueux}
\[
    \begin{cases}
        \dfrac{D\bm{U}}{Dt} = -\dfrac{1}{\rho}\,\boldsymbol{\nabla}p + \bm{g}
        \\[10pt]
        \boldsymbol{\nabla}\cdot\bm{U} = 0.
    \end{cases}
\]
\end{importantbox}

\textbf{Conditions aux limites.} Sans viscosité, on ne peut imposer qu'une seule
condition sur $\bm{U}$ à la paroi : l'\textbf{imperméabilité}, c'est-à-dire que le
fluide ne traverse pas la paroi,
\[
    \bm{U}\cdot\bm{n} = 0 \quad \text{(paroi imperméable)}.
\]
Le fluide peut donc glisser tangentiellement le long de la paroi, contrairement à la
réalité visqueuse. Pour une interface entre deux milieux $A$ et $B$, la condition
devient $\bm{U}_A\cdot\bm{n} = \bm{U}_B\cdot\bm{n}$.

\subsection{L'équation de Bernoulli}

On dispose des équations d'Euler. Pour en extraire une information sur l'énergie,
on prend le produit scalaire avec la vitesse $\bm{U}$ :
\[
    \bm{U}\cdot\frac{D\bm{U}}{Dt}
    = \bm{U}\cdot\!\left(-\frac{1}{\rho}\,\boldsymbol{\nabla}p + \bm{g}\right).
\]

Le membre gauche est la puissance de l'accélération. On reconnaît l'identité
$\bm{U}\cdot\dfrac{D\bm{U}}{Dt} = \dfrac{D}{Dt}\!\left(\dfrac{U^2}{2}\right)$.

Pour un écoulement \textbf{stationnaire} ($\partial_t = 0$), la dérivée matérielle
agit comme $D/Dt = \bm{U}\cdot\boldsymbol{\nabla}$. Le membre droit se réorganise
grâce à deux observations :
\begin{itemize}
    \item Pour un fluide incompressible ($\rho = \text{cst}$) :
    $\dfrac{1}{\rho}\,\boldsymbol{\nabla}p = \boldsymbol{\nabla}\!\left(\dfrac{p}{\rho}\right)$.
    \item La pesanteur est une force conservative :
    $\bm{g} = \boldsymbol{\nabla}(\bm{g}\cdot\bm{x})$, soit avec
    $\bm{g} = -g\bm{e}_3$ : $\bm{g} = \boldsymbol{\nabla}(-gx_3)$.
\end{itemize}

L'équation d'énergie cinétique se récrit :
\[
    \bm{U}\cdot\boldsymbol{\nabla}\!\left(\frac{U^2}{2} + \frac{p}{\rho} - gx_3\right) = 0.
\]

Cette relation signifie que la quantité entre parenthèses est constante dans la
direction de $\bm{U}$, c'est-à-dire \textit{le long d'une ligne de courant}.
En multipliant par $\rho$ :

\begin{importantbox}{Équation de Bernoulli}
Pour un écoulement \textbf{non visqueux}, \textbf{stationnaire},
\textbf{incompressible}, sous \textbf{forces volumiques conservatives}
(pesanteur) :
\[
    \mathcal{H} \;\equiv\; p + \rho\,g\,x_3 + \rho\,\frac{U^2}{2} = \text{cst}
    \quad \text{le long d'une ligne de courant.}
\]
$\mathcal{H}$ (en J/m$^3$) est la \emph{charge totale} : somme de la pression
statique, de la pression hydrostatique et de la \emph{pression dynamique}
$q = \tfrac{1}{2}\rho U^2$.
\end{importantbox}

La physique est limpide : si le fluide accélère (car une conduite se rétrécit, ou
car l'extrados de l'aile oblige les lignes de courant à se resserrer), la pression
diminue pour compenser. C'est exactement ce qui génère la portance sur l'aile du
Rafale : l'air s'écoule plus vite sur l'extrados que sur l'intrados, donc la
pression est plus faible au-dessus qu'en dessous, et cette différence de pression
crée une force ascendante. Voilà pourquoi le Rafale ne coule pas dans l'océan.

$\mathcal{H}$ prend des valeurs potentiellement différentes d'une ligne de courant
à l'autre. Pour un écoulement stationnaire, les lignes de courant coïncident avec
les trajectoires : $\mathcal{H}$ est donc conservé pour une particule fluide donnée
tout au long de son trajet.

\subsection{Application : le tube de Pitot}

Le Rafale embarque plusieurs \textbf{tubes de Pitot} sur le nez et le fuselage.
Ces sondes mesurent la pression en deux points d'une même ligne de courant et en
déduisent la vitesse par application directe de Bernoulli.

\vspace{0.5cm}
\begin{center}
\begin{tikzpicture}[scale=1.0]
    % Corps du tube
    \fill[gray!30] (-4, 0.55) rectangle (2, -0.55);
    % Pointe (point d'arrêt)
    \fill[gray!50] (2, 0.55) -- (3, 0) -- (2, -0.55) -- cycle;
    \draw[thick] (-4, 0.55) -- (2, 0.55) -- (3, 0) -- (2,-0.55) -- (-4,-0.55);
    % Point d'arrêt
    \fill[red!80!black] (3, 0) circle (3pt);
    \node[right=4pt] at (3.05, 0) {$p_0$ (point d'arrêt, $U=0$)};
    % Prise de pression statique
    \draw[thick] (-0.5, 0.55) -- (-0.5, 1.4);
    \fill[white] (-0.5, 1.2) circle (6pt);
    \node[above] at (-0.5, 1.5) {$p_1$ (pression statique)};
    % Flèches d'écoulement amont
    \foreach \y in {1.1, 0.55, 0, -0.55, -1.1} {
        \draw[->, >=Stealth, thick, blue!70!black] (-6.5, \y) -- (-5.5, \y);
    }
    \node[left=4pt] at (-6.5, 0) {$U_1,\; p_1,\; \rho_1$};
    % Capteur différentiel (manomètre)
    \draw[thick] (-0.5, -0.55) -- (-0.5, -1.5) -- (3.0, -1.5) -- (3.0, 0);
    \node[below] at (1.25, -1.5) {capteur $\Delta p = p_0 - p_1$};
\end{tikzpicture}
\end{center}
\vspace{0.3cm}

Le long de la ligne de courant centrale, qui aboutit au point d'arrêt (où
$U = 0$), Bernoulli donne (à $x_3$ constant, donc sans terme hydrostatique) :
\[
    p_1 + \rho_1\,\frac{U_1^2}{2} = p_0 + \rho_1\cdot\frac{0^2}{2} = p_0.
\]

\begin{importantbox}{Mesure de vitesse par tube de Pitot}
\[
    U_1 = \sqrt{\dfrac{2(p_0 - p_1)}{\rho_1}}.
\]
$p_0$ est la \textbf{pression totale} (ou d'arrêt), $p_1$ la \textbf{pression
statique}, et $q_1 = \tfrac{1}{2}\rho_1 U_1^2$ la \textbf{pression dynamique}.
La relation $p_0 = p_1 + q_1$ n'est valable \emph{qu'en régime incompressible}
($M_a \lesssim 0{,}3$) ; une formulation corrigée pour le régime compressible
sera vue au chapitre~7.
\end{importantbox}

Les calculateurs avioniques du Rafale évaluent $U_1$ en temps réel à partir de
la mesure de $\Delta p = p_0 - p_1$. Henri Pitot avait développé cet instrument
en 1732 pour mesurer la vitesse de la Seine à Paris. Aujourd'hui, le même principe
physique équipe les Airbus A350, les F-22, et les sondes spatiales.

% ------------------------------------------
\section*{Résumé du chapitre}
\addcontentsline{toc}{section}{Résumé du chapitre}
% ------------------------------------------

\begin{importantbox}{Ce qu'il faut retenir}
\begin{itemize}
    \item \textbf{Milieu continu} : valide dès que $Kn = \lambda/L \ll 1$.
    Incompressible si $M_a \lesssim 0{,}3$.

    \item \textbf{Dérivée matérielle} :
    $\dfrac{Df}{Dt} = \dfrac{\partial f}{\partial t} + \bm{U}\cdot\boldsymbol{\nabla}f$
    --- variation vue en suivant une particule fluide.

    \item \textbf{Théorème de Reynolds} :
    $\dfrac{d}{dt}\displaystyle\int_\mathcal{D}\varphi\,d\mathcal{V}
    = \displaystyle\int_\mathcal{D}\dfrac{\partial\varphi}{\partial t}\,d\mathcal{V}
    + \displaystyle\int_\mathcal{S}\varphi\,\bm{U}_\mathcal{S}\cdot\bm{n}\,ds$.

    \item \textbf{Conservation de la masse} :
    $\dfrac{\partial\rho}{\partial t} + \boldsymbol{\nabla}\cdot(\rho\bm{U}) = 0$.
    En incompressible : $\boldsymbol{\nabla}\cdot\bm{U} = 0$.

    \item \textbf{Tenseur des contraintes} : $\bm{T} = \ten{\sigma}\cdot\bm{n}$,
    $\ten{\sigma}$ symétrique.

    \item \textbf{Conservation de la QDM} :
    $\rho\dfrac{D\bm{U}}{Dt} = \boldsymbol{\nabla}\cdot\ten{\sigma} + \rho\bm{g}$.

    \item \textbf{Fluide non visqueux} :
    $\ten{\sigma} = -p\,\ten{I}$, $\bm{T} = -p\bm{n}$.
    Hydrostatique : $p = p_0 - \rho g x_3$.
    Archimède : $\bm{F}_A = -M_f\bm{g}$.

    \item \textbf{Équations d'Euler} :
    $\dfrac{D\bm{U}}{Dt} = -\dfrac{1}{\rho}\boldsymbol{\nabla}p + \bm{g}$
    et $\boldsymbol{\nabla}\cdot\bm{U} = 0$.

    \item \textbf{Bernoulli} (non visqueux, stationnaire, incompressible,
    forces conservatives) :
    $p + \rho g x_3 + \dfrac{\rho U^2}{2} = \text{cst}$ le long d'une ligne
    de courant. Application directe : tube de Pitot.
\end{itemize}
\end{importantbox}

% ==========================================
\chapter{Fluide visqueux newtonien}
% ==========================================

\noindent
\textit{Fil conducteur --- le Rafale Marine.}
Au chapitre précédent, nous avons posé les équations de conservation de la masse
et de la quantité de mouvement, en traitant le fluide comme parfait (non visqueux).
Ce modèle était suffisant pour comprendre la portance en régime de vol rapide via
Bernoulli, et pour calculer la poussée d'Archimède lors d'un éjection en mer.
Mais il reste muet sur une réalité fondamentale~: la résistance de l'air.
Un Rafale en vol de convoyage consomme du kérosène non pas pour vaincre l'inertie
(il vole à vitesse constante) mais pour compenser la \emph{traînée}, cette force
qui s'oppose à son déplacement et qui naît des frottements visqueux entre l'air et
la voilure. Pour décrire ces frottements, il faut enrichir notre modèle du tenseur
des contraintes~: c'est l'objet de ce chapitre.

% ==========================================
\section{Du tenseur de Cauchy au fluide newtonien}
% ==========================================

\subsection{Décomposition du tenseur des contraintes}

Au chapitre 1 nous avons introduit le tenseur de Cauchy $\ten{\sigma}$ qui encode
l'intégralité des efforts de contact exercés par le fluide environnant sur une
particule fluide. Un résultat général de la mécanique des milieux continus est que
ce tenseur peut toujours se décomposer en deux contributions de nature distincte~:

\begin{equation}
    \ten{\sigma} = -p\,\ten{I} + \ten{\tau}
    \label{eq:decomp_sigma}
\end{equation}

Cette décomposition \emph{n'est pas une définition arbitraire}~: elle reflète une
réalité physique. Considérons un fluide au repos. Il exerce une pression $p$ sur
toutes les faces d'un volume de contrôle, de manière isotrope (égale dans toutes
les directions). Cette contribution isotrope s'écrit précisément $-p\,\ten{I}$~: le
signe moins vient de la convention que $\bm{n}$ pointe vers l'extérieur, et la
pression pousse vers l'intérieur. Dès que le fluide se met en mouvement, des
contraintes \emph{déviatoriques} $\ten{\tau}$ apparaissent en plus~: ce sont les
contraintes visqueuses, liées aux gradients de vitesse et donc nulles au repos.

\begin{defbox}{Décomposition du tenseur des contraintes}
\[
    \ten{\sigma} = -p\,\ten{I} + \ten{\tau}
\]
$p$ est la pression thermodynamique (Pa), $\ten{I}$ l'identité, et $\ten{\tau}$
le \textbf{tenseur des contraintes visqueuses} (Pa), nul pour un fluide parfait.
\end{defbox}

\subsection{Tenseur des taux de déformation}

La question qui se pose naturellement est~: de quoi dépend $\ten{\tau}$~?
Physiquement, les frottements visqueux naissent de la \emph{déformation} du fluide~:
c'est parce que les couches fluides glissent les unes sur les autres que des forces
tangentielles apparaissent. La mesure locale de cette déformation est encodée dans
le gradient de vitesse $\nabla\bm{U}$, dont la composante $(i,j)$ est
$\partial U_i/\partial x_j$.

Ce tenseur gradient peut être décomposé de manière unique en une partie symétrique
et une partie antisymétrique~:

\begin{equation}
    \nabla\bm{U} = \ten{D} + \ten{\Omega}
\end{equation}

avec~:

\begin{align}
    D_{ij} &= \frac{1}{2}\left(\pd{U_i}{x_j} + \pd{U_j}{x_i}\right)
    & \ten{D} &~: \text{tenseur des taux de déformation (symétrique)} \\[6pt]
    \Omega_{ij} &= \frac{1}{2}\left(\pd{U_i}{x_j} - \pd{U_j}{x_i}\right)
    & \ten{\Omega} &~: \text{tenseur des taux de rotation (antisymétrique)}
\end{align}

Pourquoi ce partage~? Parce que $\ten{\Omega}$ décrit une rotation de corps rigide
locale~: elle ne déforme pas le fluide, donc elle ne génère aucune contrainte
visqueuse. Seule $\ten{D}$, qui mesure la déformation irréversible, peut produire
des frottements. On peut le montrer formellement~: le vecteur tourbillon
$\bm{\omega} = \frac{1}{2}\nabla\times\bm{U}$ est lié à $\ten{\Omega}$ par
$\Omega_{ij} = -\varepsilon_{ijk}\omega_k$.

\begin{defbox}{Tenseur des taux de déformation}
\[
    D_{ij} = \frac{1}{2}\left(\pd{U_i}{x_j} + \pd{U_j}{x_i}\right)
    \qquad \Longleftrightarrow \qquad
    \ten{D} = \frac{1}{2}\left(\nabla\bm{U} + (\nabla\bm{U})^T\right)
\]
$\ten{D}$ est symétrique~; ses composantes diagonales $D_{ii} = \partial U_i/\partial x_i$
sont des taux d'élongation, et ses composantes hors-diagonale $D_{ij}$ ($i\neq j$) sont
des taux de cisaillement.
\end{defbox}

\textit{Illustration.} L'écoulement de Couette plan---air entre deux plaques, la
plaque supérieure se déplaçant à $U_e$---donne $\bm{U} = (U_e x_2/h,\,0,\,0)$.
Seul $D_{12}=D_{21}=U_e/(2h)$ est non nul~: l'écoulement est un pur cisaillement.
C'est exactement la situation qui règne dans la couche limite contre la peau du Rafale.

\subsection{Loi constitutive newtonienne}

Il reste à relier $\ten{\tau}$ à $\ten{D}$. La relation la plus simple---et la plus
répandue---est la proportionnalité linéaire~: c'est la \textbf{loi de Newton} pour les
fluides visqueux. Elle énonce que les contraintes visqueuses sont proportionnelles aux
taux de déformation~:

\begin{importantbox}{Loi de Newton pour un fluide visqueux newtonien (incompressible)}
\[
    \ten{\tau} = 2\mu\,\ten{D}
    \qquad\Longleftrightarrow\qquad
    \tau_{ij} = \mu\left(\pd{U_i}{x_j} + \pd{U_j}{x_i}\right)
\]
$\mu$ (Pa$\cdot$s) est la \textbf{viscosité dynamique}, propriété du fluide, mesurée
expérimentalement~: $\mu_\text{air}\approx 1{,}8\times10^{-5}$~Pa$\cdot$s à $20\,^{\circ}$C,
$\mu_\text{eau}\approx 10^{-3}$~Pa$\cdot$s.
\end{importantbox}

Ce modèle est dit \emph{newtonien} car il postule la linéarité entre contrainte et taux
de déformation. Il fonctionne très bien pour l'air et l'eau, mais pas pour tous les
fluides (cf.\ section suivante).

\subsection{Lois rhéologiques --- au-delà du modèle newtonien}

La loi de Newton n'est qu'un cas particulier. Pour un écoulement de Couette
($\tau = \tau_{12}$, $\dot{\gamma}=\partial U_1/\partial x_2$), on peut classer les
comportements rhéologiques~:

\begin{center}
\renewcommand{\arraystretch}{1.4}
\begin{tabular}{lll}
\toprule
\textbf{Comportement} & \textbf{Loi} & \textbf{Exemples} \\
\midrule
Fluide parfait & $\tau=0$ & air à grande échelle \\
Newtonien & $\tau=\mu\dot{\gamma}$ & air, eau, huile \\
Rhéo-épaississant & $\tau=k\dot{\gamma}^n$, $n>1$ & amidon dans l'eau \\
Rhéo-fluidifiant & $\tau=k\dot{\gamma}^n$, $n<1$ & sang, ketchup, moutarde \\
Bingham & $\tau=\tau_0+\mu_B\dot{\gamma}$ si $\tau>\tau_0$ & chocolat, dentifrice \\
\bottomrule
\end{tabular}
\end{center}

Le sang est rhéo-fluidifiant~: sa viscosité effective diminue quand le cœur pompe
plus fort (gradient de vitesse élevé dans les artères). Le ketchup obéit à un
modèle de Bingham~: il faut exercer un seuil de contrainte $\tau_0$ pour le faire
couler. Dans ce chapitre, nous travaillons exclusivement avec des fluides newtoniens
(air, eau)~; le Rafale ne carbure pas au ketchup.

% ==========================================
\section{Écoulement de Couette plan}
% ==========================================

Avant de poser les équations générales, ancrons nos idées sur un cas canonique~:
l'écoulement de Couette plan entre deux plaques parallèles distantes de $h$, la
plaque inférieure fixe, la plaque supérieure entraînée à vitesse $U_e$.

Ce profil de vitesse est~: $\bm{U}=(U_e x_2/h,\, 0,\, 0)$.
On peut vérifier qu'il satisfait $\nabla\cdot\bm{U}=0$ (incompressible) et que les
termes d'accélération sont nuls (écoulement établi, stationnaire).

\begin{center}
\begin{tikzpicture}[scale=1.2]
    % plaques
    \fill[gray!30] (-0.2,0) rectangle (4.2,-0.15);
    \fill[gray!30] (-0.2,2) rectangle (4.2,2.15);
    \draw[thick] (-0.2,0) -- (4.2,0);
    \draw[thick] (-0.2,2) -- (4.2,2);
    % hachures plaque inférieure
    \foreach \x in {0,0.4,...,4} {
        \draw[gray] (\x,0) -- (\x-0.2,-0.2);
    }
    % hachures plaque supérieure
    \foreach \x in {0,0.4,...,4} {
        \draw[gray] (\x,2) -- (\x-0.2,2.2);
    }
    % profil linéaire
    \foreach \y in {0,0.4,0.8,1.2,1.6,2.0} {
        \pgfmathsetmacro{\len}{1.6*\y/2}
        \draw[-{Latex}, blue!70!black, thick] (1,\y) -- (1+\len,\y);
    }
    % accolade h
    \draw[<->,gray] (-0.05,0) -- (-0.05,2);
    \node[left,gray] at (-0.1,1) {$h$};
    % étiquettes
    \node[right] at (2.6,2) {$U_e$};
    \node[right] at (1,0.2) {\small $U_1(x_2)=U_e\dfrac{x_2}{h}$};
    % axe
    \draw[->,gray] (3.8,-0.3) -- (3.8,2.4) node[right] {$x_2$};
\end{tikzpicture}
\end{center}

Le tenseur des contraintes se calcule directement~:
\[
    \ten{D} = \begin{pmatrix} 0 & \frac{U_e}{2h} & 0 \\ \frac{U_e}{2h} & 0 & 0 \\ 0 & 0 & 0 \end{pmatrix},
    \qquad
    \ten{\tau} = 2\mu\ten{D} = \begin{pmatrix} 0 & \frac{\mu U_e}{h} & 0 \\ \frac{\mu U_e}{h} & 0 & 0 \\ 0 & 0 & 0 \end{pmatrix}
\]

La contrainte de cisaillement est uniforme dans tout l'écoulement~:

\begin{importantbox}{Contrainte pariétale en Couette plan}
\[
    \tau_w = \tau_{12} = \mu\,\frac{U_e}{h}
\]
Plus le gradient de vitesse est grand (couche mince, vitesse élevée), plus le
frottement est intense. C'est précisément ce qui se passe dans la couche limite
contre la peau du Rafale~: $U_e/h \to \infty$ car $h$ est très petit.
\end{importantbox}

% ==========================================
\section{Forces aérodynamiques~: portance et traînée}
% ==========================================

Un avion en vol est soumis à deux forces aérodynamiques résultantes~: la
\textbf{portance} $F_L$ perpendiculaire à la direction de vol $\bm{U}_\infty$, et la
\textbf{traînée} $F_D$ parallèle à $\bm{U}_\infty$ (s'opposant au mouvement).
Ces forces résultent de l'intégration du vecteur contrainte $\bm{T}=\ten{\sigma}\cdot\bm{n}$
sur la surface de l'aile~:

\[
    \bm{F} = \oint_{\text{aile}} \ten{\sigma}\cdot\bm{n}\,dS
    = \underbrace{-\oint p\,\bm{n}\,dS}_{\text{contrib.\ pression}}
      + \underbrace{\oint \ten{\tau}\cdot\bm{n}\,dS}_{\text{contrib.\ visqueuse}}
\]

La portance résulte principalement de la différence de pression entre l'intrados
et l'extrados (Bernoulli~: l'air accélère sur l'extrados, la pression diminue).
La traînée a deux origines~: la traînée de forme (liée aux recirculations en aval)
et la traînée de frottement (liée à $\ten{\tau}\cdot\bm{n}$ sur la surface).

\begin{defbox}{Finesse aérodynamique}
\[
    f = \frac{F_L}{F_D}
\]
La finesse mesure l'efficacité aérodynamique d'un aéronef~: pour un angle de
descente $\alpha$, la distance planée vaut $L = h/\tan\alpha \approx f \times h$.
\end{defbox}

\medskip
\noindent Quelques ordres de grandeur instructifs~:

\begin{center}
\renewcommand{\arraystretch}{1.3}
\begin{tabular}{lc}
\toprule
\textbf{Aéronef} & \textbf{Finesse} $F_L/F_D$ \\
\midrule
Planeur ETA & 70 \\
Solar Impulse 2 & $\approx 34$ \\
Airbus A380 & 20 \\
Concorde & 7{,}5 \\
\textbf{Rafale} & \textbf{$\approx$ 5--9} \\
Navette spatiale & $\approx 1$ \\
\bottomrule
\end{tabular}
\end{center}

\medskip
La faible finesse du Rafale (comparée à un liner) n'est pas un défaut~: c'est le
prix de la manœuvrabilité. Un avion de chasse doit pouvoir virer à $9g$, accélérer
de $0$ à Mach~2 en quelques secondes, et voler à très basse altitude à grande
vitesse. Ces performances exigent une voilure courte, épaisse, à faible allongement,
ce qui génère beaucoup de traînée induite. Le moteur M88-2 (73~kN de poussée unitaire)
compense cette traînée en permanence.

% ==========================================
\section{Équations de Navier-Stokes}
% ==========================================

Nous avons maintenant tous les ingrédients pour écrire les équations gouvernant
un fluide visqueux newtonien incompressible. L'équation de conservation de la
quantité de mouvement du chapitre 1 est~:
\[
    \rho\,\DD{\bm{U}} = \nabla\cdot\ten{\sigma} + \rho\bm{g}
\]
On substitue la décomposition $\ten{\sigma}=-p\ten{I}+\ten{\tau}$ et la loi de
Newton $\ten{\tau}=2\mu\ten{D}$. Pour un fluide incompressible
($\nabla\cdot\bm{U}=0$), on montre que $\nabla\cdot(2\mu\ten{D})=\mu\nabla^2\bm{U}$
(le terme croisé disparaît grâce à la condition d'incompressibilité). On obtient~:

\begin{importantbox}{Équations de Navier-Stokes --- fluide newtonien incompressible}
\[
    \nabla\cdot\bm{U} = 0
\]
\[
    \DD{\bm{U}} = -\frac{1}{\rho}\nabla p + \nu\nabla^2\bm{U} + \bm{g}
\]
avec la \textbf{viscosité cinématique} $\nu = \mu/\rho$ (m$^2$/s)~:
\quad $\nu_\text{air} \approx 1{,}5\times10^{-5}$~m$^2$/s,
\quad $\nu_\text{eau} \approx 1{,}0\times10^{-6}$~m$^2$/s.
\end{importantbox}

Ces équations ont été établies indépendamment par \textbf{Navier} (1822) et
\textbf{Stokes} (1845). Elles forment le modèle de référence pour tous les
écoulements de fluides newtoniens~: elles décrivent aussi bien le décollement de
la couche limite sur l'aile du Rafale que la formation des ouragans dans
l'Atlantique. Leur solution générale reste un problème ouvert~: l'existence et la
régularité des solutions en trois dimensions est l'un des sept problèmes du
millénaire (Clay Mathematics Institute, 2000, doté de 1~million de dollars).

Le terme $\nu\nabla^2\bm{U}$ est le terme visqueux (diffusion de quantité de
mouvement). Il tend à \emph{homogénéiser} le profil de vitesse, comme la chaleur
diffuse des zones chaudes vers les zones froides.

\subsection{Conditions aux limites}

Les équations de Navier-Stokes sont des équations aux dérivées partielles du second
ordre en espace~: elles nécessitent des conditions aux limites sur tout le bord du
domaine fluide.

\medskip
\noindent\textbf{À la paroi solide} (peau du Rafale, paroi d'une canalisation)~:
la condition de non-glissement exprime que le fluide adhère à la paroi.
Contrairement au modèle parfait qui autorisait le glissement, ici~:
\[
    \bm{U}_\text{fluide} = \bm{U}_\text{paroi}
    \qquad\text{(condition de non-glissement)}
\]
C'est l'existence de cette couche de fluide \emph{immobile} contre la paroi qui
génère le gradient de vitesse $\partial U/\partial n$ et donc les contraintes
visqueuses.

\medskip
\noindent\textbf{À l'interface entre deux fluides} (surface libre eau/air, interface
entre deux huiles)~: on impose la continuité de la vitesse et la continuité du
vecteur contrainte~:
\[
    \bm{U}_A = \bm{U}_B,
    \qquad
    \ten{\sigma}_A\cdot\bm{n} = \ten{\sigma}_B\cdot\bm{n},
    \qquad
    \bm{U}_S\cdot\bm{n} = \bm{U}\cdot\bm{n}
\]

\medskip
\noindent\textbf{Pression effective.} Il est souvent commode de travailler avec la
\emph{pression effective} $p_\text{eff}=p-\rho\bm{g}\cdot\bm{x}$, qui absorbe
les effets de la gravité dans le terme de pression et simplifie les équations en
régimes horizontaux.

% ==========================================
\section{Bilan de l'énergie cinétique}
% ==========================================

\subsection{Dérivation depuis la conservation de la quantité de mouvement}

Pour comprendre les pertes de charge et le travail des machines, on a besoin d'un
bilan d'énergie mécanique. Ce bilan n'est pas une nouvelle hypothèse~: on l'obtient
en prenant le produit scalaire de l'équation de Newton avec la vitesse locale $\bm{U}$.

En effet, $\bm{U}\cdot(m\bm{a}) = \bm{U}\cdot\bm{F} = \text{puissance des forces}$
est le théorème travail-puissance de la mécanique du point. Son analogue pour une
particule fluide de densité $\rho$ est~:

\[
    \bm{U}\cdot\left\{\rho\DD{\bm{U}} = \nabla\cdot\ten{\sigma}+\rho\bm{g}\right\}
    \quad\Longrightarrow\quad
    \rho\,\DD{}\!\left(\frac{U^2}{2}\right) = \bm{U}\cdot\left(\nabla\cdot\ten{\sigma}+\rho\bm{g}\right)
\]

On applique ensuite le théorème de Reynolds à l'énergie cinétique $\chi=U^2/2$
sur un domaine fixe $\mathcal{D}$ de surface $\mathcal{S}$~:

\[
    \frac{d}{dt}\int_{\mathcal{D}}\rho\frac{U^2}{2}\,dv
    = \underbrace{\int_{\mathcal{D}}\bm{U}\cdot\left(\nabla\cdot\ten{\sigma}+\rho\bm{g}\right)dv}_{=(a)}
    - \int_{\mathcal{S}}\rho\frac{U^2}{2}\,\bm{U}\cdot\bm{n}\,ds
\]

Pour transformer le terme (a), on utilise l'identité vectorielle (que l'on
démontre par développement en indices, en exploitant la symétrie de $\ten{\sigma}$)~:
\[
    \bm{U}\cdot(\nabla\cdot\ten{\sigma})
    = \nabla\cdot(\bm{U}\cdot\ten{\sigma}) - \ten{\sigma}:\ten{D}
\]

où le \textbf{double produit contracté} $\ten{\sigma}:\ten{D}=\sigma_{ij}D_{ij}$
est un scalaire obtenu par multiplication terme à terme et sommation.
On peut alors séparer (a) en une puissance des forces extérieures $P_\text{ext}$
et une puissance des forces intérieures $P_\text{int}$~:

\[
    (a) = \underbrace{\int_{\mathcal{S}}\bm{U}\cdot(\ten{\sigma}\cdot\bm{n})\,ds
          + \int_{\mathcal{D}}\rho\bm{U}\cdot\bm{g}\,dv}_{P_\text{ext}}
          \underbrace{- \int_{\mathcal{D}}\ten{\sigma}:\ten{D}\,dv}_{P_\text{int}}
\]

On obtient finalement~:

\begin{importantbox}{Bilan intégral d'énergie cinétique}
\[
    \frac{d}{dt}\int_{\mathcal{D}}\rho\frac{U^2}{2}\,dv
    = P_\text{ext} + P_\text{int} - \int_{\mathcal{S}}\rho\frac{U^2}{2}\,\bm{U}\cdot\bm{n}\,ds
    \tag{13}
\]
La variation d'énergie cinétique dans $\mathcal{D}$ est alimentée par la puissance
des forces extérieures et intérieures, et diminuée par le flux d'énergie cinétique
sortant.
\end{importantbox}

\subsection{Dissipation visqueuse}

Pour un fluide newtonien incompressible, le terme $P_\text{int}$ se simplifie.
En substituant $\ten{\sigma}=-p\ten{I}+\ten{\tau}$ et en notant que
$(-p\ten{I}):\ten{D} = -p\,\text{tr}(\ten{D}) = -p\nabla\cdot\bm{U}=0$
(incompressible), on obtient~:
\[
    -P_\text{int} = \int_{\mathcal{D}}\ten{\sigma}:\ten{D}\,dv
    = \int_{\mathcal{D}}\ten{\tau}:\ten{D}\,dv
\]

En substituant $\ten{\tau}=2\mu\ten{D}$~:

\begin{importantbox}{Dissipation visqueuse}
\[
    \mathcal{D}_v \equiv \int_{\mathcal{D}}\ten{\tau}:\ten{D}\,dv
    = \int_{\mathcal{D}} 2\mu\ten{D}:\ten{D}\,dv
    = \int_{\mathcal{D}} \frac{\mu}{2}\left(\pd{U_i}{x_j}+\pd{U_j}{x_i}\right)^2 dv \geq 0
\]
La dissipation visqueuse $\mathcal{D}_v \geq 0$ est toujours \textbf{positive}~:
c'est une transformation irréversible d'énergie mécanique en chaleur par
frottements visqueux. Elle est nulle dans un fluide parfait ($\mu=0$).
\end{importantbox}

La positivité de $\mathcal{D}_v$ traduit le second principe de la thermodynamique~:
les frottements ne créent pas d'énergie mécanique, ils la dégradent en agitation
thermique. Pour le Rafale, la dissipation visqueuse dans la couche limite représente
une part non négligeable de la traînée totale~: même si elle est petite comparée à
l'énergie cinétique de l'écoulement, c'est ce terme qui oblige les moteurs à
tourner en croisière.

\subsection{Illustration~: écoulement de Poiseuille en canal plan}

L'écoulement de Poiseuille entre deux plaques parallèles (distantes de $2h$,
écoulement selon $x_1$) admet le profil parabolique~:
\[
    \bm{U} = \left(\frac{U_0}{h^2}(h^2-x_2^2),\;0,\;0\right)
\]
Les tenseurs correspondants sont~:

\[
    \ten{D} = \begin{pmatrix} 0 & -\frac{U_0 x_2}{h^2} & 0 \\[4pt] -\frac{U_0 x_2}{h^2} & 0 & 0 \\ 0 & 0 & 0 \end{pmatrix},
    \quad
    \ten{\tau} = 2\mu\ten{D},
    \quad
    \ten{\sigma} = \begin{pmatrix} -p & -\frac{2\mu U_0 x_2}{h^2} & 0 \\[4pt] -\frac{2\mu U_0 x_2}{h^2} & -p & 0 \\ 0 & 0 & -p \end{pmatrix}
\]

Sur la paroi supérieure ($x_2=h$, normale $\bm{n}=(0,1,0)$), le vecteur contrainte est~:
\[
    \bm{T} = \ten{\sigma}\cdot\bm{n} = \begin{pmatrix} -2\mu U_0/h \\ -p \\ 0 \end{pmatrix}
\]
La force visqueuse par unité de surface exercée par le fluide sur la paroi supérieure
vaut $\bm{F}=2(\mu U_0/h)\bm{e}_1$, et par symétrie la force totale (deux parois) est
$4\mu(U_0/h)\bm{e}_1$.

La dissipation visqueuse intégrée sur la hauteur du canal~:
\[
    \int_{-h}^{+h}\ten{\tau}:\ten{D}\,dx_2
    = \int_{-h}^{+h}4\mu\left(\frac{U_0 x_2}{h^2}\right)^2 dx_2
    = \frac{8\mu U_0^2}{3h} \quad \text{(W.m}^{-2}\text{)}
\]

\textbf{Résultat clé~:} la dissipation visqueuse est proportionnelle à $U_0^2$.
Doubler la vitesse dans une conduite quadruple les pertes par frottement.
Pour le Rafale, cela signifie que voler à Mach~1{,}6 au lieu de Mach~0{,}8
ne consomme pas deux fois plus pour les frottements de peau~: il en consomme
quatre fois plus.

% ==========================================
\section{Énergétique d'un système à flux continu}
% ==========================================

\subsection{Bilan d'énergie mécanique}

La plupart des systèmes industriels---tuyauteries, pompes, turbines---sont des
\emph{systèmes à flux continu}~: le fluide entre par une section $\mathcal{S}_1$,
traverse un domaine $\mathcal{D}$ (éventuellement muni d'une machine), et sort
par $\mathcal{S}_2$.

\begin{center}
\begin{tikzpicture}[scale=1.0]
    % domaine
    \draw[thick, blue!50!black, rounded corners=8pt]
        (0,0.8) -- (0,2.2) -- (4,2.2) -- (4,0.8) -- (3.2,0.8)
        arc[start angle=90, end angle=-90, radius=0.5]
        -- (0.8,0.8) arc[start angle=-90, end angle=90, radius=0.5]
        -- cycle;
    \node[blue!50!black] at (2,1.5) {$\mathcal{D}$};
    % machine
    \draw[fill=gray!30] (2.1,0.5) ellipse (0.5 and 0.6);
    \draw[thick, dashed] (2.1,1.1) -- (2.1,2.5);
    \node[above] at (2.1,2.5) {$P_m$};
    \draw[->, thick] (2.1,2.5) -- (2.1,2.1);
    % sections
    \draw[fill=blue!20] (-0.1,0.8) rectangle (0,2.2);
    \node[left] at (0,1.5) {$\mathcal{S}_1$};
    \draw[->,blue!70!black, thick] (-0.5,1.5) -- (0.3,1.5);
    \draw[fill=blue!20] (4,0.8) rectangle (4.1,2.2);
    \node[right] at (4.1,1.5) {$\mathcal{S}_2$};
    \draw[->,blue!70!black, thick] (3.7,1.5) -- (4.5,1.5);
    % paroi
    \node[below] at (2.1,0.0) {$\mathcal{S}_w$ (paroi fixe)};
\end{tikzpicture}
\end{center}

Depuis le bilan d'énergie cinétique, en utilisant que la gravité est conservative
($\bm{g}=\nabla(\bm{g}\cdot\bm{x})$) pour absorber l'énergie potentielle~:

\[
    \frac{d}{dt}\int_{\mathcal{D}}\rho\left(\frac{U^2}{2}-\bm{g}\cdot\bm{x}\right)dv
    = P_m - \mathcal{D}_v
    - \int_{\mathcal{S}_1\cup\mathcal{S}_2}\rho\left(\frac{U^2}{2}-\bm{g}\cdot\bm{x}\right)\bm{U}\cdot\bm{n}\,ds
\]

En développant $P_\text{ext}^{\mathcal{S}}$ et en regroupant la pression dans
l'intégrale surfacique, on fait apparaître la densité d'énergie de Bernoulli
$\mathcal{H}=p-\rho\bm{g}\cdot\bm{x}+\rho U^2/2$~:

\[
    \frac{d}{dt}\int_{\mathcal{D}}\rho\left(\frac{U^2}{2}-\bm{g}\cdot\bm{x}\right)dv
    = P_m - \mathcal{D}_v
    - \int_{\mathcal{S}_1\cup\mathcal{S}_2}\mathcal{H}\,\bm{U}\cdot\bm{n}\,ds
\]

En régime permanent, la dérivée temporelle s'annule. En prenant la
\textbf{moyenne temporelle} (pour traiter indifféremment les régimes laminaires et
turbulents), on obtient~:

\begin{importantbox}{Bilan d'énergie mécanique --- régime permanent (moyenne temporelle)}
\[
    \bar{P}_m = \int_{\mathcal{S}_1\cup\mathcal{S}_2}\overline{\mathcal{H}\bm{U}\cdot\bm{n}}\,ds
    + \bar{\mathcal{D}}_v
\]
La puissance fournie à l'écoulement par la machine alimente deux termes~: la
variation du flux d'énergie entre l'entrée $\mathcal{S}_1$ et la sortie $\mathcal{S}_2$,
et la dissipation visqueuse $\bar{\mathcal{D}}_v \geq 0$.
\end{importantbox}

\subsection{Perte de charge}

Pour modéliser simplement l'intégrale de flux, on introduit la
\textbf{charge hydraulique} $H$ (attention~: $H \neq \mathcal{H}$), définie par~:
\[
    \bar{Q}_v H \equiv \int_{\mathcal{S}}\overline{\mathcal{H}\,\bm{U}\cdot\bm{n}}\,ds
\]

$H$ est une valeur moyennée en temps et en espace de $\mathcal{H}/Q_v$
sur la section~: c'est l'énergie mécanique transportée par unité de volume
de fluide. Le bilan d'énergie s'écrit alors sous forme compacte~:

\begin{importantbox}{Bilan énergétique d'un système à flux continu}
\[
    \bar{P}_m + \bar{Q}_v(H_1 - H_2) = \bar{\mathcal{D}}_v
    \tag{17}
\]
\end{importantbox}

Dans une conduite simple sans machine ($\bar{P}_m=0$), on a nécessairement
$\bar{Q}_v(H_1-H_2) = \bar{\mathcal{D}}_v > 0$, c'est-à-dire $H_1>H_2$~:
la charge diminue dans le sens de l'écoulement. Cette perte de charge est
la manifestation macroscopique de la dissipation visqueuse~: l'énergie mécanique du
fluide est dégradée irréversiblement en chaleur par frottements contre les parois.

\subsection{Rendement des machines hydrauliques}

Lorsqu'une machine est présente, on introduit son \textbf{rendement} $\eta$
pour relier la puissance mécanique de l'arbre $P_m$ à la puissance réellement
transmise au fluide.

\medskip
\textbf{Pompe (ou ventilateur, compresseur)~:} la machine fournit de l'énergie au
fluide pour l'élever en charge ($H_2>H_1$)~:
\[
    \bar{Q}_v(H_2-H_1) = \eta P_m > 0, \qquad \bar{Q}_v\Delta H_P = \eta P_m
\]

\textbf{Turbine~:} le fluide cède de l'énergie à la machine ($H_1>H_2$)~:
\[
    P_T = -P_m = \eta\,\bar{Q}_v(H_1-H_2) > 0, \qquad P_T = \eta Q_v\Delta H_T
\]

\subsection{Application~: centrale hydroélectrique de Malgovert (barrage de Tignes)}

\begin{center}
\begin{tikzpicture}[scale=0.9]
    % barrage
    \fill[blue!40] (0,4) rectangle (2,5);
    \draw[thick] (0,4) -- (2,4) -- (2,5) -- (0,5);
    \node[left] at (-0.1,4.5) {\small barrage de Tignes};
    \node[right] at (2.1,4.7) {\small $\mathcal{H}_D = p_0+\rho g z_D$};
    % conduite forcée
    \draw[thick, blue!70!black] (2,4.2) -- (5,1.5);
    \draw[thick, blue!70!black] (2,4.0) -- (5,1.3);
    \node[below left] at (3.2,2.9) {\small conduite forcée $\Delta H_L$};
    % turbine
    \draw[fill=gray!30] (5,1.3) rectangle (5.6,1.5);
    \node[below] at (5.3,1.3) {\small turbine $\Delta H_T$};
    \draw[->,thick] (5.3,1.5) -- (5.3,2.2);
    \node[right] at (5.3,1.9) {\small $\sim$};
    % sortie
    \fill[blue!40] (5.6,1.0) rectangle (7,1.5);
    \node[right] at (7.0,1.3) {\small $\mathcal{H}_R = p_0+\rho g z_R$};
    % altitude z
    \draw[->,gray] (-0.3,1.0) -- (-0.3,5.3);
    \node[left,gray] at (-0.3,3.1) {$z$};
\end{tikzpicture}
\end{center}

Le bilan de charge entre le réservoir amont (Tignes, altitude $z_D$) et le canal de
fuite (Malgovert, altitude $z_R$) s'écrit~:
\[
    H_D - H_R = \rho g(z_D - z_R) = \Delta H_L + \Delta H_T + \Delta H_O
\]

où $\Delta H_L$ et $\Delta H_O$ sont les pertes de charge dans la conduite forcée
et à la sortie, et $\Delta H_T$ est la charge récupérée par la turbine. La puissance
électrique produite est~:
\[
    P_T = \eta_T Q_v \Delta H_T = \eta_T Q_v\left[\rho g(z_D-z_R) - \Delta H_L - \Delta H_O\right]
\]

\textit{Remarque~:} à la surface du barrage ($x_2=z_D$), la vitesse $U_D$ est
négligeable (énorme surface $\mathcal{S}$, petit débit $Q_v=U_D\times\mathcal{S}$),
donc $\mathcal{H}_D\approx p_0+\rho g z_D$ (la contribution cinétique $\rho U_D^2/2$
est infime). C'est la même approximation que lors du calcul Bernoulli de la
vidange d'un réservoir.

L'énergie est donc essentiellement de l'énergie potentielle gravitationnelle~:
le barrage de Tignes (situé à environ 1800~m d'altitude) alimente la centrale
de Malgovert (environ 900~m) via une conduite forcée de plus de 4~km.
La puissance installée de la centrale est d'environ 180~MW~: c'est le poids de
l'eau qui fait tourner la turbine.

% ==========================================
\section*{Résumé du chapitre}
\addcontentsline{toc}{section}{Résumé du chapitre}
% ==========================================

\begin{importantbox}{Ce qu'il faut retenir}
\textbf{Loi constitutive newtonienne}~:
\[
    \ten{\sigma} = -p\ten{I} + 2\mu\ten{D},
    \qquad
    D_{ij} = \tfrac{1}{2}\!\left(\pd{U_i}{x_j}+\pd{U_j}{x_i}\right)
\]

\vspace{4pt}
\textbf{Équations de Navier-Stokes} (incompressible)~:
\[
    \nabla\cdot\bm{U}=0,
    \qquad
    \rho\DD{\bm{U}} = -\nabla p + \mu\nabla^2\bm{U} + \rho\bm{g}
\]

\vspace{4pt}
\textbf{Dissipation visqueuse} (toujours $\geq 0$)~:
\[
    \mathcal{D}_v = \int_{\mathcal{D}}2\mu\ten{D}:\ten{D}\,dv \geq 0
\]

\vspace{4pt}
\textbf{Perte de charge en système à flux continu}~:
\[
    \bar{P}_m + \bar{Q}_v(H_1-H_2) = \bar{\mathcal{D}}_v
\]
La charge $H$ diminue toujours dans le sens de l'écoulement en l'absence
de machine. La viscosité prend, la machine rend.

\vspace{4pt}
\textbf{Rendement}~: pompe $Q_v\Delta H_P = \eta P_m$~;
turbine $P_T = \eta Q_v \Delta H_T$.
\end{importantbox}

% ==========================================
\chapter{Analyse dimensionnelle --- nombre de Reynolds}
% ==========================================

\noindent
\textit{Fil conducteur --- le Rafale Marine.}
La soufflerie F1 de l'ONERA à Fauga-Mauzac peut accueillir une maquette d'avion
à l'échelle~$1/10$. La question que pose l'ingénieur est simple~: les données
mesurées sur la maquette sont-elles directement transposables au Rafale réel~?
Plus généralement, peut-on identifier de manière \emph{systématique} les quelques
paramètres adimensionnels qui gouvernent la physique d'un écoulement, sans avoir
à réaliser des milliers d'expériences en faisant varier indépendamment chaque
grandeur physique~?

La réponse est oui, et elle repose sur un théorème profond~: le
\textbf{théorème $\Pi$} de Vaschy-Buckingham. Ce chapitre en montre la puissance,
en dégage le nombre adimensionnel central de la mécanique des fluides --- le
\textbf{nombre de Reynolds} --- et l'applique aux conduites (Poiseuille).

% ==========================================
\section{Analyse dimensionnelle}
% ==========================================

\subsection{Unités et dimensions}

Avant de construire des nombres sans dimension, il faut clarifier ce qu'est
une dimension. La physique reconnaît quatre dimensions primitives en mécanique
des fluides~:

\begin{center}
\renewcommand{\arraystretch}{1.3}
\begin{tabular}{lcc}
\toprule
\textbf{Grandeur} & \textbf{Dimension} & \textbf{Unité SI} \\
\midrule
Masse      & $[M]$ & kg \\
Longueur   & $[L]$ & m  \\
Temps      & $[T]$ & s  \\
Température& $[\Theta]$ & K \\
\bottomrule
\end{tabular}
\end{center}

\noindent Le choix des unités est arbitraire (on pourrait travailler en pieds,
en livres et en secondes), mais les \emph{dimensions} sont universelles.
Toute grandeur physique s'exprime comme un produit de puissances de ces quatre
dimensions de base~:

\[
[V] = LT^{-1}, \quad
[F] = MLT^{-2}, \quad
[p] = ML^{-1}T^{-2}, \quad
[\nu] = L^2T^{-1}, \quad
[\mu] = ML^{-1}T^{-1}
\]

\subsection{Théorème $\Pi$ de Vaschy-Buckingham}

Considérons un problème physique impliquant $n$ grandeurs dimensionnelles
$a_1, a_2, \ldots, a_n$. On suppose qu'il existe une relation implicite
\[
    \mathcal{F}(a_1, a_2, \ldots, a_n) = 0
\]
entre toutes ces grandeurs (la traînée d'un avion, par exemple, \emph{doit}
dépendre d'un certain nombre de paramètres physiques, même si la forme exacte
de la loi est inconnue).

Chaque grandeur $a_i$ possède une dimension de la forme
$[a_i] = L^{\alpha_i}T^{\beta_i}M^{\gamma_i}\Theta^{\delta_i}$.
On cherche les combinaisons $a_1^{\lambda_1}\cdots a_n^{\lambda_n}$
qui sont sans dimension, c'est-à-dire telles que
$[a_1]^{\lambda_1}\cdots[a_n]^{\lambda_n} = L^0T^0M^0\Theta^0$.
Cela revient à résoudre le système linéaire~:

\[
\begin{pmatrix} \alpha_1 & \cdots & \alpha_n \\ \beta_1 & & \beta_n \\ \gamma_1 & & \gamma_n \\ \delta_1 & \cdots & \delta_n \end{pmatrix}
\begin{pmatrix} \lambda_1 \\ \vdots \\ \lambda_n \end{pmatrix} = \mathbf{0}
\]

Notons $r$ le rang de cette matrice dimensionnelle~: le noyau du système est
de dimension $n-r$, ce qui signifie qu'il existe exactement $n-r$ combinaisons
adimensionnelles indépendantes $\Pi_1, \ldots, \Pi_{n-r}$.

\begin{importantbox}{Théorème Pi --- Vaschy (1892) et Buckingham (1914)}
Tout problème physique impliquant $n$ grandeurs et de rang dimensionnel $r$
peut être entièrement décrit par une loi reliant $n-r$ paramètres
adimensionnels~:
\[
    \Pi_1 = \mathcal{F}_0(\Pi_2, \ldots, \Pi_{n-r})
\]
où $\Pi_1$ est le nombre sans dimension associé à la grandeur d'intérêt.
Cas particulier~: si $n - r = 1$, alors $\Pi_1 = \mathrm{cst}$.
\end{importantbox}

La force du théorème est de réduire drastiquement l'espace paramétrique~:
au lieu de faire varier $n$ paramètres dimensionnels indépendamment,
il suffit d'explorer $n-r$ groupes adimensionnels.

\subsection{Application~: fréquence de détachement tourbillonnaire}

\textit{Problème.} Derrière un cylindre placé dans un écoulement uniforme
$U_\infty$, des tourbillons se détachent alternativement de chaque côté
(allée de Von Kármán). Quelle est la fréquence $f$ de ce détachement~?

Les paramètres physiques pertinents sont~: $f$, $U_\infty$, $\rho$, $\mu$, $D$
(diamètre du cylindre). Soit $n=5$. La matrice dimensionnelle est~:

\begin{center}
\renewcommand{\arraystretch}{1.2}
\begin{tabular}{c|ccccc}
 & $U_\infty$ & $\rho$ & $\mu$ & $D$ & $f$ \\
\hline
$L$ & $1$ & $-3$ & $-1$ & $1$ & $0$ \\
$T$ & $-1$ & $0$ & $-1$ & $0$ & $-1$ \\
$M$ & $0$ & $1$ & $1$ & $0$ & $0$ \\
$\Theta$ & $0$ & $0$ & $0$ & $0$ & $0$ \\
\end{tabular}
\end{center}

Le rang de cette matrice est $r=3$, donc $n-r=2$ paramètres adimensionnels
à construire. On obtient naturellement~:

\begin{importantbox}{Nombre de Strouhal et nombre de Reynolds}
\[
\Pi_1 = \frac{fD}{U_\infty} \equiv \mathrm{St} \quad \text{(nombre de Strouhal)},
\qquad
\Pi_2 = \frac{\rho U_\infty D}{\mu} \equiv \mathrm{Re} \quad \text{(nombre de Reynolds)}
\]
\[
\Longrightarrow \quad \mathrm{St} = \mathcal{F}_0(\mathrm{Re})
\]
\end{importantbox}

Ce résultat est remarquable~: la fréquence de détachement, qui dépend
\textit{a priori} de cinq paramètres, ne dépend en réalité que d'un seul
nombre adimensionnel, Re. Les expériences confirment~:

\[
\mathrm{St} \approx 0.198\left(1 - \frac{19.7}{\mathrm{Re}}\right)
\quad \text{(Lienhard 1966)}
\]

Pour les antennes radar d'un Rafale en vol à basse altitude --- profil de mât
de diamètre $D\sim 2$~cm, $U_\infty\sim 100$~m/s~--- on obtient
$\mathrm{Re} \approx 1.3\times10^5$ et $f \approx \mathrm{St}\cdot U_\infty/D \approx 1$~kHz.
Si cette fréquence coïncide avec un mode propre de la structure, on risque
la résonance~: c'est un problème réel en conception aéronautique.

% ==========================================
\section{Nombre de Reynolds}
% ==========================================

\subsection{Définition et interprétation physique}

Le nombre de Reynolds est \emph{le} paramètre adimensionnel central de la
mécanique des fluides des fluides incompressibles. Il est défini à partir
d'une vitesse caractéristique $U$ et d'une longueur caractéristique $L$
de l'écoulement~:

\begin{defbox}{Nombre de Reynolds}
\[
    \mathrm{Re} = \frac{UL}{\nu} = \frac{\rho UL}{\mu}
\]
$U$~: vitesse caractéristique (m/s),
$L$~: longueur caractéristique (m),
$\nu$~: viscosité cinématique (m$^2$/s).
\end{defbox}

Pour comprendre ce que Re mesure physiquement, regardons les ordres de grandeur
des termes dans les équations de Navier-Stokes~:

\[
\underbrace{\pd{\bm{U}}{t} + \bm{U}\cdot\nabla\bm{U}}_{\sim\,U^2/L\;(\text{inertie})}
= -\frac{1}{\rho}\nabla p + \underbrace{\nu\nabla^2\bm{U}}_{\sim\,\nu U/L^2\;(\text{viscosité})} + \bm{g}
\]

Le rapport de ces deux termes est~:
\[
    \frac{\text{terme convectif (inertiel)}}{\text{terme visqueux}}
    \sim \frac{U^2/L}{\nu U/L^2} = \frac{UL}{\nu} = \mathrm{Re}
\]

Re est donc le rapport des forces inertielles sur les forces visqueuses.
On peut aussi l'interpréter comme un rapport de temps~:
\[
    \mathrm{Re} = \frac{L^2/\nu}{L/U}
    = \frac{\text{temps de diffusion visqueuse}}{\text{temps de convection}}
\]

\medskip
\noindent Deux régimes extrêmes~:

\medskip
$\mathrm{Re} \ll 1$ --- \textbf{écoulement de Stokes} (ou rampant)~: la viscosité
domine, le terme convectif est négligeable. C'est le régime des microorganismes
nageant dans l'eau~: une bactérie s'arrête instantanément si elle cesse de
battre ses flagelles.

\medskip
$\mathrm{Re} \gg 1$ --- \textbf{écoulement quasi-parfait}~: la viscosité est
négligeable sauf dans des couches minces près des parois (couches limites,
chapitre~4). C'est le régime du Rafale en vol~:

\[
    \mathrm{Re}_\text{Rafale} = \frac{U_\infty \cdot c}{\nu}
    \approx \frac{200\,\text{m/s} \times 5.3\,\text{m}}{1.5\times10^{-5}\,\text{m}^2/\text{s}}
    \approx 7\times10^7
\]

À ce nombre de Reynolds, l'écoulement autour de l'aile est très loin d'être
laminaire~: la couche limite est turbulente sur presque toute la corde.

\subsection{Expérience de Reynolds (1883)}

Osborne Reynolds a mis en évidence expérimentalement la transition
laminaire-turbulent dans une conduite en injectant un filet de colorant~:
\begin{itemize}
    \item À faible $U$ (Re petit)~: le filet reste parfaitement droit --- écoulement \textbf{laminaire}.
    \item Au-dessus d'un seuil Re\,$\approx 2300$~: le filet se mélange brusquement
    à l'eau environnante --- écoulement \textbf{turbulent}.
\end{itemize}
Cette expérience, publiée en 1883 dans les \textit{Phil.\ Trans.\ Roy.\ Soc.},
est l'une des plus importantes de l'histoire de la mécanique des fluides.

\subsection{Équation de Navier-Stokes adimensionnée}

En adimensionnant les variables par les échelles caractéristiques du problème~:
\[
    \tilde{U}_i = \frac{U_i}{U_\infty}, \quad
    \tilde{p} = \frac{p}{\rho U_\infty^2}, \quad
    \tilde{x}_i = \frac{x_i}{L}, \quad
    \tilde{t} = \frac{t}{L/U_\infty}
\]

l'équation de Navier-Stokes devient~:

\begin{importantbox}{NS adimensionnée --- Re seul paramètre gouvernant}
\[
    \frac{\partial \tilde{\bm{U}}}{\partial \tilde{t}}
    + \tilde{\bm{U}}\cdot\tilde{\nabla}\tilde{\bm{U}}
    = -\tilde{\nabla}\tilde{p}
    + \frac{1}{\mathrm{Re}}\tilde{\nabla}^2\tilde{\bm{U}},
    \qquad \tilde{\nabla}\cdot\tilde{\bm{U}} = 0
\]
avec Re $= U_\infty L/\nu$.
\end{importantbox}

Ce résultat est fondamental~: une fois adimensionné, l'écoulement autour d'un
obstacle de géométrie donnée ne dépend que d'un seul paramètre, Re.
Deux expériences de géométrie similaire ayant le même Re produisent
\emph{exactement} le même champ de vitesse adimensionnel.
C'est le \textbf{principe de similitude dynamique}~: c'est lui qui justifie
les essais en soufflerie.

% ==========================================
\section{Écoulement autour d'un obstacle}
% ==========================================

\subsection{Coefficient de traînée}

Appliquons le théorème $\Pi$ à la force de traînée $F_D$ exercée par un
écoulement sur un obstacle de taille caractéristique $L$. Les paramètres
sont $F_D,\,\mu,\,\rho,\,U_\infty,\,L$~; $n=5$, rang $r=3$, donc $n-r=2$.

On construit les deux groupes adimensionnels~:

\[
    \Pi_1 = \frac{F_D}{q_\infty S} \equiv C_D,
    \qquad
    \Pi_2 = \frac{\rho U_\infty L}{\mu} \equiv \mathrm{Re}
\]

où $q_\infty = \frac{1}{2}\rho U_\infty^2$ est la pression dynamique
(déjà rencontrée au tube de Pitot) et $S \propto L^2$ est la surface
de référence (maître-couple).

\begin{defbox}{Coefficient de traînée}
\[
    C_D \equiv \frac{F_D}{\dfrac{1}{2}\rho U_\infty^2 S}
    \qquad \Longrightarrow \qquad
    C_D = C_D(\mathrm{Re})
\]
La traînée est entièrement déterminée par Re pour une géométrie donnée.
\end{defbox}

Pour une sphère lisse de diamètre $D$ ($S=\pi D^2/4$)~:
\begin{itemize}
    \item $\mathrm{Re}_D \ll 1$ (Stokes)~: $C_D = 24/\mathrm{Re}_D$ --- la traînée
    est due entièrement à la viscosité.
    \item $10^2 \lesssim \mathrm{Re}_D \lesssim 10^5$~: $C_D \approx 0.4$ --- la traînée
    est dominée par la dépression dans le sillage (recirculation).
    \item $\mathrm{Re}_D^c \approx 3\times10^5$~: \textbf{crise de traînée} --- la couche
    limite passe de laminaire à turbulente, la séparation recule et $C_D$ chute
    brutalement de $\sim0.4$ à $\sim0.1$.
\end{itemize}

La crise de traînée explique pourquoi les balles de golf sont texturées~:
les alvéoles déclenchent la transition turbulente de la couche limite à
$\mathrm{Re}$ modéré, réduisant la traînée et augmentant la portée.

\subsection{Maquette et similitude en soufflerie}

Pour mesurer correctement $C_D = C_D(\mathrm{Re}')$ du Rafale à l'échelle~$1$
sur une maquette de corde $c < c'$, il faut reproduire le nombre de Reynolds
réel~:

\[
    \mathrm{Re}' = \frac{\rho'_\infty U'_\infty c'}{\mu'_\infty}
    \overset{!}{=}
    \mathrm{Re} = \frac{\rho_\infty U_\infty c}{\mu_\infty}
\]

En aéronautique, $U_\infty \approx U'_\infty$ (le nombre de Mach doit aussi être
conservé), mais $c < c'$~: il faut donc augmenter $\rho_\infty/\mu_\infty$ en
soufflerie. Deux stratégies~:

\medskip
\noindent \textbf{Soufflerie pressurisée} ($p_\infty \nearrow$)~: en utilisant
$\rho_\infty = p_\infty/(rT_\infty)$, augmenter $p_\infty$ augmente $\rho_\infty$
et donc Re. Inconvénient~: la pression dynamique $q_\infty = \frac{\gamma}{2}p_\infty M^2$
augmente aussi, ce qui génère des contraintes mécaniques élevées sur la maquette.

\medskip
\noindent \textbf{Soufflerie cryogénique} ($T_\infty \searrow$)~: abaisser $T_\infty$
diminue $\mu_\infty$ (viscosité dynamique) et augmente $\rho_\infty$, les deux
effets augmentant Re. L'\textbf{European Transonic Windtunnel} (ETW, Cologne)
injecte de l'azote liquide pour atteindre $T_\infty$ entre 110 et 313~K, à des
pressions de 1.25 à 4.5~bar, permettant des nombres de Reynolds de corde
jusqu'à $80\times10^6$ --- comparable aux conditions de vol réelles des
grands avions de ligne.

\medskip
La viscosité dynamique de l'air suit la \textbf{loi de Sutherland}~:
\[
    \mu(T) \approx \mu_0\left(\frac{T}{T_0}\right)^{3/2}\frac{T_0+T_s}{T+T_s}
\]
avec $T_0=273$~K, $T_s=111$~K, $\mu_0=1.716\times10^{-5}$~kg/(m$\cdot$s).
Ainsi $\mu$ augmente avec $T$ (contrairement aux liquides~!), ce qui signifie
que l'air froid est moins visqueux~: c'est l'effet recherché en soufflerie
cryogénique.

% ==========================================
\section{Applications du théorème $\Pi$}
% ==========================================

\subsection{Traînée de la coque d'un bateau}

Un navire est soumis à deux contributions de traînée de natures très différentes~:
le frottement visqueux le long de la coque, et la traînée de vagues (ondes de
surface créées par le déplacement du bateau). Les variables pertinentes sont~:
$F_D,\,\mu,\,\rho,\,U_\infty,\,L,\,g$ --- la gravité apparaît car elle régit
la dynamique des ondes de surface. Soit $n=6$, rang $r=3$, donc $n-r=3$.

\[
    \Pi_1 = \frac{F_D}{\frac{1}{2}\rho U^2 S} = C_D, \qquad
    \Pi_2 = \frac{\rho U L}{\mu} = \mathrm{Re}, \qquad
    \Pi_3 = \frac{U}{\sqrt{gL}} = \mathrm{Fr} \quad \text{(Froude)}
\]

On obtient $C_D = C_D(\mathrm{Re}, \mathrm{Fr})$. Le nombre de Froude compare
la vitesse du navire à la vitesse des ondes de surface~: $\mathrm{Fr} > 1$
signifie que le bateau va plus vite que ses propres vagues (régime de
planage). Le sillage en V que font les canards sur un lac a toujours
le même angle ($\approx 39°$)~: c'est une conséquence directe du nombre
de Froude et de la dispersion des ondes de surface.

Puisqu'on ne peut pas simultanément conserver Re et Fr entre la maquette
et le bateau réel (cela imposerait des contraintes incompatibles sur $U$
et $\nu$), on utilise \textbf{l'hypothèse de Froude}~:

\[
    C_D \approx C_F(\mathrm{Re}) + C_R(\mathrm{Fr})
\]

Les deux contributions obéissent à des physiques distinctes~: on les mesure
séparément et on les recombine. En essais sur maquette~: Fr est conservé
(ce qui fixe la vitesse de la maquette), $C_D$ total est mesuré, et $C_F$ est
calculé analytiquement (frottement laminaire ou turbulent connu), ce qui
permet d'extraire $C_R$. À l'échelle~$1$, Re réel est connu, $C_F$ est
recalculé, et on recompose $C_D$ réel.

\subsection{Vitesse d'une embarcation d'aviron}

Pourquoi un équipage de 8 rameurs va-t-il plus vite qu'un équipage de 1~?
L'analyse dimensionnelle apporte une réponse élégante (McMahon, \textit{Sciences}, 1971).

On suppose $N$ rameurs identiques (puissance individuelle $P$, poids et
morphologie similaires), et une géométrie d'embarcation similaire. Le frottement
est dominé par la traînée de peau ($C_D \approx C_F(\mathrm{Re})$, ondes négligeables
pour un bateau profilé dans la plage de vitesse d'aviron).

La puissance fournie doit équilibrer la puissance de traînée~:
\[
    N \times P \propto C_F \cdot \frac{1}{2}\rho U^2 S \cdot U
\]

Or le volume submergé est $\mathcal{V} \sim N$ (un équipage de 8 pèse 8 fois plus
qu'un rameur seul), d'où la surface mouillée $S \sim \mathcal{V}^{2/3} \sim N^{2/3}$.
En supposant $C_F$ approximativement constant (Re varie peu)~:

\[
    N \times P \propto N^{2/3} \cdot U^3
    \quad\Longrightarrow\quad
    U \propto N^{1/9}
\]

Ce résultat prédit que la vitesse croit très lentement avec le nombre de
rameurs~: de 1 à 8 rameurs, le gain de vitesse théorique est
$8^{1/9} \approx 1.28$, soit $+28\%$. Les données des championnats du monde
(Lucerne 2010) confirment remarquablement cette loi en $N^{1/9}$.

% ==========================================
\section{Écoulement de Poiseuille en conduite circulaire}
% ==========================================

\subsection{Solution analytique}

Revenons à un écoulement canonique~: l'écoulement laminaire établi dans une
conduite circulaire de rayon $R$ et de diamètre $D=2R$. Le circuit hydraulique
du Rafale (carburant, fluide de refroidissement, circuit de commande de vol)
est constitué de tels tuyaux~: comprendre la perte de charge qui s'y développe
est indispensable pour dimensionner les pompes.

En coordonnées cylindriques $(r, \theta, z)$ avec $\bm{U}=(0,0,U(r))$,
les équations de Navier-Stokes (gravité négligée, conduite horizontale)
se réduisent à~:

\[
    0 = -\pd{p}{r}, \qquad 0 = -\frac{1}{r}\pd{p}{\theta}, \qquad
    0 = -\pd{p}{z} + \frac{\mu}{r}\frac{d}{dr}\!\left(r\frac{dU}{dr}\right)
\]

Les deux premières équations indiquent que $p = p(z)$ uniquement. La troisième
égalise une fonction de $z$ et une fonction de $r$~: les deux doivent être
égales à une constante $G$~:

\[
    \frac{dp}{dz} = \mu\frac{1}{r}\frac{d}{dr}\!\left(r\frac{dU}{dr}\right) = G = \text{cst}
\]

En intégrant deux fois avec les conditions aux limites $U(R)=0$ (non-glissement)
et $dU/dr|_{r=0}=0$ (symétrie axiale)~:

\begin{importantbox}{Profil de Poiseuille (conduite circulaire)}
\[
    U(r) = -\frac{GR^2}{4\mu}\left[1-\left(\frac{r}{R}\right)^2\right]
    = U_0\left[1-\left(\frac{r}{R}\right)^2\right]
\]
Profil \textbf{parabolique}~; $U_0=-GR^2/(4\mu) > 0$ (car $G<0$~: la pression
décroît dans le sens de l'écoulement).
La pression varie linéairement~: $p = Gz$.
\end{importantbox}

\begin{center}
\begin{tikzpicture}[scale=1.1]
    \draw[thick] (-0.3,1.2) -- (4.5,1.2);
    \draw[thick] (-0.3,-1.2) -- (4.5,-1.2);
    \foreach \x in {-0.3,-0.1,...,0.1} {
        \draw[gray] (\x,1.2) -- (\x-0.2,1.4);
        \draw[gray] (\x,-1.2) -- (\x-0.2,-1.4);
    }
    \foreach \x in {4.3,4.1,3.9} {
        \draw[gray] (\x,1.2) -- (\x+0.2,1.4);
        \draw[gray] (\x,-1.2) -- (\x+0.2,-1.4);
    }
    \foreach \y in {-1.1,-0.7,-0.3,0,0.3,0.7,1.1} {
        \pgfmathsetmacro{\len}{2.2*(1-(\y/1.2)^2)}
        \draw[-{Latex}, blue!70!black] (1.0,\y) -- (1.0+\len,\y);
    }
    \draw[<->, gray] (0.2,-1.2) -- (0.2,1.2);
    \node[left, gray] at (0.2,0) {$D$};
    \node[right] at (3.2,0) {\small $U_0$};
    \node[below] at (2.5,-1.2) {\small $U(r)=U_0[1-(r/R)^2]$};
    \node at (0.7,0) {$z$};
    \draw[->, gray, thin] (0.6,0) -- (0.95,0);
\end{tikzpicture}
\end{center}

\subsection{Vitesse de débit et contrainte pariétale}

La \textbf{vitesse de débit} $U_d$ (vitesse moyenne sur la section) est~:
\[
    U_d \equiv \frac{1}{\pi R^2}\int_0^R U(r)\,2\pi r\,dr = \frac{U_0}{2}
\]
La vitesse maximale (en $r=0$) est le double de la vitesse de débit.

La \textbf{contrainte pariétale} en $r=R$~:
\[
    \tau_w = \mu\left.\frac{dU}{dr}\right|_{r=R} = -\frac{2\mu U_0}{R}
    = -\frac{4\mu U_d}{R}
\]

\subsection{Coefficient de frottement et modèle 1-D}

On définit le \textbf{coefficient de frottement} par analogie avec le
coefficient de traînée~:

\begin{defbox}{Coefficient de frottement en conduite}
\[
    C_f \equiv \frac{|\tau_w|}{\rho U_d^2/2} = \frac{16}{\mathrm{Re}_D}
    \qquad \text{avec} \quad \mathrm{Re}_D = \frac{U_d D}{\nu}
\]
\end{defbox}

Ce résultat est valable en régime laminaire ($\mathrm{Re}_D \lesssim 2300$).
En régime turbulent ($\mathrm{Re}_D \gtrsim 4000$), la loi de Blasius donne
$C_f \approx 0.0791\,\mathrm{Re}_D^{-1/4}$, et à très grand Reynolds la
relation implicite de Prandtl-Kármán~:
\[
    \frac{1}{\sqrt{C_f}} \approx 3.86\log_{10}\!\left(\mathrm{Re}_D\sqrt{C_f}\right) - 0.088
\]

Relions maintenant $C_f$ à la perte de charge macroscopique.
En appliquant le bilan d'énergie du chapitre~2 (éq.~17) à la conduite
sans machine~:

\[
    Q_v(H_1 - H_2) = \mathcal{D}_v
\]

Seule la pression contribue à la variation de charge $H$ (les termes cinétiques
sont identiques à l'entrée et à la sortie d'une conduite de section constante).
Le calcul donne~:

\begin{importantbox}{Perte de charge en conduite --- modèle 1-D}
\[
    \Delta H = [p]^1_2 = -GL = \frac{4\mu U_0}{R^2} L
    = \underbrace{\frac{64}{\mathrm{Re}_D}}_{\psi(\mathrm{Re}_D)}
    \cdot\frac{L}{D}\cdot\frac{1}{2}\rho U_d^2
\]

Plus généralement, pour tout régime d'écoulement en conduite~:
\[
    \Delta H = \psi(\mathrm{Re}_D)\,\frac{L}{D}\,\frac{1}{2}\rho U_d^2
\]
avec $\psi = 64/\mathrm{Re}_D$ (laminaire) ou $\psi \approx 0.316\,\mathrm{Re}_D^{-1/4}$
(turbulent --- Blasius).
\end{importantbox}

Cette formule est directement utilisable pour dimensionner les circuits de
carburant du Rafale. Supposons un tube de rayon $R=5$~mm, longueur $L=1$~m,
débit volumique $Q_v=10^{-4}$~m$^3$/s de kérosène ($\rho\approx800$~kg/m$^3$,
$\nu\approx 2\times10^{-6}$~m$^2$/s)~:
\[
    U_d = \frac{Q_v}{\pi R^2} \approx 1.27\,\text{m/s},
    \quad
    \mathrm{Re}_D = \frac{U_d D}{\nu} \approx \frac{1.27\times0.01}{2\times10^{-6}} \approx 6350
\]
L'écoulement est turbulent. La perte de charge sur 1~m est de l'ordre de
$\Delta H \approx 0.316\times6350^{-1/4}\times(1/0.01)\times\frac{1}{2}\times800\times1.27^2
\approx 640$~Pa. Pour 20~m de tuyauterie, on obtient~$\sim$12~kPa~: la pompe
à carburant doit fournir cette pression en plus de la pression d'injection.

% ==========================================
\section*{Résumé du chapitre}
\addcontentsline{toc}{section}{Résumé du chapitre}
% ==========================================

\begin{importantbox}{Ce qu'il faut retenir}
\textbf{Théorème $\Pi$} (Vaschy-Buckingham)~: $n$ variables, rang dimensionnel $r$
$\Rightarrow$ $n-r$ groupes adimensionnels, loi de la forme
$\Pi_1 = \mathcal{F}_0(\Pi_2,\ldots,\Pi_{n-r})$.

\vspace{4pt}
\textbf{Nombre de Reynolds}~:
\[
    \mathrm{Re} = \frac{UL}{\nu} = \frac{\text{forces inertielles}}{\text{forces visqueuses}}
    = \frac{\text{temps visqueux}}{\text{temps convectif}}
\]

\vspace{4pt}
\textbf{Nombres sans dimension classiques}~:
\quad St $= fD/U_\infty$,\quad
Fr $= U/\!\sqrt{gL}$,\quad
$C_D = F_D/(\frac{1}{2}\rho U^2 S)$.

\vspace{4pt}
\textbf{NS adimensionnée}~: l'écoulement ne dépend que de Re (pour géométrie fixée).

\vspace{4pt}
\textbf{Poiseuille}~:
$U(r)=U_0[1-(r/R)^2]$,\quad $U_d=U_0/2$,\quad $C_f=16/\mathrm{Re}_D$.

\vspace{4pt}
\textbf{Perte de charge}~:
$\Delta H = \psi(\mathrm{Re}_D)\dfrac{L}{D}\dfrac{1}{2}\rho U_d^2$,\quad
$\psi = 64/\mathrm{Re}_D$ (laminaire).
\end{importantbox}

% ==========================================
\chapter{R\'egimes d'\'ecoulement en fonction du nombre de Reynolds}
% ==========================================

Le chapitre pr\'ec\'edent a \'etabli que le nombre de Reynolds $\mathrm{Re}_L = UL/\nu$
est le param\`etre sans dimension gouvernant les \'ecoulements incompressibles --- il est
seul \`a appara\^itre dans les \'equations de Navier-Stokes adimensionn\'ees. Ce chapitre
explore les deux r\'egimes limites oppos\'es: $\mathrm{Re}_L \ll 1$ (viscosit\'e dominante,
\'ecoulement \emph{rampant} de Stokes) et $\mathrm{Re}_L \gg 1$ (inertie dominante,
apparition d'une \emph{couche limite} pr\`es des parois).

Le Rafale Marine vole \`a $\mathrm{Re} \approx 7\times10^7$ sur son aile (corde $L\approx
3\,\text{m}$, $U_\infty\approx 200\,\text{m/s}$, $\nu_\text{air}\approx
1{,}5\times10^{-5}\,\text{m}^2/\text{s}$), ce qui place sans ambigu\"{i}t\'e l'avion dans
le r\'egime \`a haut nombre de Reynolds. Ses moteurs M88-2 ing\`erent cependant des
particules microm\'etriques (poussi\`eres, bruine saline) dont le nombre de Reynolds
particulaire est inf\'erieur \`a~1 --- preuve que les deux r\'egimes peuvent coexister au
sein d'un m\^eme syst\`eme.

%===========================================================================
\section{\'Ecoulement \`a bas nombre de Reynolds}
%===========================================================================

\subsection{Le terme convectif peut \^etre n\'eglig\'e}

Rappelons les \'equations de Navier-Stokes pour un fluide incompressible de densit\'e
$\rho$ et viscosit\'e cin\'ematique $\nu$~:
\[
  \pd{\bm{U}}{t} + \bm{U}\cdot\nabla\bm{U}
  = -\frac{1}{\rho}\nabla p + \nu\nabla^2\bm{U},
  \qquad \nabla\cdot\bm{U}=0.
\]
En notant $U$ et $L$ les \'echelles caract\'eristiques de vitesse et de longueur, on
estime les ordres de grandeur~:
\[
  \underbrace{\bm{U}\cdot\nabla\bm{U}}_{\sim\, U^2/L}
  \qquad\text{contre}\qquad
  \underbrace{\nu\nabla^2\bm{U}}_{\sim\, \nu U/L^2}.
\]
Le rapport terme convectif / terme visqueux vaut $UL/\nu = \mathrm{Re}_L$. Lorsque
$\mathrm{Re}_L \ll 1$, le terme convectif est n\'egligeable.

\subsection{\'Ecoulement de Stokes ($\mathrm{Re}_L \ll 1$)}

Pour \'etudier ce r\'egime, on choisit comme temps caract\'eristique le \emph{temps de
diffusion visqueuse} $L^2/\nu$ (et non $L/U$ adapt\'e au cas inertiel), car c'est la
viscosit\'e qui dicte la dynamique. En posant~:
\[
  \tilde{U}_i = \frac{U_i}{U_\infty}, \quad
  \tilde{t} = \frac{t}{L^2/\nu}, \quad
  \tilde{p} = \frac{p}{\mu U_\infty/L},
  \quad \mathrm{Re}_L = \frac{\rho L U_\infty}{\mu},
\]
les \'equations de Navier-Stokes adimensionn\'ees deviennent~:
\[
  \pd{\tilde{\bm{U}}}{\tilde{t}} + \mathrm{Re}_L\,\tilde{\bm{U}}\cdot
  \tilde{\nabla}\tilde{\bm{U}}
  = -\tilde{\nabla}\tilde{p} + \tilde{\nabla}^2\tilde{\bm{U}}.
\]
Pour $\mathrm{Re}_L \ll 1$, le terme $\mathrm{Re}_L\,\tilde{\bm{U}}\cdot
\tilde{\nabla}\tilde{\bm{U}} \ll 1$ (sous r\'eserve que la solution approch\'ee reste
d'ordre~1 en variables adimensionn\'ees --- condition \`a v\'erifier \emph{a posteriori},
tr\`es restrictive en pratique). On parle d'\textbf{\'ecoulement rampant} ou
\textbf{\'ecoulement de Stokes}.

\subsection{\'Ecoulement de Stokes stationnaire}

En r\'egime permanent, les \'equations se r\'eduisent \`a~:
\[
  \left\{\begin{aligned}
    \mu\nabla^2\bm{U} &= \nabla p \\[4pt]
    \nabla\cdot\bm{U} &= 0
  \end{aligned}\right.
  \qquad + \text{ conditions aux limites.}
\]
Ce syst\`eme est lin\'eaire --- une simplification spectaculaire par rapport aux
\'equations de Navier-Stokes compl\`etes. Stokes (1851) en trouva la solution analytique
exacte pour une sph\`ere, donnant le champ de vitesse et les iso-contours de pression
de part et d'autre du corps. Le caract\`ere rampant se traduit par une absence de
d\'ecollement: les lignes de courant restent attach\'ees \`a la sph\`ere et l'\'ecoulement
est quasi-sym\'etrique amont/aval (visualisation de Taneda, $\mathrm{Re}_D = 0{,}16$).

\subsection{Force de Stokes --- coefficient de tra\^in\'ee d'une sph\`ere}

La force exerc\'ee par le fluide sur une sph\`ere de rayon $a$ (diam\`etre $D=2a$)
plong\'ee dans un \'ecoulement uniforme $U_\infty$ se d\'ecompose en tra\^in\'ee de
forme (pression) et de frottement (visqueux)~:
\[
  \bm{F}_{\text{fluide}\to\text{sph\`ere}}
  = \int_{\mathcal{S}} \bigl(-p\,\bm{n}\bigr)\,ds
  + \int_{\mathcal{S}} \ten{\tau}\cdot\bm{n}\,ds.
\]
La r\'esolution analytique des \'equations de Stokes donne la force de tra\^in\'ee~:
\[
  F_D = 6\pi\mu a\, U_\infty
  \qquad\text{(force de Stokes, 1851)},
\]
o\`u le facteur $6\pi = 2\pi + 4\pi$: les $2\pi$ viennent de la diff\'erence de pression
amont/aval (tra\^in\'ee de forme), les $4\pi$ du frottement visqueux pari\'etal (tra\^in\'ee
de peau). Le coefficient de tra\^in\'ee en d\'ecoule~:
\[
  C_D \equiv \frac{F_D}{\dfrac{1}{2}\rho U_\infty^2 \cdot \pi a^2}
  = \frac{6\pi\mu a\, U_\infty}{\dfrac{1}{2}\rho U_\infty^2 \cdot \pi a^2}
  = \frac{24}{\rho U_\infty D/\mu}
  = \frac{24}{\mathrm{Re}_D}.
\]

\begin{importantbox}{Force de Stokes --- sph\`ere lisse \`a $\mathrm{Re}_D \ll 1$}
\[
  F_D = 6\pi\mu a\, U_\infty,
  \qquad
  C_D = \dfrac{24}{\mathrm{Re}_D}.
\]
La tra\^in\'ee est domin\'ee par les effets visqueux et cro\^it \emph{lin\'eairement}
avec la vitesse. Sur la courbe $C_D(\mathrm{Re}_D)$ pour une sph\`ere lisse, la loi
$24/\mathrm{Re}_D$ est valide jusqu'\`a $\mathrm{Re}_D \approx 1$; pour
$10^3 \lesssim \mathrm{Re}_D \lesssim 3\times10^5$, $C_D \approx 0{,}4$ (plateau);
puis crise de tra\^in\'ee \`a $\mathrm{Re}_D^c \approx 3\times10^5$.
\end{importantbox}

\subsection{Applications de la force de Stokes}

La condition $\mathrm{Re} = U_\infty d_p/\nu \leq 1$ est satisfaite dans de nombreux
contextes~:

\begin{itemize}
  \item \textbf{Ensemencement en m\'ecanique des fluides exp\'erimentale (PIV, LDV)~:}
    gouttes d'huile d'olive de diam\`etre $d_p = 1{,}5\,\mu$m, $\rho = 970\,\text{kg/m}^3$,
    $\nu = 1{,}5\times10^{-5}\,\text{m}^2/\text{s}$. La condition $\mathrm{Re} \leq 1$
    impose $U_\infty < 9{,}7\,\text{m/s}$ --- largement v\'erifi\'e dans les souffleries
    o\`u ces particules servent de traceurs fid\`eles.

  \item \textbf{Exp\'erience de Millikan (1913)~:} des gouttelettes d'huile \'electris\'ees
    tombent sous gravit\'e, frein\'ees par la force de Stokes. L'application d'un champ
    \'electrique compensant exactement les deux forces permit \`a Millikan de mesurer la
    charge \'el\'ementaire $e = 1{,}6\times10^{-19}\,\text{C}$.

  \item \textbf{Moteur M88-2 du Rafale~:} l'air est aspir\'e \`a ${\approx}60\,\text{m/s}$
    \`a l'entr\'ee. Des particules de poussi\`ere de diam\`etre $d_p\sim 10\,\mu$m voient
    une vitesse relative $\Delta U\sim 1\,\text{m/s}$ (elles ne suivent pas imm\'ediatement
    l'acc\'el\'eration du fluide). Alors $\mathrm{Re}_p = \Delta U\,d_p/\nu \approx 0{,}7
    \ll 1$~: la force de Stokes s'applique. C'est ce principe qui fonde les s\'eparateurs
    centrifuges plac\'es en entr\'ee de turbor\'eacteurs militaires pour limiter
    l'\'erosion des aubes de compresseur.
\end{itemize}

%===========================================================================
\section{\'Ecoulement \`a haut nombre de Reynolds}
%===========================================================================

\subsection{Retour des \'equations d'Euler --- et leurs limites}

\`A l'oppos\'e, pour $\mathrm{Re} \gg 1$, le terme visqueux $\nu\nabla^2\bm{U}
\sim \nu U/L^2$ devient n\'egligeable devant le terme convectif $\sim U^2/L$. Le
mod\`ele non visqueux du chapitre~1 peut \^etre r\'eappliqu\'e~:
\[
  \left\{\begin{aligned}
    \pd{\bm{U}}{t} + \bm{U}\cdot\nabla\bm{U} &= -\dfrac{1}{\rho}\nabla p \\[6pt]
    \nabla\cdot\bm{U} &= 0
  \end{aligned}\right.
  \qquad\text{(\'equations d'Euler).}
\]

Mais cette simplification a une limite cruciale~: les \'equations d'Euler ne peuvent
pas satisfaire la condition d'adh\'erence $\bm{U}=\bm{0}$ \`a une paroi solide. La
vitesse doit passer de $U_\infty$ (\`a l'infini) \`a~0 (paroi) sur une distance tr\`es
courte, g\'en\'erant des gradients de vitesse $\partial U/\partial x_2 \gg U/L$ et donc
des contraintes visqueuses localement importantes m\^eme si $\mu$ est petit~: c'est la
\textbf{couche limite}.

\subsection{Concept de couche limite (Prandtl, 1904)}

En 1904, Ludwig Prandtl proposa en dix pages seulement une id\'ee r\'evolutionnaire~:
l'\'ecoulement se structure en deux r\'egions distinctes.

\begin{center}
\begin{tikzpicture}[scale=1.05]
  % Corps profilé stylisé
  \draw[thick,fill=gray!25]
    (0,0) .. controls (1,0.45) and (3,0.55) .. (5,0.1)
          .. controls (3,-0.05) and (1,-0.05) .. (0,0);
  % Couche limite (bande rouge au-dessus du profil)
  \draw[red!75,thick,dashed]
    (0.1,0.05) .. controls (1,0.78) and (3,0.88) .. (5,0.38);
  % Flèche delta
  \draw[<->,red!70,thick] (2.5,0.48) -- (2.5,0.78);
  \node[red!70,right] at (2.55,0.63) {$\delta \ll L$};
  % Lignes d'écoulement externe
  \foreach \y in {1.0,1.4,1.8} {
    \draw[->,blue!65,thick] (-1.4,\y) -- (-0.1,\y);
  }
  \node[blue!75] at (-0.9,1.4) {$U_\infty$};
  % Labels
  \draw[<->] (0,-0.4) -- (5,-0.4);
  \node[below] at (2.5,-0.4) {$L$};
  \node[blue!75,above] at (3.5,1.85)
    {\footnotesize \'ecoulement externe (Euler)};
  \node[red!75,below] at (3.5,-0.1)
    {\footnotesize couche limite~: effets visqueux dominants};
\end{tikzpicture}
\end{center}

\begin{itemize}
  \item \textbf{L'\'ecoulement externe} (loin de la paroi)~: non visqueux, r\'egi par
    les \'equations d'Euler. L'\'epaisseur de couche limite $\delta \to 0$ quand
    $\mathrm{Re}\to\infty$.
  \item \textbf{La couche limite} (pr\`es de la paroi)~: mince couche d'\'epaisseur
    $\delta \ll L$ dans laquelle les effets visqueux sont localement importants malgr\'e
    un nombre de Reynolds global \'elev\'e.
\end{itemize}

Deux remarques~:
\begin{enumerate}
  \item Lorsque $\mathrm{Re}$ d\'ecro\^it, $\delta$ grossit~; pour $\mathrm{Re}
    \lesssim \mathcal{O}(1)$, il n'y a plus de structuration particuli\`ere et c'est
    l'\'ecoulement de Stokes global.
  \item \`A $\mathrm{Re}$ suffisamment \'elev\'e, la couche limite laminaire peut
    transitionner vers une couche limite \textbf{turbulente} --- aux cons\'equences
    majeures sur la tra\^in\'ee.
\end{enumerate}

\subsection{Corps profil\'e --- couche limite attach\'ee}

Pour un corps profil\'e (comme l'aile du Rafale), la g\'eom\'etrie est con\c{c}ue pour
que la couche limite reste attach\'ee sur toute (ou presque toute) la surface~:

\begin{itemize}
  \item Au bord d'attaque, l'\'ecoulement s'arr\^ete au \textbf{point d'arr\^et}, puis
    acc\'el\`ere le long de l'extrados et de l'intrados.
  \item La couche limite se d\'eveloppe le long de la corde et se s\'epare seulement
    pr\`es du bord de fuite, g\'en\'erant un \textbf{sillage mince}.
  \item La tra\^in\'ee est alors domin\'ee par la tra\^in\'ee de frottement (tra\^in\'ee
    de peau), faible car $\delta \ll L$.
\end{itemize}

Si l'angle d'attaque $\alpha$ augmente, la couche limite se s\'epare pr\'ematur\'ement~:
le profil non d\'ecorch\'e (corps profil\'e, sillage mince) devient d\'ecorn\'e (corps
non profil\'e, large sillage turbulent). Pour le profil ONERA~64A015, le d\'ecrochage
est quasi-total d\`es $\alpha \approx 17{,}5^\circ$ (visualisation Werl\'e, 1974).

\subsection{Corps non profil\'e --- crise de tra\^in\'ee}

Pour un corps non profil\'e (sph\`ere, cylindre, arri\`ere tronqu\'e d'un v\'ehicule),
les dimensions longitudinale et transversale sont comparables. La couche limite se
s\'epare bien avant l'arri\`ere, g\'en\'erant un \textbf{large sillage turbulent}
source d'une forte tra\^in\'ee de forme.

\paragraph{Crise de tra\^in\'ee.}
Pour un nombre de Reynolds critique $\mathrm{Re}^c$, $C_D$ chute brusquement. Pour
la sph\`ere lisse, $\mathrm{Re}_D^c \approx 3\times10^5$. Le m\'ecanisme est le
suivant~:

\begin{enumerate}
  \item En dessous de $\mathrm{Re}^c$~: couche limite laminaire, s\'eparation pr\'ecoce
    (vers $\theta\approx 80^\circ$ depuis le bord d'attaque), \textbf{large sillage}
    et forte tra\^in\'ee de forme.
  \item Au-dessus de $\mathrm{Re}^c$~: la couche limite transionne en r\'egime
    \textbf{turbulent} avant le point de s\'eparation. L'\'energie cin\'etique turbulente
    permet au fluide de r\'esister plus longtemps au gradient de pression adverse~;
    la s\'eparation est repouss\'ee vers l'aval (vers $\theta\approx 120^\circ$). Le
    sillage est \textbf{beaucoup plus mince}~: la tra\^in\'ee de forme chute. La l\'eg\`ere
    augmentation de tra\^in\'ee de frottement (CL turbulente) est amplement compens\'ee.
\end{enumerate}

\begin{importantbox}{Crise de tra\^in\'ee --- bilan physique}
Transition laminaire $\to$ turbulent de la couche limite $\Rightarrow$ point de s\'eparation
repouss\'e en aval $\Rightarrow$ sillage plus mince $\Rightarrow$ chute brutale de $C_D$.
Pour la sph\`ere~: $C_D \approx 0{,}4$ (subcritique) $\to$ $C_D \approx 0{,}1$
(supercritique).
\end{importantbox}

\paragraph{La balle de golf.}
Les alv\'eoles (\emph{dimples}) induisent une transition turbulente pr\'ematur\'ee,
abaissant $\mathrm{Re}^c$ \emph{en dessous} du Reynolds de jeu ($D=4{,}3\,\text{cm}$,
$U\approx 67\,\text{m/s}$, $\mathrm{Re}\approx 1{,}9\times10^5$). Sans alv\'eoles,
la balle serait subcritique~: $C_D$ quatre fois plus grand, distance de vol divis\'ee
par quatre.

\paragraph{Le cylindre lisse.}
Pour le cylindre, $\mathrm{Re}_D^c \approx 2\times10^5$. Son sillage pr\'esente de
surcro\^it le ph\'enom\`ene de l'\textbf{all\'ee de tourbillons de K\'arm\'an}~: un
l\^acher p\'eriodique et alternant de tourbillons, caract\'eris\'e par le nombre de
Strouhal $\mathrm{St} = fD/U \approx 0{,}2$ (d\'eriv\'e au chapitre~3 par analyse
dimensionnelle). Von K\'arm\'an rapporte lui-m\^eme que l'id\'ee lui fut inspir\'ee par
un tableau m\'edi\'eval de Saint Christophe traversant un fleuve (\'eglise
Saint-Dominique, Bologne), dont les pieds cr\'eent exactement une telle all\'ee dans
l'eau. Ces all\'ees de K\'arm\'an sont visibles depuis l'espace~: des nuages organis\'es
en double rang\'ee altern\'ee se forment en aval de l'\^ile Alexander Selkirk (Pacifique
Sud) ou des \^iles Kouriles (entre le Kamtchatka et le Japon).

\paragraph{G\'en\'erateurs de vortex.}
Sur l'aile du Rafale --- et plus g\'en\'eralement sur les avions \`a hautes performances
--- des \textbf{g\'en\'erateurs de vortex} (petites ailettes plac\'ees en amont de la
zone \`a risque) induisent d\'elib\'er\'ement une transition turbulente locale. La couche
limite turbulente r\'esiste mieux au gradient de pression adverse, retardant le d\'ecrochage
et \'elargissant le domaine de vol utilisable. On les distingue nettement sur le Boeing~777
le long de l'intrados de la voilure, et ils \'equipent \'egalement certaines versions
du Beechcraft Baron.

%===========================================================================
\section{Couche limite laminaire}
%===========================================================================

\subsection{G\'eom\'etrie et hypoth\`ese d'\'ecoulement mince}

Consid\'erons une plaque plane d'\'envergure unitaire soumise \`a un \'ecoulement uniforme
de vitesse $U_{e1}$ parall\`ele \`a la plaque. On note $x_1$ la direction de l'\'ecoulement,
$x_2$ la normale, $\delta(x_1)$ l'\'epaisseur locale de la couche limite et $L$ la
distance depuis le bord d'attaque.

\begin{center}
\begin{tikzpicture}[scale=1.0]
  % Plaque
  \fill[gray!20] (-0.1,-0.18) rectangle (7,0);
  \draw[thick] (0,0) -- (7,0);
  % Entrée
  \foreach \y in {0.4,0.9,1.4,1.9} {
    \draw[->,blue!65,thick] (-1.6,\y) -- (-0.1,\y);
  }
  \node[blue!75] at (-1.0,1.3) {$U_{e1}$};
  % Couche limite
  \draw[red!80,thick] (0,0) .. controls (2.0,0.28) and (4.5,0.52) .. (7,0.72);
  \node[red!75,above] at (5.5,0.78) {\small couche limite};
  % delta
  \draw[<->,red!65] (3.5,0) -- (3.5,0.46);
  \node[red!65,right] at (3.55,0.23) {$\delta$};
  % L arrow
  \draw[<->] (0,-0.38) -- (3.5,-0.38);
  \node[below] at (1.75,-0.38) {$L$};
  % Profil de vitesse (à x = 3.5)
  \draw[blue!50,thick,->] (3.5,0) -- (3.5,1.8) node[right]{\small $x_2$};
  \draw[blue!60,thick]
    (3.5,0) .. controls (3.8,0.2) and (3.92,0.4) .. (3.93,0.46) -- (3.93,1.6);
  \node[blue!55,right] at (4.0,1.3) {\small $U_1(x_2)$};
  \node[right] at (7.1,0.36) {\small $U_{e1}$};
  % Axes x1
  \draw[->] (0,-0.6) -- (7.5,-0.6) node[right]{$x_1$};
\end{tikzpicture}
\end{center}

L'\textbf{hypoth\`ese d'\'ecoulement mince} est $\delta \ll L$. Typiquement~: $\delta
\approx 1\,\text{cm}$ pour $L\approx 1\,\text{m}$ sur une aile d'avion. Pour le Rafale
($L\approx 3\,\text{m}$, $\mathrm{Re}_L \approx 4\times10^7$)~:
\[
  \frac{\delta}{L} \sim \mathrm{Re}_L^{-1/2} \approx \frac{1}{\sqrt{4\times10^7}}
  \approx 1{,}6\times10^{-4}
  \quad\Rightarrow\quad
  \delta \approx 0{,}5\,\text{mm.}
\]
Moins d'un millim\`etre sur trois m\`etres de corde --- une minceur qui l\'egitime
toutes les approximations qui suivent.

\subsection{Analyse d'ordre de grandeur}

Les \'equations de Navier-Stokes pour un \'ecoulement 2-D stationnaire incompressible
($U_3=0$) sont~:
\[
  \left\{\begin{aligned}
    U_1\pd{U_1}{x_1} + U_2\pd{U_1}{x_2} &=
      -\dfrac{1}{\rho}\pd{p}{x_1}
      + \nu\!\left(\pd{^2U_1}{x_1^2}+\pd{^2U_1}{x_2^2}\right)\\[6pt]
    U_1\pd{U_2}{x_1} + U_2\pd{U_2}{x_2} &=
      -\dfrac{1}{\rho}\pd{p}{x_2}
      + \nu\!\left(\pd{^2U_2}{x_1^2}+\pd{^2U_2}{x_2^2}\right)\\[6pt]
    \pd{U_1}{x_1} + \pd{U_2}{x_2} &= 0
  \end{aligned}\right.
\]

\paragraph{Continuit\'e.}
Les deux termes de l'\'equation de continuit\'e sont n\'ecessairement du m\^eme ordre~:
$\mathcal{U}_1/L \sim \mathcal{U}_2/\delta$, d'o\`u~:
\[
  \mathcal{U}_2 \sim \frac{\delta}{L}\,\mathcal{U}_1 \ll \mathcal{U}_1.
\]
La composante transverse est bien plus petite que la composante longitudinale --- coh\'erent
avec un \'ecoulement quasi-parall\`ele \`a la plaque.

\paragraph{\'Equation selon $x_1$.}
On estime chaque terme~:
\[
  \underbrace{U_1\pd{U_1}{x_1}}_{\sim\,\mathcal{U}_1^2/L}
  +\underbrace{U_2\pd{U_1}{x_2}}_{\sim\,\mathcal{U}_1^2/L}
  = -\dfrac{1}{\rho}\pd{p}{x_1}
  + \nu\!\underbrace{\pd{^2U_1}{x_1^2}}_{\sim\,\nu\mathcal{U}_1/L^2}
  + \nu\!\underbrace{\pd{^2U_1}{x_2^2}}_{\sim\,\nu\mathcal{U}_1/\delta^2}.
\]
Le premier terme de viscosit\'e ($\sim\nu\mathcal{U}_1/L^2$) est n\'egligeable devant le
second ($\sim\nu\mathcal{U}_1/\delta^2$) car $\delta \ll L$. Les termes convectifs doivent
\'equilibrer le terme visqueux dominant~; ni plus grand (sinon l'\'ecoulement serait
purement non visqueux, contradiction avec l'adh\'erence), ni plus petit (sinon on aurait
un \'ecoulement de Stokes, valide seulement pour $\mathrm{Re}\leq 1$). Cet \'equilibre
impose~:
\[
  \frac{\mathcal{U}_1^2}{L} \sim \nu\frac{\mathcal{U}_1}{\delta^2}
  \quad\Rightarrow\quad
  \left(\frac{\delta}{L}\right)^2 \sim \frac{\nu}{\mathcal{U}_1 L}
  = \frac{1}{\mathrm{Re}_L}
  \quad\Rightarrow\quad
  \boxed{\dfrac{\delta}{L} \sim \mathrm{Re}_L^{-1/2}}.
\]

\begin{importantbox}{\'Epaisseur de couche limite --- loi d'\'echelle}
\[
  \dfrac{\delta}{L} \sim \mathrm{Re}_L^{-1/2}, \qquad
  \delta \propto \sqrt{x_1}.
\]
Pour que $\delta \ll L$, il faut n\'ecessairement $\mathrm{Re}_L \gg 1$~: la th\'eorie
de couche limite n'est valide qu'\`a grand nombre de Reynolds. Par ailleurs $\mathrm{Re}_\delta
= \mathrm{Re}_L^{1/2}$, interm\'ediaire entre~1 et $\mathrm{Re}_L$.
\end{importantbox}

\paragraph{\'Equation selon $x_2$.}
La m\^eme analyse montre que tous les termes de la composante transversale sont plus petits
d'un facteur $(\delta/L)^2 \sim 1/\mathrm{Re}_L$ par rapport \`a l'\'equation selon $x_1$.
En particulier, la variation de pression \`a travers la couche limite est d'ordre
$\rho\mathcal{U}_1^2/\mathrm{Re}_L$ --- n\'egligeable. On en conclut~:
\[
  \pd{p}{x_2} \approx 0
  \quad\Rightarrow\quad
  p(x_1,x_2) \approx p_e(x_1).
\]
La pression est \textbf{constante \`a travers la couche limite}~: elle est impos\'ee par
l'\'ecoulement externe eul\'erien. Bernoulli appliqu\'e \`a l'ext\'erieur donne~:
\[
  \pd{p}{x_1} = \frac{dp_e}{dx_1} = -\rho U_{e1}\frac{dU_{e1}}{dx_1}.
\]

\subsection{\'Equations de Prandtl (1904)}

En retenant les simplifications \'etablies ci-dessus, Navier-Stokes se r\'eduit aux
\textbf{\'equations de Prandtl}~:

\begin{importantbox}{\'Equations de Prandtl --- couche limite laminaire 2-D}
\[
  \left\{\begin{aligned}
    U_1\pd{U_1}{x_1} + U_2\pd{U_1}{x_2}
      &= U_{e1}\dfrac{dU_{e1}}{dx_1} + \nu\pd{^2U_1}{x_2^2}\\[8pt]
    \pd{U_1}{x_1} + \pd{U_2}{x_2} &= 0
  \end{aligned}\right.
\]
avec $U_1 = U_2 = 0$ en $x_2 = 0$~~(adh\'erence)~~et~~$U_1\to U_{e1}(x_1)$ quand
$x_2\to\infty$.
\end{importantbox}

Ces \'equations sont la \emph{forme parabolois\'ee} des \'equations de Navier-Stokes~:
le terme de diffusion selon $x_1$ a disparu, et l'\'equation peut \^etre int\'egr\'ee
pas \`a pas depuis le bord d'attaque vers l'aval, pourvu que l'\'ecoulement externe
$(U_{e1}, p_e)$ soit connu.

\subsection{R\^ole du gradient de pression}

La pression impos\'ee par l'\'ecoulement externe joue un r\^ole d\'eterminant~:
\begin{itemize}
  \item \textbf{Gradient favorable} ($dp_e/dx_1 < 0$, \'ecoulement acc\'el\'er\'e)~:
    la pression pousse le fluide vers l'aval, favorisant l'attachement de la couche limite.
    Aucun d\'ecollement ne se produit.
  \item \textbf{Gradient adverse} ($dp_e/dx_1 > 0$, \'ecoulement d\'ec\'el\'er\'e)~:
    la pression s'oppose \`a l'\'ecoulement. Le fluide proche de la paroi, d\'ej\`a
    ralenti par le frottement, peut \^etre stopp\'e puis invers\'e~: c'est la
    \textbf{s\'eparation} (d\'ecollement) de la couche limite, accompagn\'ee d'une zone
    de recirculation et d'un large sillage.
\end{itemize}

Sur l'aile du Rafale, l'extrados pr\'esente une acc\'el\'eration forte pr\`es du bord
d'attaque (gradient favorable) puis une d\'ec\'el\'eration vers le bord de fuite (gradient
adverse). Le d\'ecrochage ($\alpha\gtrsim 25^\circ$ en configuration basse vitesse)
survient quand ce gradient adverse devient trop fort pour la couche limite.

\subsection{Solution de Blasius (1908)}

Pour la plaque plane \`a gradient de pression nul ($U_{e1} = \text{cst}$,
$dp_e/dx_1=0$), les \'equations de Prandtl admettent une \textbf{solution autosimilaire}~:
Blasius (premier doctorant de Prandtl) montra en 1908 que le profil $U_1/U_{e1}$ est une
fonction universelle de la variable $\eta = x_2\sqrt{U_{e1}/(\nu x_1)}$ seulement. La
r\'esolution num\'erique de l'\'equation diff\'erentielle ordinaire r\'esultante donne~:
\[
  \frac{\delta}{x_1} \simeq \frac{4{,}92}{\sqrt{\mathrm{Re}_{x_1}}},
  \qquad
  C_f = \frac{\tau_w}{\dfrac{1}{2}\rho U_{e1}^2}
  \simeq \frac{0{,}664}{\sqrt{\mathrm{Re}_{x_1}}},
\]
o\`u $\mathrm{Re}_{x_1} = U_{e1}x_1/\nu$ et $\tau_w = \mu\,\partial U_1/\partial
x_2\big|_{x_2=0}$ est la contrainte pari\'etale.

\begin{defbox}{Solution de Blasius --- plaque plane sans gradient de pression}
\[
  \dfrac{\delta}{x_1} \simeq \dfrac{4{,}92}{\sqrt{\mathrm{Re}_{x_1}}},
  \qquad
  C_f \simeq \dfrac{0{,}664}{\sqrt{\mathrm{Re}_{x_1}}}.
\]
La couche limite laminaire s'\'epaissit comme $\sqrt{x_1}$ et le frottement pari\'etal
d\'ecro\^it comme $\mathrm{Re}_{x_1}^{-1/2}$. Le profil de vitesse universel de Blasius
a \'et\'e v\'erifi\'e exp\'erimentalement avec une pr\'ecision remarquable d\`es 1942
(Dhawan).
\end{defbox}

\paragraph{Application au Rafale.}
Pour $L = 3\,\text{m}$, $U_\infty = 200\,\text{m/s}$, $\nu = 1{,}5\times10^{-5}\,
\text{m}^2/\text{s}$, $\mathrm{Re}_L = 4\times10^7$~:
\[
  \delta(L) = \frac{4{,}92 \times 3}{\sqrt{4\times10^7}} \approx 2{,}3\,\text{mm},
  \qquad
  C_f(L) \approx \frac{0{,}664}{\sqrt{4\times10^7}} \approx 1{,}0\times10^{-4}.
\]
La surface mouill\'ee du Rafale \'etant $\approx 70\,\text{m}^2$, la tra\^in\'ee de
peau totale est loin d'\^etre n\'egligeable m\^eme si $C_f$ est tr\`es faible.

% ==========================================
\section*{Résumé du chapitre}
\addcontentsline{toc}{section}{Résumé du chapitre}
% ==========================================

\begin{importantbox}{\`A retenir --- Chapitre 4}
\textbf{Bas $\mathrm{Re}$~:} \'ecoulement de Stokes, $F_D = 6\pi\mu a\,U_\infty$,
$C_D = 24/\mathrm{Re}_D$. Applicable aux petites particules (PIV, Millikan, s\'eparateurs
de moteurs).

\textbf{Haut $\mathrm{Re}$~:} \'ecoulement d'Euler \`a l'ext\'erieur + couche limite
d'\'epaisseur $\delta/L \sim \mathrm{Re}_L^{-1/2}$. Corps profil\'e~: faible tra\^in\'ee
de peau. Corps non profil\'e~: forte tra\^in\'ee de forme --- crise de tra\^in\'ee \`a
$\mathrm{Re}^c$ par transition turbulente de la CL (balle de golf, g\'en\'erateurs de
vortex).

\textbf{\'Equations de Prandtl~:} pression impos\'ee par l'ext\'erieur~;
gradient adverse $\Rightarrow$ s\'eparation. Solution de Blasius~: $\delta \propto
\sqrt{x_1}$, $C_f \propto \mathrm{Re}_{x_1}^{-1/2}$.
\end{importantbox}

% ==========================================
\chapter{Vorticit\'e}
% ==========================================

Lorsqu'un Rafale Marine aponte sur le \emph{Charles de Gaulle}, les pilotes suivants
doivent attendre plusieurs dizaines de secondes avant de s'engager dans la m\^eme
trajectoire. La raison~: l'air derri\`ere l'appareil s'enroule en deux puissants
tourbillons de saumon, visibles par condensation de la vapeur d'eau, capables de
d\'estabiliser un a\'eronef l\'eger qui les traverserait.

Ces tourbillons ne sont pas un accident~: ils sont la cons\'equence directe de la
\emph{portance}. Plus pr\'ecis\'ement, ils sont li\'es \`a une grandeur fondamentale
que nous n'avons pas encore isol\'ee en tant que telle~: la \textbf{vorticit\'e},
c'est-\`a-dire la rotation locale du fluide. Ce chapitre lui est enti\`erement consacr\'e.

On va d'abord d\'efinir proprement la vorticit\'e et la circulation, puis montrer
comment elles \'evoluent (\'equation de transport), d'o\`u elles viennent (les parois),
et enfin comment elles permettent de calculer la portance d'un profil d'aile
(th\'eor\`eme de Kutta-Joukowski). Le fil conducteur est clair~: comprendre pourquoi
le Rafale vole, c'est comprendre la vorticit\'e qu'il g\'en\`ere.

\begin{defbox}{Fil conducteur}
La vorticit\'e est la grandeur qui relie la couche limite (chapitre~4) \`a la portance
(chapitre~1, Bernoulli). Elle na\^it aux parois par cisaillement visqueux, est transport\'ee
par l'\'ecoulement, et sa circulation autour du profil d\'etermine directement la force de
sustentation. Les tourbillons de saumon du Rafale en sont la manifestation tridimensionnelle.
\end{defbox}

% ------------------------------------------
\section{D\'efinition et interpr\'etation physique}
% ------------------------------------------

Au chapitre~2, nous avons d\'ecompos\'e le gradient de vitesse $\nabla\bm{U}$ en une partie
sym\'etrique $\ten{D}$ (taux de d\'eformation) et une partie antisym\'etrique $\ten{\Omega}$
(taux de rotation). La partie antisym\'etrique d\'ecrit une rotation locale de corps rigide.
Il est naturel de vouloir condenser cette information en un seul vecteur~: c'est exactement
ce que fait la vorticit\'e.

Formellement, on d\'efinit~:

\begin{defbox}{Vorticit\'e}
\[
  \boldsymbol{\omega} = \nabla \times \bm{U}
\]
Le vecteur $\boldsymbol{\omega}$ est \emph{deux fois} la vitesse angulaire instantan\'ee
d'une particule fluide~: $\boldsymbol{\omega} = 2\boldsymbol{\Omega}_\text{rotation}$.
On rappelle que le tenseur des taux de rotation $\ten{\Omega}$
(partie antisym\'etrique de $\nabla\bm{U}$, vue au chapitre~2) v\'erifie
$\Omega_{ij} = -\tfrac{1}{2}\varepsilon_{ijk}\omega_k$.
\end{defbox}

En deux dimensions, o\`u $\bm{U} = U_1(x_1,x_2)\,\bm{e}_1 + U_2(x_1,x_2)\,\bm{e}_2$,
la vorticit\'e se r\'eduit \`a une seule composante scalaire selon $\bm{e}_3$~:
\[
  \omega_3 = \pd{U_2}{x_1} - \pd{U_1}{x_2}.
\]

Pour fixer les id\'ees, consid\'erons l'\'ecoulement de Couette plan du chapitre~2~:
$\bm{U} = (U_e x_2/h,\,0,\,0)$. On calcule imm\'ediatement $\omega_3 = -U_e/h$~: la
vorticit\'e est uniforme et n\'egative (rotation horaire), proportionnelle au gradient
de vitesse. Plus le cisaillement est fort, plus les particules fluides tournent vite.

Un \'ecoulement est dit \textbf{irrotationnel} si $\boldsymbol{\omega} = \bm{0}$
partout dans le domaine. On verra que c'est une cl\'e pour mod\'eliser les \'ecoulements
\`a grand Reynolds autour de profils portants~: loin des parois, l'\'ecoulement est
irrotationnel, et cela d\'everrouille une version puissante de l'\'equation de Bernoulli.

% ------------------------------------------
\section{Tourbillon circulaire et circulation}
% ------------------------------------------

Pour ancrer la notion de vorticit\'e sur un exemple concret, consid\'erons le tourbillon
circulaire --- mod\`ele id\'ealis\'e des tourbillons de saumon du Rafale, mais aussi
des tornades, des tourbillons de vidange, et des cyclones.

\subsection{Tourbillon de Rankine}

Consid\'erons un \'ecoulement purement azimutal \`a sym\'etrie axiale~:
$\bm{U} = u(r)\,\bm{e}_\theta$ en coordonn\'ees cylindriques. La vorticit\'e
axiale s'\'ecrit, en remarquant que $\nabla\times\bm{U}$ en coordonn\'ees cylindriques
donne
\[
  \omega_z = \frac{1}{r}\pd{(r\,u)}{r}.
\]
C'est la d\'erivation de $ru$ qui appara\^it, non pas de $u$ seul~: cela traduit le fait que
m\^eme une rotation solide (o\`u la vitesse cro\^it lin\'eairement) peut \^etre rotationnelle.

Le \textbf{tourbillon de Rankine} est le mod\`ele canonique d'un vortex de rayon caract\'eristique
$r_0$~:
\[
  u(r) = \begin{cases} \Omega_0\, r & r \le r_0 \quad \text{(rotation solide, }
  \omega_z = 2\Omega_0\text{)} \\ \dfrac{\Omega_0 r_0^2}{r} & r > r_0 \quad
  \text{(irrotationnel, } \omega_z = 0\text{)} \end{cases}
\]
La vitesse est continue en $r = r_0$, et la vorticit\'e est concentr\'ee \`a l'int\'erieur
du c\oe ur de rayon $r_0$, nulle \`a l'ext\'erieur.

\begin{center}
\begin{tikzpicture}[scale=1.1]
  % Axes
  \draw[->] (-0.2,0) -- (4.2,0) node[right] {$r$};
  \draw[->] (0,-0.2) -- (0,2.5) node[above] {$u(r)$};
  % r0 marker
  \draw[dashed, gray] (1.5,0) -- (1.5,1.5);
  \node[below] at (1.5,0) {$r_0$};
  % solid rotation part
  \draw[blue, very thick] (0,0) -- (1.5,1.5);
  % irrotational part
  \draw[red, very thick, domain=1.5:4.0, samples=60]
    plot (\x, {1.5*1.5/\x});
  % labels
  \node[blue] at (0.7,1.1) {\small rotation solide};
  \node[red] at (3.0,0.8) {\small irrotationnel};
  \node[left] at (0,1.5) {$\Omega_0 r_0$};
  % omega annotation
  \draw[->] (0.5, 0.05) arc (0:340:0.25);
  \node at (0.5, 0.5) {\small $\omega_z = 2\Omega_0$};
  \node at (2.8, 0.25) {\small $\omega_z = 0$};
\end{tikzpicture}
\end{center}

\subsection{Circulation et th\'eor\`eme de Stokes}

La vorticit\'e d\'ecrit une propri\'et\'e \emph{locale} du champ de vitesse. Mais en
pratique, on a souvent besoin d'une mesure \emph{globale} de la rotation d'un \'ecoulement
autour d'un contour ferm\'e --- par exemple autour d'un profil d'aile. C'est le r\^ole
de la \textbf{circulation}~:

\begin{defbox}{Circulation}
\[
  \Gamma_{\mathcal{C}} = \oint_{\mathcal{C}} \bm{U} \cdot d\bm{x}
\]
Le th\'eor\`eme de Stokes relie imm\'ediatement circulation et vorticit\'e~:
\[
  \Gamma_{\mathcal{C}} = \iint_{\mathcal{S}} \boldsymbol{\omega} \cdot \bm{n}\, dS
\]
o\`u $\mathcal{S}$ est une surface quelconque s'appuyant sur $\mathcal{C}$.
\end{defbox}

Pour le tourbillon circulaire, en int\'egrant sur un cercle de rayon $r$~:
\[
  \Gamma = 2\pi r\, u(r).
\]
Dans la zone irrotationnelle ($r > r_0$), $\Gamma = 2\pi\Omega_0 r_0^2 = \text{cst}$,
ind\'ependamment du rayon du contour. C'est une propri\'et\'e remarquable~: la circulation
autour de tout contour entourant le c\oe ur vaut la m\^eme constante $\Gamma$, m\^eme si
l'\'ecoulement est partout irrotationnel hors du c\oe ur.

\subsection{Tourbillon ponctuel}

En faisant tendre $r_0 \to 0$ \`a $\Gamma$ fix\'e, tout le flux de vorticit\'e se concentre en
un point~: on obtient un \textbf{tourbillon ponctuel} dont le champ de vitesse est
\[
  u(r) = \frac{\Gamma}{2\pi r}, \qquad
  \omega_3 = \Gamma\,\delta(x_1)\delta(x_2).
\]
Ce mod\`ele singulier est l'analogue en 2D du tourbillon rectiligne infini, et il servira de
brique \'el\'ementaire pour construire l'\'ecoulement autour d'un profil portant.

% ------------------------------------------
\section{Version forte de l'\'equation de Bernoulli}
% ------------------------------------------

Au chapitre~1, l'\'equation de Bernoulli a \'et\'e \'etablie le long d'une ligne de courant,
en supposant l'\'ecoulement parfait et stationnaire. La constante $\mathcal{H}$ pouvait
diff\'erer d'une ligne de courant \`a l'autre. On va montrer ici que si l'\'ecoulement est
de surcro\^it \emph{irrotationnel}, cette constante est la m\^eme dans tout le domaine~:
un r\'esultat consid\'erablement plus puissant.

\subsection{Reformulation des \'equations d'Euler}

Le point de d\'epart est une identit\'e vectorielle qui fait appara\^itre explicitement
la vorticit\'e dans le terme convectif. On peut montrer (en d\'eveloppant composante par
composante) que~:
\[
  (\bm{U}\cdot\nabla)\bm{U} = \boldsymbol{\omega} \times \bm{U}
  + \nabla\!\left(\tfrac{1}{2}U^2\right).
\]
En substituant dans les \'equations d'Euler (fluide parfait incompressible, avec
$\bm{g} = -g\bm{e}_3 = \nabla(-gx_3)$), on regroupe les gradients~:
\[
  \pd{\bm{U}}{t} + \boldsymbol{\omega}\times\bm{U}
  = -\nabla\!\left(\frac{p}{\rho} + \frac{U^2}{2} + gx_3\right).
\]
On reconna\^it dans le gradient la charge totale $H/\rho$ d\'ej\`a rencontr\'ee au
chapitre~1. En d\'efinissant $H = p + \tfrac{1}{2}\rho U^2 + \rho g x_3$, l'\'equation
prend la \textbf{forme de Crocco}~:
\[
  \pd{\bm{U}}{t} + \boldsymbol{\omega}\times\bm{U}
  = -\frac{1}{\rho}\nabla H.
\]

\subsection{Cons\'equence pour un \'ecoulement irrotationnel}

Si l'\'ecoulement est \textbf{stationnaire} ($\partial_t\bm{U} = \bm{0}$) et
\textbf{irrotationnel} ($\boldsymbol{\omega} = \bm{0}$), le membre de gauche s'annule
et l'on obtient imm\'ediatement~:

\begin{importantbox}{Bernoulli fort --- \'ecoulement irrotationnel}
Si l'\'ecoulement est \textbf{irrotationnel} ($\boldsymbol{\omega}=\bm{0}$),
\textbf{stationnaire} ($\partial_t=0$) et \textbf{non visqueux}~:
\[
  \nabla H = \bm{0} \quad \Longrightarrow \quad
  H = p + \tfrac{1}{2}\rho U^2 + \rho g x_3 = \mathrm{cst} \quad \text{partout dans le domaine.}
\]
Contrairement au Bernoulli classique (valable le long d'une ligne de courant seulement),
cette version forte est uniforme dans tout le fluide.
\end{importantbox}

La physique est la suivante~: dans le Bernoulli ``faible'' du chapitre~1, la constante
pouvait diff\'erer entre deux lignes de courant car le terme $\boldsymbol{\omega}\times\bm{U}$
pouvait transf\'erer de l'\'energie d'une ligne \`a l'autre. En l'absence de vorticit\'e,
ce couplage dispara\^it~: toutes les lignes de courant partagent la m\^eme \'energie. On
exploitera ce r\'esultat dans la section~5.7 pour calculer les forces de pression sur un
profil d'aile.

% ------------------------------------------
\section{\'Equation de transport de la vorticit\'e}
% ------------------------------------------

On sait maintenant ce qu'est la vorticit\'e et ce qu'elle implique (Bernoulli fort quand
elle est nulle). La question naturelle est~: comment \'evolue-t-elle~? Sous quelles conditions
dispara\^it-elle, se conserve-t-elle, ou s'amplifie-t-elle~? Pour r\'epondre, on d\'erive
une \'equation d'\'evolution pour $\boldsymbol{\omega}$ directement \`a partir des \'equations
de Navier-Stokes.

\subsection{D\'erivation}

On part de la forme de Crocco (g\'en\'eralis\'ee au cas visqueux)~:
\[
  \pd{\bm{U}}{t} + \boldsymbol{\omega}\times\bm{U}
  = -\frac{1}{\rho}\nabla H + \nu\nabla^2\bm{U}.
\]
On prend le \emph{rotationnel} des deux membres. Le terme de pression dispara\^it
imm\'ediatement car $\nabla\times\nabla H = \bm{0}$ (rotationnel d'un gradient). Pour le
terme visqueux, $\nabla\times(\nu\nabla^2\bm{U}) = \nu\nabla^2(\nabla\times\bm{U})
= \nu\nabla^2\boldsymbol{\omega}$ (les op\'erateurs commutent pour un fluide incompressible).
Le membre de gauche donne $\partial_t\boldsymbol{\omega}
+ \nabla\times(\boldsymbol{\omega}\times\bm{U})$.

Il reste \`a d\'evelopper $\nabla\times(\boldsymbol{\omega}\times\bm{U})$. En utilisant
l'identit\'e vectorielle
$\nabla\times(\bm{A}\times\bm{B}) = \bm{A}(\nabla\cdot\bm{B})
- \bm{B}(\nabla\cdot\bm{A}) + (\bm{B}\cdot\nabla)\bm{A} - (\bm{A}\cdot\nabla)\bm{B}$,
avec $\bm{A}=\boldsymbol{\omega}$ et $\bm{B}=\bm{U}$, et en utilisant
$\nabla\cdot\bm{U}=0$ (incompressibilit\'e) et $\nabla\cdot\boldsymbol{\omega}=0$
(toujours vrai car divergence d'un rotationnel), on obtient~:
\[
  \nabla\times(\boldsymbol{\omega}\times\bm{U})
  = (\bm{U}\cdot\nabla)\boldsymbol{\omega} - (\boldsymbol{\omega}\cdot\nabla)\bm{U}.
\]

En rassemblant tous les termes~:

\begin{importantbox}{\'Equation de transport de la vorticit\'e}
\[
  \pd{\boldsymbol{\omega}}{t} + \bm{U}\cdot\nabla\boldsymbol{\omega}
  = \underbrace{\boldsymbol{\omega}\cdot\nabla\bm{U}}_{\text{\'etirement/basculement}}
  + \underbrace{\nu\nabla^2\boldsymbol{\omega}}_{\text{diffusion visqueuse}}
\]
Le membre de gauche est la d\'eriv\'ee mat\'erielle $D\boldsymbol{\omega}/Dt$~: c'est la
variation de vorticit\'e vue en suivant une particule fluide.
\end{importantbox}

\subsection{Interpr\'etation physique des termes}

Les termes du membre de droite ont des significations physiques distinctes~:

\medskip
\noindent\textbf{1. Convection} (membre de gauche, $\bm{U}\cdot\nabla\boldsymbol{\omega}$)
--- La vorticit\'e est advect\'ee par l'\'ecoulement, transport\'ee d'un point \`a l'autre
comme n'importe quelle propri\'et\'e mat\'erielle. C'est ce terme qui explique que les
tourbillons de saumon du Rafale sont emport\'es vers l'aval.

\medskip
\noindent\textbf{2. \'Etirement et basculement} ($\boldsymbol{\omega}\cdot\nabla\bm{U}$)
--- Terme purement \textbf{tridimensionnel} (nul en 2D). Imaginons un tube de vorticit\'e
\'etir\'e par l'\'ecoulement~: par conservation du flux de vorticit\'e $\Gamma = \omega\cdot S$,
si la section $S$ diminue, l'intensit\'e $\omega$ augmente. C'est l'analogue exact du
patineur qui ram\`ene les bras et tourne plus vite. Ce m\'ecanisme d'\'etirement est
le principal responsable de l'amplification de la vorticit\'e dans les \'ecoulements
turbulents tridimensionnels.

\medskip
\noindent\textbf{3. Diffusion visqueuse} ($\nu\nabla^2\boldsymbol{\omega}$) --- La viscosit\'e
\'etale la vorticit\'e spatialement, exactement comme la chaleur diffuse dans un solide
(m\^eme forme math\'ematique). Le temps caract\'eristique de diffusion sur une longueur
$\ell$ est $\ell^2/\nu$. C'est ce terme qui \'epaissit progressivement le c\oe ur des
tourbillons au fil du temps.

\medskip
\noindent\textbf{Cas particuliers importants.}
\begin{itemize}
  \item Fluide parfait ($\nu=0$) en 2D~: pas d'\'etirement, pas de diffusion.
    La vorticit\'e de chaque particule est conserv\'ee~: $\DD{\omega_3} = 0$.
  \item Fluide parfait en 3D~: la vorticit\'e peut \^etre amplifi\'ee par \'etirement,
    m\^eme sans viscosit\'e.
  \item Fluide visqueux en 2D~: la vorticit\'e ob\'eit \`a une \'equation de
    convection-diffusion scalaire (section suivante).
\end{itemize}

% ------------------------------------------
\section{\'Ecoulement 2D et tourbillon de Lamb-Oseen}
% ------------------------------------------

Le cas bidimensionnel est particuli\`erement instructif car l'\'equation de transport
se simplifie consid\'erablement. En deux dimensions, le terme d'\'etirement
$\boldsymbol{\omega}\cdot\nabla\bm{U}$ est nul (la vorticit\'e $\boldsymbol{\omega}
= \omega_3\bm{e}_3$ est perpendiculaire au plan de l'\'ecoulement et ne peut pas
\^etre \'etir\'ee par les gradients de vitesse dans ce plan). L'\'equation de transport
se r\'eduit \`a un bilan scalaire~:
\[
  \pd{\omega_3}{t} + \bm{U}\cdot\nabla\omega_3 = \nu\nabla^2\omega_3.
\]
C'est une \'equation de convection-diffusion pour $\omega_3$, formellement identique
\`a l'\'equation de la chaleur avec advection. Sans viscosit\'e, $D\omega_3/Dt = 0$~:
chaque particule conserve sa vorticit\'e. Avec viscosit\'e, la vorticit\'e s'\'etale
progressivement --- exactement comme une tache d'encre diffuse dans l'eau.

\subsection{Solution de Lamb-Oseen}

On part d'un tourbillon ponctuel de circulation $\Gamma$ \`a l'instant $t=0$
(vorticit\'e concentr\'ee en un point), et on cherche
comment la viscosit\'e le diffuse au cours du temps. C'est un probl\`eme mod\`ele
fondamental, analogue \`a la solution de Green pour l'\'equation de la chaleur.

En sym\'etrie axiale ($\omega_3$ ne d\'epend que de $r$
et $t$, pas de convection radiale par sym\'etrie), l'\'equation se r\'eduit \`a
\[
  \pd{\omega_3}{t} = \nu\left(\frac{\partial^2\omega_3}{\partial r^2}
  + \frac{1}{r}\pd{\omega_3}{r}\right).
\]
C'est l'\'equation de la chaleur en coordonn\'ees cylindriques (sans source). La condition
initiale est $\omega_3(r,0) = \Gamma\,\delta(\bm{x})$, et la condition aux limites
$\omega_3 \to 0$ quand $r\to\infty$. Par analogie avec la solution fondamentale de la
chaleur en 2D, on trouve la solution gaussienne~:

\begin{importantbox}{Tourbillon de Lamb-Oseen}
\[
  \omega_z(r,t) = \frac{\Gamma}{4\pi\nu t}\exp\!\left(-\frac{r^2}{4\nu t}\right),
  \qquad
  u_\theta(r,t) = \frac{\Gamma}{2\pi r}\left[1 - \exp\!\left(-\frac{r^2}{4\nu t}\right)\right].
\]
Le rayon caract\'eristique du c\oe ur vortical cro\^it comme~:
$r_c(t) = \sqrt{4\nu t} \propto \sqrt{t}$.
\end{importantbox}

Ce r\'esultat illustre la diffusion visqueuse de la vorticit\'e~: le tourbillon initialement
singulier s'\'etale progressivement, mais la circulation totale $\Gamma$ reste conserv\'ee
pour tout $t > 0$. C'est ce m\'ecanisme qui explique la longue persistance des tourbillons
de saumon derri\`ere les avions~--- sur le Rafale Marine en approche, ces tourbillons peuvent
persister plusieurs dizaines de secondes apr\`es le passage de l'appareil, ce qui impose des
distances de s\'eparation minimales sur le pont du porte-avions.

\medskip
Un exemple spectaculaire~: les anneaux de vorticit\'e \'eject\'es par l'Etna atteignent
${\sim}200\,\mathrm{m}$ de diam\`etre et montent jusqu'\`a $1000\,\mathrm{m}$ d'altitude,
leur c\oe ur s'\'etalant lentement par diffusion visqueuse.

%============================================================
\section{Vorticit\'e produite par les parois}
%============================================================

On a \'etabli comment la vorticit\'e \'evolue (convection, \'etirement, diffusion). Mais
d'o\`u vient-elle initialement~? Un \'ecoulement uniforme $\bm{U} = U_\infty\bm{e}_1$
est irrotationnel ($\boldsymbol{\omega} = \bm{0}$). L'\'equation de transport montre que
si $\boldsymbol{\omega} = \bm{0}$ partout \`a $t=0$ et qu'il n'y a aucune source, la
vorticit\'e reste nulle pour tout $t$. Alors d'o\`u viennent les tourbillons~?

La r\'eponse est~: des \textbf{parois}. En pr\'esence de la condition de non-glissement
(chapitre~2), le fluide adh\`ere \`a la surface solide~: $\bm{U} = \bm{0}$ \`a la paroi.
La vitesse doit passer de z\'ero (paroi) \`a $U_\infty$ (\'ecoulement externe) sur
l'\'epaisseur $\delta$ de la couche limite. Ce gradient de vitesse \'enorme
($\partial U_1/\partial x_2 \sim U_\infty/\delta$) engendre une vorticit\'e intense
pr\`es de la paroi.

Consid\'erons une paroi plane en $x_2 = 0$, avec $U_1(x_1,0) = 0$. La contrainte de
cisaillement vaut $\tau = \mu\,\partial U_1/\partial x_2$, et la vorticit\'e proche paroi est
\[
  \omega_3\big|_{\text{paroi}} = \pd{U_2}{x_1} - \pd{U_1}{x_2}
  \approx -\pd{U_1}{x_2} = -\frac{\tau}{\mu}.
\]
En vecteur, on peut \'ecrire de fa\c con \'el\'egante~: pour une paroi de normale
$\boldsymbol{n}$,
\[
  \boldsymbol{\tau} = \mu\,\boldsymbol{\omega} \times \boldsymbol{n}.
\]

\begin{importantbox}{Origine pari\'etale de la vorticit\'e}
La vorticit\'e est \textbf{cr\'e\'ee exclusivement aux parois} (par le cisaillement visqueux),
puis diffus\'ee et convect\'ee dans le fluide. Dans un \'ecoulement sans paroi (fluide infini,
$\boldsymbol{\omega}=\bm{0}$ \`a $t=0$), la vorticit\'e reste nulle pour tout $t$.

Ce r\'esultat justifie \emph{a posteriori} le mod\`ele d'\'ecoulement potentiel (irrotationnel)
utilis\'e loin des parois~: la vorticit\'e n'a aucune raison d'appara\^itre l\`a o\`u elle
n'a pas \'et\'e produite.
\end{importantbox}

Sur le Rafale Marine, la couche limite sur l'aile est la source primaire de vorticit\'e~: elle
na\^it au bord d'attaque sous l'effet du cisaillement visqueux, est convect\'ee jusqu'au bord
de fuite, et se retrouve \emph{in fine} dans le sillage et les tourbillons marginaux.
La g\'eom\'etrie de l'aile delta \`a double fl\`eche favorise la formation de tourbillons de
bord d'attaque \`a grande incidence --- une source de portance non lin\'eaire absente des
ailes conventionnelles.

% ------------------------------------------
\section{\'Ecoulement autour d'un corps \`a grand Reynolds}
% ------------------------------------------

On dispose maintenant de tous les outils pour comprendre la structure de l'\'ecoulement
autour d'un corps \`a grand nombre de Reynolds~: la vorticit\'e est confin\'ee dans la
couche limite (cr\'e\'ee aux parois, diffus\'ee sur $\delta \ll L$) et dans le sillage
(convect\'ee depuis le bord de fuite). Le reste du domaine est irrotationnel.
On peut donc mod\'eliser l'\'ecoulement ext\'erieur comme \emph{potentiel} (section~5.7),
et la couche limite comme une fine nappe de vorticit\'e.

\subsection{Corps non profil\'e}

\`A grand Reynolds, la couche limite se d\'ecolle t\^ot sur un corps non profil\'e (sph\`ere,
cylindre), cr\'eant un sillage turbulent large. La vorticit\'e produite \`a la paroi est
inject\'ee dans ce sillage~: elle s'organise en all\'ee de tourbillons altern\'es (all\'ee de
B\'enard-K\'arm\'an, vue au chapitre~4), source d'une tra\^in\'ee \'elev\'ee et de vibrations.

\subsection{Corps profil\'e}

Sur un corps profil\'e (\`a faible incidence), la couche limite reste attach\'ee jusqu'au bord
de fuite~: le sillage est mince, la traîn\'ee est faible. La vorticit\'e produite \`a la paroi
est convect\'ee dans ce sillage~; elle peut \^etre mod\'elis\'ee comme une \textbf{nappe de
vorticit\'e} $\mathcal{S}_\omega$ d'\'epaisseur nulle, qui s'\'ecoule depuis le bord de fuite.

\subsection{Mod\'elisation par nappe de vorticit\'e}

\`A la limite $Re \to \infty$, la couche limite est infiniment mince et on la remplace par une
nappe de vorticit\'e $\mathcal{S}_\omega$ (une surface de discontinuit\'e de la vitesse
tangentielle). L'intensit\'e locale de la nappe est
\[
  \gamma_s = U_1^+ - U_1^-, \qquad \omega_3 = \gamma_s\,\delta(x_2),
\]
o\`u $U_1^\pm$ sont les vitesses tangentielles de part et d'autre de la nappe. Les conditions
de raccord \`a la nappe sont~:
\[
  p_A = p_B \quad \text{(continuit\'e de la pression)},
  \qquad \boldsymbol{U}_A\cdot\boldsymbol{n} = \boldsymbol{U}_B\cdot\boldsymbol{n}
  \quad \text{(imperm\'eabilit\'e)}.
\]
Le syst\`eme \`a r\'esoudre dans le reste du domaine est~:
\[
  \begin{cases} \displaystyle\pd{\boldsymbol{U}}{t} + \boldsymbol{U}\cdot\nabla\boldsymbol{U}
  = -\dfrac{1}{\rho}\nabla p \\[6pt] \nabla\cdot\boldsymbol{U} = 0 \end{cases}
\]
avec les conditions $\boldsymbol{U}\cdot\boldsymbol{n}=0$ sur la surface $\mathcal{S}$ du
profil et $\boldsymbol{U}\to\boldsymbol{U}_\infty$ \`a l'amont.

% ------------------------------------------
\section{Bases de l'a\'erodynamique}
% ------------------------------------------

On arrive au c\oe ur de ce chapitre~: le calcul de la portance d'un profil d'aile.
L'id\'ee est d'exploiter conjointement deux r\'esultats \'etablis plus haut~:
(1)~l'\'ecoulement ext\'erieur est irrotationnel (vorticit\'e confin\'ee aux parois),
et (2)~le Bernoulli fort s'applique dans tout le domaine irrotationnel. On va montrer
que le probl\`eme se ram\`ene \`a r\'esoudre une \'equation de Laplace, et que la
portance d\'epend d'un seul param\`etre~: la circulation $\Gamma$ autour du profil.

\subsection{\'Ecoulement potentiel autour d'un profil d'aile}

On consid\`ere un \'ecoulement 2D stationnaire autour d'un profil d'aile \`a grand nombre de
Reynolds~: $\boldsymbol{U} = (U_1(x_1,x_2),\, U_2(x_1,x_2))$.

La continuit\'e de la pression \`a travers la nappe $\mathcal{S}_\omega$ impose, via le
Bernoulli fort, que la norme de la vitesse $|\boldsymbol{U}|$ soit continue au travers de
$\mathcal{S}_\omega$. Combin\'e avec la condition $\boldsymbol{U}\cdot\boldsymbol{n}$ continue,
on conclut que le \emph{vecteur} $\boldsymbol{U}$ lui-m\^eme est continu~: la nappe de
vorticit\'e peut \^etre retir\'ee du probl\`eme, et \textbf{l'\'ecoulement est irrotationnel
dans tout le domaine}.

\medskip
\noindent\textbf{Fonction de courant.}
Pour un \'ecoulement 2D incompressible, l'\'equation de continuit\'e
$\partial U_1/\partial x_1 + \partial U_2/\partial x_2 = 0$ est automatiquement satisfaite
en introduisant la \textbf{fonction de courant} $\psi$ d\'efinie par
\[
  U_1 = \pd{\psi}{x_2}, \qquad U_2 = -\pd{\psi}{x_1}.
\]
La vorticit\'e s'exprime alors comme~:
\[
  \omega_3 = \pd{U_2}{x_1} - \pd{U_1}{x_2}
  = -\frac{\partial^2\psi}{\partial x_1^2} - \frac{\partial^2\psi}{\partial x_2^2}
  = -\nabla^2\psi.
\]
La propri\'et\'e $\boldsymbol{U}\cdot\nabla\psi = 0$ montre que $\psi$ est constante le long
de chaque ligne de courant.

L'irrotationnalité ($\omega_3 = 0$) entra\^ine alors $\nabla^2\psi = 0$. Le probl\`eme se
ram\`ene \`a~:

\begin{importantbox}{\'Ecoulement potentiel 2D autour d'un profil}
\[
  \nabla^2\psi = \frac{\partial^2\psi}{\partial x_1^2}
  + \frac{\partial^2\psi}{\partial x_2^2} = 0,
  \qquad \psi = 0 \text{ sur } \mathcal{S}, \qquad
  \psi \to U_\infty x_2 \text{ loin du profil.}
\]
C'est l'\'equation de Laplace avec des conditions aux limites mixtes~: bien pos\'ee
num\'eriquement, r\'esolvable par \'el\'ements finis ou panels.
\end{importantbox}

Une fois $\psi$ d\'etermin\'ee, le champ de vitesse s'obtient par d\'erivation et la pression
via le Bernoulli fort $p + \tfrac{1}{2}\rho U^2 = \mathrm{cst}$.

\subsection{Th\'eor\`eme de Kutta-Joukowski}

Sans conna\^itre la g\'eom\'etrie pr\'ecise du profil, on peut d\'eterminer la \emph{forme
universelle} de la solution loin du profil. On passe en coordonn\'ees polaires $(r,\theta)$~:
\[
  \frac{1}{r}\pd{}{r}\!\left(r\pd{\psi}{r}\right)
  + \frac{1}{r^2}\frac{\partial^2\psi}{\partial\theta^2} = 0.
\]
La solution $2\pi$-p\'eriodique en $\theta$ se cherche comme une s\'erie de Fourier
$\psi = \sum_{n=-\infty}^{+\infty} C_n(r)\,e^{in\theta}$, o\`u chaque $C_n$ v\'erifie
\[
  r\frac{d}{dr}\!\left(r\frac{dC_n}{dr}\right) - n^2 C_n = 0
  \quad\Longrightarrow\quad
  \begin{cases} C_0 = A_0\ln r + B_0 & (n=0) \\
  C_n = A_n r^{|n|} + B_n r^{-|n|} & (n\ne 0). \end{cases}
\]
Les termes en $r^{|n|}$ (rouges dans les slides) contribuent uniquement au champ proche,
li\'e \`a la g\'eom\'etrie particuli\`ere du profil~: ils sont n\'egligeables loin du profil.
Le raccordement \`a l'\'ecoulement uniforme \`a l'infini,
$\psi_\infty = U_\infty x_2 = U_\infty r\sin\theta$, impose
\[
  A_1 = A_{-1} = \frac{U_\infty}{2i}, \quad A_n = 0 \text{ pour } |n|>1.
\]
La fonction de courant loin du profil se r\'eduit donc \`a
\[
  \psi = U_\infty r\sin\theta + A_0\ln r + B_0 + \mathcal{O}(r^{-1}).
\]
On en d\'eduit le champ de vitesse en coordonn\'ees cylindriques
($U_r = (1/r)\partial\psi/\partial\theta$, $U_\theta = -\partial\psi/\partial r$)~:
\[
  \boldsymbol{U} = U_\infty\cos\theta\,\boldsymbol{e}_r
  + \left(-U_\infty\sin\theta - \frac{A_0}{r}\right)\boldsymbol{e}_\theta
  + \mathcal{O}(r^{-2})
  = \boldsymbol{U}_\infty - \frac{A_0}{r}\,\boldsymbol{e}_\theta + \mathcal{O}(r^{-2}).
\]
Le terme en $A_0/r$ est exactement le champ de vitesse d'un tourbillon ponctuel (vu en
\S5.2)~! La circulation $\Gamma$ autour d'un cercle $\mathcal{C}$ de rayon $R$ entourant le
profil (parcouru dans le sens horaire) est~:
\[
  \Gamma = \oint_{\mathcal{C}} \boldsymbol{U}\cdot d\boldsymbol{x}
  = -\int_0^{2\pi} \left(-\frac{A_0}{R}\,\boldsymbol{e}_\theta\right)\cdot R\,d\theta\,
  \boldsymbol{e}_\theta = 2\pi A_0.
\]
Ainsi $A_0 = \Gamma/2\pi$, et loin du profil~:
\[
  \boldsymbol{U} = \boldsymbol{U}_\infty - \frac{\Gamma}{2\pi r}\,\boldsymbol{e}_\theta
  + \mathcal{O}(r^{-2}).
\]
\textbf{Le profil se comporte comme un tourbillon ponctuel de circulation $\Gamma$ en champ
lointain}, quelle que soit sa forme pr\'ecise. La pression, via le Bernoulli fort, vaut~:
\[
  p = p_\infty + \frac{\rho\Gamma}{2\pi r}\,\boldsymbol{U}_\infty\cdot\boldsymbol{e}_\theta
  + \mathcal{O}(r^{-2}).
\]

\medskip
\noindent\textbf{Calcul de la force.}
On consid\`ere le domaine $\mathcal{D}$ born\'e par la surface du profil $\mathcal{S}$ et
un grand cercle $\mathcal{S}^\star$ de rayon $R$. Pour un \'ecoulement stationnaire non
visqueux, le th\'eor\`eme de Reynolds appliqu\'e \`a la quantit\'e de mouvement donne~:
\[
  \int_{\mathcal{S}\cup\mathcal{S}^\star}
  \left(p\,\boldsymbol{n} + \rho\boldsymbol{U}(\boldsymbol{U}\cdot\boldsymbol{n})\right)dS
  = \boldsymbol{0}.
\]
La force $\boldsymbol{F}$ exerc\'ee sur le profil (par unit\'e d'envergure) est donc
calcul\'ee sur le grand cercle $\mathcal{S}^\star$ (avec $\boldsymbol{n} = -\boldsymbol{e}_r$)~:
\[
  \boldsymbol{F}
  = -\int_0^{2\pi}\Big[p\,\boldsymbol{e}_r
  + \rho\boldsymbol{U}({\boldsymbol{U}\cdot\boldsymbol{e}_r})\Big]_{r=R} R\,d\theta.
\]
En substituant les expressions asymptotiques de $p$ et $\boldsymbol{U}$ \`a grand $R$
(termes crois\'es $\propto U_\infty\Gamma/R$ dominant), les calculs donnent apr\`es
int\'egration (les termes en $U_\infty^2$ et $\Gamma^2/R^2$ s'annulent par sym\'etrie)~:
\[
  F_1 = 0, \qquad F_2 = \rho\Gamma U_\infty.
\]

\begin{importantbox}{Th\'eor\`eme de Kutta-Joukowski}
La force exerc\'ee sur un profil d'aile en \'ecoulement potentiel 2D stationnaire est
perpendiculaire \`a la vitesse amont~:
\[
  \boldsymbol{F} = \rho\,\Gamma\,U_\infty\,\boldsymbol{e}_2
\]
Il n'y a \textbf{pas de tra\^in\'ee} (paradoxe de d'Alembert) et la \textbf{portance} est
directement proportionnelle \`a la circulation $\Gamma$ autour du profil.
\end{importantbox}

Ce r\'esultat est fondamental~: il montre que la portance n'est pas une propri\'et\'e
intrins\`eque du profil, mais une cons\'equence de la \emph{circulation} que l'\'ecoulement
acquiert autour de lui. Sur le Rafale Marine, les tourbillons de saumon visibles \`a
l'appontage t\'emoignent justement de cette circulation~--- les volutes de vapeur d'eau
photographi\'ees derri\`ere les ailes sont des lignes de courant qui s'enroulent autour des
tourbillons marginaux, eux-m\^emes li\'es \`a la circulation g\'en\'eratrice de portance.
La condition de Kutta (la vitesse reste finie au bord de fuite) impose la valeur de $\Gamma$
et ferme le probl\`eme.

\begin{center}
\begin{tikzpicture}[scale=1.0]
  % Airfoil schematic
  \fill[gray!20] plot[smooth, tension=0.7] coordinates
    {(-2,0) (-1.5,0.35) (-0.5,0.5) (0.5,0.4) (1.5,0.15) (2,0) (1.5,-0.1)
     (0.5,-0.15) (-0.5,-0.1) (-1.5,-0.05) (-2,0)};
  \draw[thick] plot[smooth, tension=0.7] coordinates
    {(-2,0) (-1.5,0.35) (-0.5,0.5) (0.5,0.4) (1.5,0.15) (2,0)};
  \draw[thick] plot[smooth, tension=0.7] coordinates
    {(-2,0) (-1.5,-0.05) (-0.5,-0.1) (0.5,-0.15) (1.5,-0.1) (2,0)};
  % Incoming flow
  \foreach \y in {-0.8, -0.4, 0, 0.4, 0.8} {
    \draw[->, blue!70] (-3.5,\y) -- (-2.3,\y);
  }
  \node[blue] at (-3.0, 1.1) {$\boldsymbol{U}_\infty$};
  % Circulation arrow
  \draw[->, red!80, thick] (0,1.5) arc (90:-270:0.5cm);
  \node[red] at (1.2,1.5) {$\Gamma$};
  % Force arrow
  \draw[->, very thick, green!60!black] (0,0) -- (0,1.2)
    node[right] {$\boldsymbol{F} = \rho\Gamma U_\infty\boldsymbol{e}_2$};
  % vortex sheet indication
  \draw[dashed, orange] (2,0) -- (3.5,-0.1) node[right, orange, font=\small]
    {sillage $\mathcal{S}_\omega$};
\end{tikzpicture}
\end{center}

\subsection{Effet Magnus}

La force de Kutta-Joukowski est aussi connue sous le nom de \textbf{force de Magnus} (1852)
dans le cas d'un objet en rotation sur lui-m\^eme (vitesse angulaire $\boldsymbol{\Omega}$).
Un objet tournant entra\^ine le fluide par viscosit\'e, cr\'eant une circulation
$\Gamma \simeq \Omega\, S$ (o\`u $S \propto L^2$ est la section transversale). La force de
Magnus est donc~:
\[
  \boldsymbol{F}_M \propto \rho L^3\,\boldsymbol{\Omega}\times\boldsymbol{U}.
\]

\medskip
\noindent\textbf{Trajectoire d'une balle en rotation.}
Pour une balle de masse $m$ et de rayon $a$ tournant \`a $\boldsymbol{\Omega}$, le bilan des
forces donne~:
\[
  m\ddot{\boldsymbol{x}} = m\boldsymbol{g}
  \underbrace{-C_D\,\tfrac{1}{2}\rho|\dot{\boldsymbol{x}}|^2 S\,
  \frac{\dot{\boldsymbol{x}}}{|\dot{\boldsymbol{x}}|}}_{\text{tra\^in\'ee}}
  + \underbrace{C_M\,\rho\pi a^3\,\boldsymbol{\Omega}\times\dot{\boldsymbol{x}}}_{\text{Magnus}}
\]
o\`u le coefficient $C_M = C_M(\mathrm{Re},\,\Omega a/U)$ d\'epend du r\'egime de rotation.

\begin{importantbox}{Exemple --- coup franc de Roberto Carlos (1997)}
France-Br\'esil \`a Gerland~: balle de rayon $a \approx 11\,\mathrm{cm}$, $\Omega$ tr\`es \'elev\'e.
La balle contourne le mur d\'efensif (trajectoire presque droite~--- l'\'etape laminaire)
avant de d\'evier brutalement par effet Magnus (transition turbulente~: $C_D$ chute,
$C_M$ augmente) pour finir dans le but. L'\'ecart lateral est de l'ordre de $3\,\mathrm{m}$.
\end{importantbox}

L'effet Magnus s'applique aussi au frisbee, au golf (effet de baguette avec spin arri\`ere) et
\`a certains missiles gyrostabilis\'es. Sur le Rafale Marine, le fuselage et les voilures
peuvent induire des effets de circulation crois\'ee \`a forte incidence, mais la ma\^itrise
de l'\'ecoulement passe principalement par le contr\^ole actif des surfaces de gouverne.

% ==========================================
\section*{Résumé du chapitre}
\addcontentsline{toc}{section}{Résumé du chapitre}
% ==========================================

\begin{importantbox}{Ce qu'il faut retenir}
\begin{itemize}
  \item \textbf{Vorticit\'e}~: $\boldsymbol{\omega} = \nabla\times\bm{U}$ mesure
        la rotation locale du fluide (deux fois la vitesse angulaire). En 2D~:
        $\omega_3 = \partial U_2/\partial x_1 - \partial U_1/\partial x_2$.

  \item \textbf{Circulation}~:
        $\Gamma_{\mathcal{C}} = \oint_{\mathcal{C}}\bm{U}\cdot d\bm{x}
        = \iint\boldsymbol{\omega}\cdot\bm{n}\,dS$ (th\'eor\`eme de Stokes).

  \item \textbf{Bernoulli fort}~: pour un \'ecoulement irrotationnel, stationnaire et
        parfait, $H = p + \tfrac{1}{2}\rho U^2 + \rho g x_3 = \mathrm{cst}$
        \emph{partout} dans le domaine (pas seulement le long d'une ligne de courant).

  \item \textbf{\'Equation de transport}~:
        $\dfrac{D\boldsymbol{\omega}}{Dt}
        = \boldsymbol{\omega}\cdot\nabla\bm{U} + \nu\nabla^2\boldsymbol{\omega}$.
        Convection + \'etirement (3D) + diffusion visqueuse.

  \item \textbf{Tourbillon de Lamb-Oseen}~: diffusion visqueuse d'un tourbillon
        ponctuel, $r_c(t) = \sqrt{4\nu t}$. La circulation totale est conserv\'ee.

  \item \textbf{Origine pari\'etale}~: la vorticit\'e est cr\'e\'ee exclusivement aux
        parois par cisaillement visqueux, puis convect\'ee et diffus\'ee. Sans paroi
        et sans vorticit\'e initiale, l'\'ecoulement reste irrotationnel.

  \item \textbf{\'Ecoulement potentiel}~: \`a grand Re, l'\'ecoulement ext\'erieur
        v\'erifie $\nabla^2\psi = 0$ (Laplace). La couche limite et le sillage sont
        mod\'elis\'es comme des nappes de vorticit\'e.

  \item \textbf{Kutta-Joukowski}~: $\bm{F} = \rho\Gamma U_\infty\bm{e}_2$.
        La portance est proportionnelle \`a la circulation~; il n'y a pas de tra\^in\'ee
        en \'ecoulement potentiel (paradoxe de d'Alembert).

  \item \textbf{Effet Magnus}~: un corps en rotation cr\'ee une circulation et donc une
        force lat\'erale, $\bm{F}_M \propto \rho L^3\boldsymbol{\Omega}\times\bm{U}$.
\end{itemize}
\end{importantbox}

% ==========================================
\chapter{\'Ecoulement turbulent}
% ==========================================

Jusqu'ici, tous les \'ecoulements \'etudi\'es \'etaient \emph{laminaires}~: les particules
fluides suivent des trajectoires r\'eguli\`eres, le champ de vitesse est stationnaire (ou
pr\'evisiblement instationnaire), et les \'equations de Navier-Stokes admettent des
solutions analytiques exactes (Couette, Poiseuille, Blasius). Mais en pratique, d\`es que
le nombre de Reynolds d\'epasse quelques milliers, ces solutions laminaires deviennent
instables~: l'\'ecoulement \emph{transite} vers un r\'egime fondamentalement diff\'erent,
la \textbf{turbulence}.

Sur le Rafale Marine en croisi\`ere ($\mathrm{Re} \approx 7\times10^7$), la totalit\'e de
la couche limite sur l'aile est turbulente~: les fluctuations de vitesse d\'epassent
$10\,\%$ de la vitesse de vol, les structures actives s'\'etendent de l'envergure
($\sim 1\,\mathrm{m}$) aux plus petits tourbillons dissipatifs (quelques microm\`etres).
Comprendre --- et mod\'eliser --- cette complexit\'e est au c\oe ur de la conception
a\'erodynamique moderne.

Ce chapitre introduit les outils statistiques n\'ecessaires (d\'ecomposition de Reynolds,
\'equations RANS) et la physique de la cascade d'\'energie (Kolmogorov, 1941). On terminera
par une estimation du co\^ut prohibitif de la simulation directe, qui justifie le recours
aux mod\`eles de turbulence en ing\'enierie.

\begin{defbox}{Fil conducteur}
La turbulence est le r\'egime \emph{naturel} de l'a\'erodynamique~: aucun avion ne vole
en r\'egime laminaire sur toute sa voilure. Comprendre la turbulence, c'est comprendre
pourquoi la couche limite du Rafale reste attach\'ee \`a forte incidence (retard du
d\'ecollement), mais aussi pourquoi sa tra\^in\'ee de frottement est plus \'elev\'ee qu'en
laminaire --- un compromis que l'ing\'enieur doit ma\^itriser.
\end{defbox}

% ------------------------------------------
\section{Ph\'enom\'enologie de la turbulence}
% ------------------------------------------

Observons un jet d'eau ou de fum\'ee~: \`a faible d\'ebit, le jet est laminaire, ses lignes de
courant r\'eguli\`eres et pr\'evisibles. Au-del\`a d'un certain d\'ebit, le jet se brise en
structures tourbillonnaires d\'esordonn\'ees qui fluctuent en permanence --- c'est la turbulence.
Trois traits la caract\'erisent~:

\medskip
\noindent\textbf{1.~Mouvement instationnaire et ap\'eriodique.}
Le champ de vitesse varie de fa\c con erratique. Dans la couche de cisaillement d'un jet
($D = 50\,\mathrm{mm}$, $U_j = 30\,\mathrm{m\,s^{-1}}$, $\mathrm{Re}_D = 10^5$), la
fluctuation longitudinale $u_1'(t)$ atteint $\pm 16\,\mathrm{m\,s^{-1}}$, soit plus de
50\,\% de la vitesse du jet, sans aucune p\'eriodicit\'e identifiable.

\medskip
\noindent\textbf{2.~Sensibilit\'e aux conditions initiales.}
Deux particules fluides initialement voisines divergent exponentiellement --- la turbulence
est chaotique. L'exp\'erience du \emph{blinking vortex} d'Aref (1984) le montre~: deux
tourbillons ponctuels alternant p\'eriodiquement (param\`etre de contr\^ole
$\beta = \Gamma T / 2\pi a^2$) g\'en\`erent un m\'elange chaotique d\`es $\beta \gtrsim 0.25$.
M\^eme un \'ecoulement simple peut donc \^etre impr\'evisible \`a long terme.

\medskip
\noindent\textbf{3.~Large gamme d'\'echelles spatio-temporelles.}
L'\'ecoulement turbulent contient simultan\'ement de grandes structures coh\'erentes (taille $\sim L$,
\'energie cin\'etique \'elev\'ee) et de minuscules filaments de vorticit\'e (taille $\sim l_\eta$)
o\`u l'\'energie est finalement dissip\'ee par viscosit\'e. C'est cette multi-\'echelle qui rend la
turbulence si riche et si difficile \`a simuler.

\medskip
\noindent\textbf{La turbulence n'est pas qu'une nuisance.}
Elle augmente le m\'elange et les transferts de chaleur, et --- point crucial pour
l'a\'erodynamique --- retarde le d\'ecollement des couches limites. Une couche limite turbulente,
dont le profil de vitesse est plus plein qu'une couche laminaire, r\'esiste mieux aux gradients
de pression adverses (comme sur le bord de fuite d'une aile \`a forte incidence). Sur le
Rafale, cette propri\'et\'e contribue \`a maintenir la portance jusqu'\`a des angles d'attaque \'elev\'es
lors de la mise en virage ou de l'appontage. La contrepartie est une augmentation de la
tra\^in\'ee de frottement pari\'etal.

% ------------------------------------------
\section{Description statistique}
% ------------------------------------------

La turbulence \'etant fondamentalement al\'eatoire, une description d\'eterministe
instantan\'ee du champ de vitesse est hors de port\'ee (et sans grand int\'er\^et
pratique). Il faut lui donner un cadre \textbf{statistique} rigoureux~: c'est l'objet
de la d\'ecomposition de Reynolds et des \'equations RANS.

Le champ instantan\'e $U_i(\bm{x},t)$ contient une information colossale --- trop pour
\^etre utile. Un ing\'enieur chez Dassault ne veut pas conna\^itre la vitesse \`a chaque
microseconde en chaque point~: il veut la tra\^in\'ee moyenne, la portance moyenne, le
frottement pari\'etal moyen. L'id\'ee est donc de s\'eparer ce qui est pr\'evisible (le champ
moyen) de ce qui ne l'est pas (les fluctuations). C'est exactement ce que fait la
d\'ecomposition de Reynolds.

\subsection{Moyenne et d\'ecomposition de Reynolds}

On d\'efinit la \textbf{moyenne d'ensemble} d'une variable $f(\bm{x},t)$
comme la limite de la moyenne sur $N$ r\'ealisations identiques du m\^eme \'ecoulement~:
\[
  \bar{F}(\bm{x},t) = \lim_{N\to\infty} \frac{1}{N}\sum_{i=1}^{N} f^{(i)}(\bm{x},t).
\]
Cette d\'efinition est rigoureuse mais difficile \`a impl\'ementer en pratique (et impossible pour
les \'ecoulements g\'eophysiques, qu'on ne peut pas reproduire). Pour une
\textbf{turbulence stationnaire} (dont les statistiques ne d\'ependent pas du temps), le
concept d'ergodicit\'e permet d'approcher la moyenne d'ensemble par une
\textbf{moyenne temporelle} sur une seule r\'ealisation~:
\[
  \bar{F}(\bm{x}) = \lim_{T\to\infty} \frac{1}{T}\int_{t_0}^{t_0+T}
  f(\bm{x},t')\,dt'.
\]

On peut alors caract\'eriser les fluctuations par leurs moments statistiques. Pour une variable
centr\'ee $x_i' = x_i - \bar{X}_i$ de d\'eviation quadratique
$x_{i,\mathrm{rms}} = \sqrt{\overline{x_i'^2}}$, la \textbf{dissym\'etrie} $S$ et
l'\textbf{aplatissement} $T$ (ou kurtosis) sont~:
\[
  S_{x_i} = \frac{\overline{x_i'^3}}{x_{i,\mathrm{rms}}^3},
  \qquad T_{x_i} = \frac{\overline{x_i'^4}}{x_{i,\mathrm{rms}}^4}.
\]
La \textbf{distribution gaussienne} est la r\'ef\'erence naturelle~:
$p(\xi) = \frac{1}{\sqrt{2\pi}\sigma}\exp(-\xi^2/2\sigma^2)$, pour laquelle $S=0$ et $T=3$.

\medskip
La brique fondamentale de l'analyse statistique est la \textbf{d\'ecomposition de Reynolds}~:

\begin{defbox}{D\'ecomposition de Reynolds}
Toute variable $f$ se d\'ecompose en valeur moyenne $\bar{F}$ et fluctuation $f'$~:
\[
  f = \bar{F} + f', \quad \text{avec} \quad \overline{f'} = 0
  \quad (f' = f - \bar{F}).
\]
La r\`egle fondamentale pour un produit $fg$ (avec $f = \bar{F}+f'$ et $g = \bar{G}+g'$)~:
\[
  \overline{fg} = \bar{F}\bar{G} + \overline{f'g'}
\]
car $\overline{f'\bar{G}} = \bar{G}\overline{f'} = 0$ et $\overline{\bar{F}g'} = \bar{F}\overline{g'} = 0$.
\end{defbox}

Appliqu\'ee \`a la vitesse~: $U_i \equiv \bar{U}_i + u_i'$ avec $\overline{u_i'} = 0$.
Le champ $\bar{U}_i$ est la partie \emph{calculable}~; le champ fluctuant $u_i'$ est la
partie qui doit \^etre \emph{mod\'elis\'ee}.

\subsection{\'Equations de Navier-Stokes Moyenn\'ees par Reynolds (RANS)}

On souhaite obtenir des \'equations d'\'evolution pour le champ moyen $\bar{U}_i$.
La strat\'egie est simple~: \emph{substituer} la d\'ecomposition de Reynolds dans les
\'equations de Navier-Stokes, puis \emph{prendre la moyenne}.

On part des NS incompressibles \`a densit\'e constante~:
\[
  \frac{\partial U_i}{\partial x_i} = 0, \qquad
  \frac{\partial(\rho U_i)}{\partial t} + \frac{\partial}{\partial x_j}(\rho U_i U_j)
  = -\frac{\partial p}{\partial x_i} + \frac{\partial \tau_{ij}}{\partial x_j},
  \quad \tau_{ij} = 2\mu D_{ij}.
\]
On pose $U_i = \bar{U}_i + u_i'$, $p = \bar{P} + p'$, $\tau_{ij} = \bar{\tau}_{ij} + \tau'_{ij}$.
L'\'equation de continuit\'e, moyenn\'ee, donne imm\'ediatement~:
\[
  \frac{\partial\bar{U}_i}{\partial x_i} = 0,
\]
et par soustraction $\partial u_i'/\partial x_i = 0$~: \textbf{le champ moyen et le champ
fluctuant sont tous deux \`a divergence nulle}.

D\'eveloppons le terme convectif pas \`a pas. On substitue $U_i = \bar{U}_i + u_i'$ dans le
produit $U_i U_j$~:
\[
  U_i U_j = (\bar{U}_i + u_i')(\bar{U}_j + u_j')
  = \bar{U}_i\bar{U}_j + \bar{U}_i u_j' + u_i'\bar{U}_j + u_i'u_j'.
\]
Prenons maintenant la moyenne de chacun de ces quatre termes~:
\begin{itemize}
  \item $\overline{\bar{U}_i\bar{U}_j} = \bar{U}_i\bar{U}_j$ (les moyennes sont des constantes vis-\`a-vis de l'op\'erateur $\overline{\cdot}$),
  \item $\overline{\bar{U}_i u_j'} = \bar{U}_i\,\overline{u_j'} = 0$ (car $\overline{u_j'}=0$),
  \item $\overline{u_i'\bar{U}_j} = \bar{U}_j\,\overline{u_i'} = 0$ (idem),
  \item $\overline{u_i'u_j'}$ survit~: c'est la \emph{corr\'elation crois\'ee des fluctuations}.
\end{itemize}
Bilan~: seuls le premier et le dernier terme survivent~:
\[
  \overline{U_i U_j} = \bar{U}_i\bar{U}_j + \overline{u_i'u_j'}.
\]
On obtient les \'equations RANS~:

\begin{importantbox}{\'Equations RANS --- Navier-Stokes moyenn\'ees}
\[
  \frac{\partial\bar{U}_i}{\partial x_i} = 0
\]
\[
  \frac{\partial(\rho\bar{U}_i)}{\partial t}
  + \frac{\partial(\rho\bar{U}_i\bar{U}_j)}{\partial x_j}
  = -\frac{\partial\bar{P}}{\partial x_i}
  + \frac{\partial}{\partial x_j}\!\left(
    \bar{\tau}_{ij} - \rho\,\overline{u_i' u_j'}
  \right)
\]
\end{importantbox}

Regardons bien ce r\'esultat. Les \'equations RANS ont \emph{exactement} la m\^eme forme que
les \'equations de Navier-Stokes du chapitre~2 --- \`a un terme pr\`es~:
$-\rho\,\overline{u_i'u_j'}$. Ce terme additionnel, le \textbf{tenseur des contraintes de
Reynolds}, joue le r\^ole d'une contrainte suppl\'ementaire exerc\'ee par les fluctuations
turbulentes sur le champ moyen. C'est comme si la turbulence agissait comme un fluide
visqueux additionnel qui frotte davantage que la viscosit\'e mol\'eculaire seule.

Le tenseur $-\rho\,\overline{u_i'u_j'}$ est sym\'etrique (6 composantes ind\'ependantes),
g\'en\'eralement bien plus grand que $\bar{\tau}_{ij}$
sauf pr\`es des parois, et --- surtout --- \emph{inconnu}~: c'est le
\textbf{probl\`eme de fermeture de la turbulence}. On dispose de 4 \'equations (continuit\'e + 3 QDM)
pour $\bar{U}_i$ et $\bar{P}$, mais les 6 composantes du tenseur de Reynolds sont 6
inconnues suppl\'ementaires. Il faut un \emph{mod\`ele}.

% ------------------------------------------
\section{Mod\'elisation de la turbulence}
% ------------------------------------------

Le probl\`eme de fermeture pos\'e, il faut proposer une relation entre le tenseur de
Reynolds $\overline{u_i'u_j'}$ et les quantit\'es connues (champ moyen $\bar{U}_i$ et ses
d\'eriv\'ees). On introduit d'abord les grandeurs physiques caract\'eristiques de la
turbulence, puis le mod\`ele le plus r\'epandu.

\subsection{\'Energie cin\'etique turbulente et dissipation}

En appliquant la r\`egle du produit \`a l'\'energie cin\'etique totale~:
\[
  \frac{\overline{U_i U_i}}{2} = \frac{\bar{U}_i\bar{U}_i}{2}
  + \frac{\overline{u_i'u_i'}}{2}.
\]
Le deuxi\`eme terme est l'\'energie cin\'etique port\'ee par les fluctuations.

\begin{defbox}{\'Energie cin\'etique turbulente et taux de dissipation}
L'\textbf{\'energie cin\'etique turbulente}~:
\[
  k_t \equiv \frac{\overline{u_i'u_i'}}{2}
  = \frac{1}{2}\!\left(\overline{u_1'^2}+\overline{u_2'^2}+\overline{u_3'^2}\right)
  \quad [\mathrm{m^2\,s^{-2}}].
\]
Le \textbf{taux de dissipation turbulente}~:
\[
  \varepsilon \equiv 2\nu\,\overline{d_{ij}'d_{ij}'}
  \quad [\mathrm{m^2\,s^{-3}}],
\]
o\`u $d_{ij}' = \tfrac{1}{2}(\partial u_i'/\partial x_j + \partial u_j'/\partial x_i)$
est le tenseur des taux de d\'eformation fluctuant.
$\varepsilon$ est le taux auquel $k_t$ est irr\'eversiblement converti en chaleur par la
viscosit\'e mol\'eculaire aux petites \'echelles.
\end{defbox}

\subsection{Viscosit\'e turbulente --- hypoth\`ese de Boussinesq}

Pour fermer les \'equations RANS, Boussinesq (1877) propose de mod\'eliser le tenseur de
Reynolds par analogie avec la loi de Newton pour les fluides visqueux.

Comment mod\'eliser $\overline{u_i'u_j'}$~? L'id\'ee la plus simple --- et historiquement la
premi\`ere --- est de s'inspirer de ce qu'on conna\^it d\'ej\`a~: la loi de Newton pour les
contraintes visqueuses. Au chapitre~2, on a vu que
$\tau_{ij} = \mu(\partial U_i/\partial x_j + \partial U_j/\partial x_i)$~: les contraintes
sont proportionnelles aux gradients de vitesse. Boussinesq (1877) propose exactement la m\^eme
chose pour les contraintes de Reynolds, en rempla\c cant $\mu$ par une viscosit\'e
\emph{turbulente} $\mu_t$ qui n'est plus une propri\'et\'e du fluide mais une propri\'et\'e de
l'\'ecoulement.

Dans un fluide
newtonien, les contraintes visqueuses sont proportionnelles aux taux de d\'eformation
moyens~: $\bar{\tau}_{ij} = 2\mu\bar{D}_{ij}$. Par analogie, on postule que les
fluctuations turbulentes agissent comme une \emph{viscosit\'e suppl\'ementaire} $\mu_t$~:

\begin{importantbox}{Hypoth\`ese de viscosit\'e turbulente --- Boussinesq (1877)}
\[
  -\rho\,\overline{u_i'u_j'} = 2\mu_t\bar{D}_{ij} - \frac{2}{3}\rho k_t\,\delta_{ij}
  = \mu_t\!\left(\frac{\partial\bar{U}_i}{\partial x_j}
  + \frac{\partial\bar{U}_j}{\partial x_i}\right) - \frac{2}{3}\rho k_t\,\delta_{ij},
\]
o\`u $\mu_t(\bm{x},t) = \rho\nu_t$ est la \textbf{viscosit\'e turbulente} ---
propri\'et\'e du champ d'\'ecoulement, non du fluide (contrairement \`a $\mu$).
\end{importantbox}

Le terme isotrope $\frac{2}{3}\rho k_t\delta_{ij}$ est absorb\'e dans une pression modifi\'ee
$\bar{P}^* = \bar{P} + \frac{2}{3}\rho k_t$, et les RANS deviennent~:
\[
  \frac{\partial(\rho\bar{U}_i)}{\partial t}
  + \frac{\partial(\rho\bar{U}_i\bar{U}_j)}{\partial x_j}
  = -\frac{\partial\bar{P}^*}{\partial x_i}
  + \frac{\partial}{\partial x_j}\!\left[2(\mu+\mu_t)\bar{D}_{ij}\right].
\]
Ces \'equations sont formellement identiques aux NS originales, avec $\mu \to \mu + \mu_t$.
Il reste \`a calculer $\mu_t(\bm{x},t)$.

\medskip
\noindent\textbf{Le mod\`ele $k_t$-$\varepsilon$.}
Par analyse dimensionnelle, la viscosit\'e turbulente est le produit d'une \'echelle de vitesse
$u' \sim k_t^{1/2}$ et d'une longueur de m\'elange $\ell \sim k_t^{3/2}/\varepsilon$~:
\[
  \nu_t = C_\mu\frac{k_t^2}{\varepsilon}, \quad C_\mu \approx 0.09.
\]
On adjoint alors deux \'equations de transport suppl\'ementaires pour $k_t$ et $\varepsilon$~:
c'est le \textbf{mod\`ele $k_t$-$\varepsilon$}, le plus r\'epandu en ing\'enierie.

\subsection{Bilan d'\'energie cin\'etique turbulente}

Pour comprendre \emph{d'o\`u vient} $k_t$ et \emph{comment} elle est d\'etruite, on \'etablit
son \'equation d'\'evolution. Comment obtenir cette \'equation~? La d\'emarche est analogue au
bilan d'\'energie cin\'etique du chapitre~2. On soustrait l'\'equation RANS de l'\'equation NS
instantan\'ee pour isoler l'\'equation du mouvement fluctuant. On multiplie ensuite par $u_i'$
(pour obtenir l'\'energie $u_i'^2/2$), on somme sur $i$, et on prend la moyenne d'ensemble.
Le calcul est technique mais le r\'esultat est physiquement limpide~:

\begin{importantbox}{Bilan d'\'energie cin\'etique turbulente}
\[
  \underbrace{\frac{\partial(\rho k_t)}{\partial t}
  + \frac{\partial(\rho k_t \bar{U}_j)}{\partial x_j}}_{\substack{\text{advection} \\
  \text{par }\bar{U}}}
  =
  \underbrace{-\rho\,\overline{u_i'u_j'}\,\frac{\partial\bar{U}_i}{\partial x_j}}_{\text{production }\mathcal{P}}
  \;-\; \underbrace{\rho\varepsilon}_{\text{dissipation}}
  \;+\; \underbrace{\frac{\partial}{\partial x_k}\!\left(
    -\tfrac{1}{2}\rho\,\overline{u_i'u_i'u_k'}
    - \overline{p'u_k'} + \overline{u_i'\tau_{ik}'}
  \right)}_{\text{transport spatial}}
\]
\end{importantbox}

\noindent\textbf{Interpr\'etation du terme de production $\mathcal{P}$.}
C'est le terme central~: il encode le transfert d'\'energie du \emph{cisaillement moyen}
vers les \emph{fluctuations turbulentes}. Consid\'erons un cisaillement
$\bar{U}_1(x_2)$ croissant avec $x_2$~:
\begin{itemize}
  \item Une particule avec $u_2' > 0$ vient d'une r\'egion \`a faible $\bar{U}_1$~: elle
        arrive dans une zone plus rapide avec un d\'eficit $u_1' < 0$.
  \item \`A l'inverse, $u_2' < 0$ am\`ene du fluide rapide vers le bas~: $u_1' > 0$.
\end{itemize}
Dans les deux cas $u_1'u_2' < 0$, donc $\overline{u_1'u_2'} < 0$, et la production vaut~:
\[
  \mathcal{P} \simeq -\rho\,\overline{u_1'u_2'}\,\frac{d\bar{U}_1}{dx_2} > 0.
\]
\textbf{La turbulence se nourrit du cisaillement moyen.} En r\'egime quasi-\'etabli, loin des
parois, production et dissipation s'\'equilibrent localement~:
\[
  \mathcal{P} \approx \rho\varepsilon.
\]

% ------------------------------------------
\section{Cascade d'\'energie et \'echelles de Kolmogorov}
% ------------------------------------------

On a vu que l'\'energie cin\'etique turbulente est produite aux grandes \'echelles (par le
cisaillement moyen) et dissip\'ee par la viscosit\'e. Mais la viscosit\'e agit aux
\emph{petites} \'echelles --- comment l'\'energie passe-t-elle des unes aux autres~?
C'est la question \`a laquelle r\'epond la \textbf{cascade d'\'energie} de Kolmogorov (1941).

\subsection{Deux familles d'\'echelles}

Un \'ecoulement turbulent contient un \emph{continuum} de structures, depuis les plus grandes
jusqu'aux plus petites.

\medskip
\noindent\textbf{Grandes \'echelles $O(L,\,u')$} --- impos\'ees par la g\'eom\'etrie de
l'\'ecoulement (diam\`etre de jet, corde d'aile, \'epaisseur de sillage...). Ce sont les plus
\'energ\'etiques~: elles re\c coivent l'\'energie de la production $\mathcal{P}$ et concentrent
l'essentiel de $k_t$.

\medskip
\noindent\textbf{Petites \'echelles de Kolmogorov $O(l_\eta,\,u_\eta)$} --- \`a tr\`es petite
\'echelle, la viscosit\'e domine les effets inertiels. Le nombre de Reynolds local vaut
$u_\eta l_\eta/\nu = 1$~: c'est exactement la condition \`a laquelle la viscosit\'e peut
dissiper l'\'energie en chaleur.

\begin{defbox}{\'Echelles de Kolmogorov}
La seule combinaison dimensionnellement coh\'erente de $\nu$ et $\varepsilon$ donnant une
longueur, une vitesse et un temps est~:
\[
  l_\eta = \left(\frac{\nu^3}{\varepsilon}\right)^{1/4}, \qquad
  u_\eta = (\nu\varepsilon)^{1/4}, \qquad
  \tau_\eta = \left(\frac{\nu}{\varepsilon}\right)^{1/2}.
\]
On v\'erifie~: $u_\eta l_\eta/\nu = 1$ --- le Reynolds local vaut bien 1.
\end{defbox}

\subsection{La cascade d'\'energie}

La relation entre grandes et petites \'echelles est la \textbf{cascade d'\'energie}, propos\'ee
par Kolmogorov (1941) en s'appuyant sur une id\'ee de Richardson~:

\begin{enumerate}
  \item L'\'energie est \textbf{inject\'ee} aux grandes \'echelles $L$ (via la production $\mathcal{P}$).
  \item Par interactions \textbf{non lin\'eaires} (terme convectif), les grandes structures
        se fragmentent en structures plus petites, qui se fragmentent \`a leur tour~: l'\'energie
        \emph{cascade} vers les petites \'echelles sans \^etre dissip\'ee.
  \item Aux \'echelles $l_\eta$ (o\`u $\mathrm{Re}=1$), la viscosit\'e \textbf{dissipe}
        l'\'energie en chaleur.
\end{enumerate}

\begin{center}
\begin{tikzpicture}[node distance=2.2cm, >=Stealth]
  \node[draw, rounded corners, fill=blue!25, minimum width=2.2cm, minimum height=1cm,
        font=\small] (L) {Grandes \'ech. $L$};
  \node[draw, rounded corners, fill=blue!12, minimum width=1.8cm, minimum height=0.85cm,
        font=\small, right=of L] (I) {Inertielles};
  \node[draw, rounded corners, fill=blue!5, minimum width=1.5cm, minimum height=0.7cm,
        font=\small, right=of I] (K) {$l_\eta$};
  \draw[->, very thick, red!70] (L) -- node[above, font=\small]{cascade} (I);
  \draw[->, very thick, red!70] (I) -- node[above, font=\small]{cascade} (K);
  \draw[->, very thick, orange] (K.south) -- ++(0,-0.9)
        node[below, font=\small]{dissipation $\varepsilon$};
  \draw[->, very thick, green!60!black] (L.north) +(-1.2,0.9)
        -- node[above, font=\small]{prod. $\mathcal{P}$} (L.north);
\end{tikzpicture}
\end{center}

Dans la zone \textbf{inertielle} ($l_\eta \ll \ell \ll L$), l'\'energie est simplement
relay\'ee d'\'echelle en \'echelle \`a taux constant $\varepsilon$, sans production ni dissipation.
L'analyse dimensionnelle de Kolmogorov donne le spectre d'\'energie~:

\begin{importantbox}{Loi de Kolmogorov en $k^{-5/3}$}
Dans la zone inertielle, le spectre d'\'energie ne d\'epend que de $\varepsilon$ et du nombre
d'onde $k$ (inverse d'une longueur). La seule forme dimensionnellement coh\'erente est~:
\[
  E(k) \sim \varepsilon^{2/3}\,k^{-5/3}.
\]
Cette loi universelle (1941) est v\'erifi\'ee dans d'innombrables exp\'eriences~: jets,
couches limites, atmosph\`ere, oc\'ean. La pente $-5/3$ sur un diagramme log-log est la
signature d'une cascade d'\'energie inertielle.
\end{importantbox}

\subsection{Rapport d'\'echelles et co\^ut de la simulation}

La cascade de Kolmogorov a une cons\'equence pratique majeure pour la simulation num\'erique.
Si l'on veut r\'esoudre \emph{toutes} les \'echelles de la turbulence (c'est ce qu'on appelle
une Simulation Num\'erique Directe, ou DNS), il faut un maillage assez fin pour capturer les
plus petites structures. Combien de points faut-il~? La r\'eponse est donn\'ee par le rapport
$L/l_\eta$.

En supposant qu'en r\'egime \'etabli $\mathcal{P} \approx \rho\varepsilon$, et que
$\mathcal{P} \sim \rho u'^3/L$, on obtient $\varepsilon \sim u'^3/L$ et donc~:
\[
  \frac{L}{l_\eta} = \frac{L}{(\nu^3/\varepsilon)^{1/4}}
  \sim \left(\frac{u'L}{\nu}\right)^{3/4} = \mathrm{Re}_L^{3/4}.
\]

\begin{importantbox}{Rapport d'\'echelles --- limite de la DNS}
\[
  \frac{L}{l_\eta} \sim \mathrm{Re}_L^{3/4}
\]
Une simulation num\'erique directe (DNS) r\'esolvant toutes les \'echelles sur un domaine 3D
n\'ecessite un nombre de points en grille $N \sim (L/l_\eta)^3 \sim \mathrm{Re}_L^{9/4}$.
\end{importantbox}

\medskip
\noindent\textbf{Application au Rafale Marine.}
Avec $L \approx 3\,\mathrm{m}$ (corde d'aile), $U_\infty \approx 200\,\mathrm{m\,s^{-1}}$,
$\nu_\mathrm{air} \approx 1.5\times10^{-5}\,\mathrm{m^2\,s^{-1}}$~:
\[
  \mathrm{Re}_L = \frac{U_\infty L}{\nu} \approx 4\times10^7,
  \qquad l_\eta \approx \frac{3}{(4\times10^7)^{3/4}} \approx 10\,\mu\mathrm{m}.
\]
Une DNS 3D de la couche limite du Rafale exigerait
$N \sim \mathrm{Re}_L^{9/4} \approx 10^{16}$ points de grille --- totalement hors de port\'ee.
La simulation industrielle du Rafale repose donc sur les \textbf{\'equations RANS} avec mod\`ele
$k_t$-$\varepsilon$ ou $k_t$-$\omega$ ($\sim 10^7$--$10^8$ points), accept\'e comme
compromis entre pr\'ecision et co\^ut.

% ==========================================
\section*{R\'esum\'e du chapitre}
\addcontentsline{toc}{section}{R\'esum\'e du chapitre}
% ==========================================

\begin{importantbox}{Ce qu'il faut retenir}
\begin{itemize}
  \item La \textbf{turbulence} est un \'ecoulement instationnaire, ap\'eriodique et
        chaotique, multi-\'echelle, qui appara\^it pour $\mathrm{Re}$ suffisamment grand.
        Elle augmente le m\'elange et retarde les d\'ecollements mais accro\^it la tra\^in\'ee
        de frottement.

  \item La \textbf{d\'ecomposition de Reynolds} $U_i = \bar{U}_i + u_i'$ s\'epare champ
        moyen et fluctuations~; r\`egle cl\'e~: $\overline{fg} = \bar{F}\bar{G} +
        \overline{f'g'}$.

  \item Les \textbf{\'equations RANS} font appara\^itre le \textbf{tenseur de Reynolds}
        $-\rho\,\overline{u_i'u_j'}$~: 6 inconnues suppl\'ementaires $\Rightarrow$
        \textbf{probl\`eme de fermeture}.

  \item L'\textbf{hypoth\`ese de Boussinesq}~:
        $-\rho\overline{u_i'u_j'} = 2\mu_t\bar{D}_{ij} - \tfrac{2}{3}\rho k_t\delta_{ij}$,
        avec $\nu_t = C_\mu k_t^2/\varepsilon$ pour le mod\`ele $k_t$-$\varepsilon$.

  \item Le \textbf{bilan de $k_t$}~: production $\mathcal{P}$ (du cisaillement moyen)
        --- dissipation $\rho\varepsilon$ + transport. \'Equilibre local~:
        $\mathcal{P} \approx \rho\varepsilon$.

  \item \textbf{\'Echelles de Kolmogorov}~: $l_\eta = (\nu^3/\varepsilon)^{1/4}$,
        $L/l_\eta \sim \mathrm{Re}_L^{3/4}$. DNS du Rafale~: $\sim 10^{16}$ points,
        impossible $\Rightarrow$ RANS incontournable.

  \item \textbf{Cascade de Kolmogorov}~: injection \`a $L$, transfert inertiel,
        dissipation \`a $l_\eta$. Spectre inertiel~: $E(k) \sim \varepsilon^{2/3} k^{-5/3}$.
\end{itemize}
\end{importantbox}

% ==========================================
\chapter{\'Energie, thermodynamique et \'ecoulements compressibles}
% ==========================================

Jusqu'ici nous avons impos\'e $\nabla \cdot \bm{U} = 0$~: la masse volumique
$\rho$ \'etait constante le long des lignes de courant, et la pression n'avait qu'un r\^ole
m\'ecanique. Ce cadre est mis en d\'efaut d\`es que la vitesse devient une fraction
significative de la vitesse du son.

Pour le Rafale Marine, dont la vitesse de croisi\`ere est $U_\infty \simeq 200\,\mathrm{m/s}$,
le nombre de Mach vaut $M_\infty = U_\infty/c \simeq 0{,}59$ (avec $c \simeq 340\,\mathrm{m/s}$
\`a altitude standard). Au-del\`a de $M \simeq 0{,}3$ l'hypoth\`ese incompressible engendre
une erreur sup\'erieure \`a 5\,\%~: il faut donc enrichir notre mod\`ele. Plus radicalement,
\`a l'int\'erieur du r\'eacteur M88-2 le fluide est chauff\'e \`a des temp\'eratures
sup\'erieures \`a $1\,500\,\mathrm{K}$ et acc\'el\'er\'e dans la tuy\`ere convergente-divergente
jusqu'\`a des vitesses supersoniques~: les ph\'enom\`enes thermiques et \'energ\'etiques sont
alors au c\oe{}ur du probl\`eme.

Ce chapitre construit pas \`a pas le cadre compressible en r\'epondant \`a deux questions~:
\begin{enumerate}
  \item Quelle \'equation gouverne l'\'echange d'\'energie entre les modes cin\'etique et
        interne du fluide~?
  \item Quelles lois thermodynamiques relient entre elles les variables d'\'etat
        ($\rho$, $p$, $T$, $e$, $s$, $h$)~?
\end{enumerate}
La r\'eponse fournit le syst\`eme complet d'\'equations de Navier-Stokes compressibles et
ouvre la voie \`a l'\'etude des ondes de choc.

\begin{defbox}{Fil conducteur}
Le Rafale vole en r\'egime subsonique \`a $M \approx 0{,}6$ en croisi\`ere, mais traverse
le r\'egime transsonique ($M \approx 1$) en acc\'el\'eration et atteint $M = 1{,}8$ en
pointe. Son r\'eacteur M88-2 fonctionne enti\`erement en r\'egime compressible (de la
prise d'air supersonique \`a la tuy\`ere convergente-divergente). Comprendre ces r\'egimes
exige de coupler la m\'ecanique \`a la thermodynamique~: c'est l'objet de ce dernier chapitre.
\end{defbox}

% ------------------------------------------
\section{Conservation de l'\'energie totale}
% ------------------------------------------

\subsection{Tenseur des contraintes visqueuses en \'ecoulement compressible}

Au chapitre~2, nous avons \'ecrit $\ten{\sigma} = -p\,\ten{I} + 2\mu\,\ten{D}$ pour un fluide newtonien incompressible. Ce tenseur ne contenait aucun terme li\'e \`a la dilatation volumique --- et pour cause~: en incompressible, $\nabla\cdot\bm{U} = 0$ par d\'efinition. Mais d\`es que le fluide se comprime ou se d\'etend (une particule d'air qui entre dans la prise d'air du M88-2 est comprim\'ee d'un facteur~30), le taux de dilatation $\nabla\cdot\bm{U} \neq 0$ et il faut enrichir la loi constitutive.

En \'ecoulement incompressible le tenseur des contraintes s'\'ecrivait
$\ten{\sigma} = -p\,\ten{I} + 2\mu\,\ten{D}$,
o\`u le terme de Newton \'etait enti\`erement d\'evi\'e. D\`es que $\nabla \cdot \bm{U} \neq 0$,
la d\'eformation \emph{sph\'erique} (compression/dilatation uniforme) produit \'egalement
des contraintes normales~: la loi de Newton doit \^etre compl\'et\'ee.

La forme g\'en\'erale du tenseur visqueux newtonien est~:

\begin{importantbox}{Tenseur des contraintes --- cas compressible}
\[
  \ten{\sigma} = -p\,\ten{I} + \ten{\tau},
  \qquad
  \ten{\tau} = 2\mu\,\ten{D} + \lambda\,(\nabla \cdot \bm{U})\,\ten{I}
\]
o\`u $\lambda$ est la \textbf{deuxi\`eme viscosit\'e} (ou viscosit\'e de volume).
\end{importantbox}

Le terme suppl\'ementaire $\lambda(\nabla \cdot \bm{U})\ten{I}$ est isotrope et
proportionnel au taux de dilatation volumique. On peut v\'erifier que la dissipation
visqueuse vaut alors~:
\[
  \ten{\tau} : \ten{D}
  = 2\mu\,\ten{D} : \ten{D} + \lambda\,(\nabla \cdot \bm{U})^2 \geq 0.
\]
L'hypoth\`ese de Stokes $\lambda = -\tfrac{2}{3}\mu$ (valable pour les gaz monoatomiques)
ram\`ene la pression thermodynamique \`a la valeur moyenne des contraintes normales.

\subsection{\'Energie interne~: un nouveau degr\'e de libert\'e}

Jusqu'au chapitre~2, la seule forme d'\'energie qui nous int\'eressait \'etait l'\'energie cin\'etique $\tfrac{1}{2}\rho U^2$~: c'est elle qui appara\^it dans le bilan d'\'energie m\'ecanique et dans Bernoulli. Mais r\'efl\'echissons~: quand on comprime un gaz, on lui fournit du travail m\^eme si sa vitesse macroscopique ne change pas. O\`u va cette \'energie~? Elle va dans l'agitation thermique des mol\'ecules --- c'est-\`a-dire dans l'\emph{\'energie interne}. De m\^eme, quand la couche limite du Rafale freine l'air par frottement visqueux, l'\'energie cin\'etique perdue ne dispara\^it pas~: elle est convertie en chaleur (\'energie interne). Pour d\'ecrire ces conversions, il faut introduire un nouveau degr\'e de libert\'e.

Les mol\'ecules du fluide poss\`edent
\'egalement de l'\'energie li\'ee \`a leur agitation thermique et aux interactions
intermol\'eculaires. Cette contribution microscopique est l'\textbf{\'energie interne
sp\'ecifique} $e(\bm{x},t)$ (en $\mathrm{J\,kg^{-1}}$).

L'\'energie totale sp\'ecifique est donc~:
\[
  e_t \equiv e + \frac{1}{2}U^2,
  \qquad
  E_t = \int_{\mathcal{D}} \rho\, e_t\, \mathrm{d}v.
\]

\subsection{Premier principe de la thermodynamique --- bilan global}

Le premier principe stipule que la variation d'\'energie totale d'un domaine
mat\'eriel est \'egale \`a la puissance des forces ext\'erieures augment\'ee du taux
d'apport thermique~:

\begin{importantbox}{Premier principe pour le fluide}
\[
  \frac{\mathrm{d}E_t}{\mathrm{d}t} = P_{\mathrm{ext}} + Q
\]
avec $P_{\mathrm{ext}}$ puissance des forces de volume et de surface,
$Q$ apport de chaleur par unit\'e de temps (en W).
\end{importantbox}

En soustrayant le bilan de l'\'energie cin\'etique (\'etabli au chapitre~2 sous forme
int\'egrale sur un domaine mat\'eriel $\mathcal{D}$)~:
\[
  \frac{\mathrm{d}}{\mathrm{d}t}\int_{\mathcal{D}} \rho\,\frac{U^2}{2}\,\mathrm{d}v
  = P_{\mathrm{ext}} + P_{\mathrm{int}},
\]
on obtient l'\'equation de l'\'energie \emph{interne}~:
\[
  \frac{\mathrm{d}}{\mathrm{d}t}\int_{\mathcal{D}} \rho\,e\,\mathrm{d}v
  = Q - P_{\mathrm{int}},
\]
o\`u $-P_{\mathrm{int}} = \int_{\mathcal{D}} \ten{\sigma} : \ten{D}\,\mathrm{d}v
= -\int_{\mathcal{D}} p\nabla \cdot \bm{U}\,\mathrm{d}v
+ \int_{\mathcal{D}} \ten{\tau} : \ten{D}\,\mathrm{d}v$
regroupe le travail r\'eversible de compression et la dissipation visqueuse irr\'eversible.

Le sch\'ema \'energ\'etique est simple~: l'\'energie cin\'etique, aliment\'ee par
$P_{\mathrm{ext}}$, se convertit irr\'eversiblement vers l'\'energie interne via la
dissipation visqueuse $\ten{\tau}:\ten{D}$, et se recharge r\'eversiblement
par compression $-p\nabla\cdot\bm{U}$ (qui change de signe selon que le fluide
est comprim\'e ou d\'etendu).

\subsection{Apport de chaleur --- flux et vecteur flux thermique}

L'apport thermique $Q$ provient de deux sources~:
\[
  Q \equiv Q_s + Q_v,
  \quad
  Q_s = -\oint_{\mathcal{S}} \bm{q} \cdot \bm{n}\,\mathrm{d}s,
  \quad
  Q_v = \int_{\mathcal{D}} \rho\,q_\star\,\mathrm{d}v,
\]
o\`u $\bm{q}(\bm{x},t)$ est le \textbf{vecteur densit\'e de flux de chaleur}
(en $\mathrm{W\,m^{-2}}$) et $q_\star$ (en $\mathrm{W\,kg^{-1}}$) un apport volumique
(rayonnement, combustion). La loi constitutive reliant $\bm{q}$ \`a l'\'etat local est la
\textbf{loi de Fourier} (1822)~:

\begin{defbox}{Loi de Fourier}
\[
  \bm{q} = -k\,\nabla T,
\]
o\`u $k \geq 0$ est la \textbf{conductivit\'e thermique} (en $\mathrm{W\,m^{-1}\,K^{-1}}$).
Le signe moins traduit le fait que la chaleur se propage du chaud vers le froid.
\end{defbox}

\subsection{Formulation locale de la conservation de l'\'energie}

En appliquant le th\'eor\`eme de Reynolds et en passant \`a la formulation locale (valable
en tout point r\'egulier), on obtient trois \'equations \'equivalentes mais compl\'ementaires~:

\begin{importantbox}{\'Equations locales de l'\'energie}
\[
  \begin{cases}
  \displaystyle\rho\,\DD{e_t} = \nabla\cdot(\bm{U}\cdot\ten{\sigma})
                                + \rho\bm{U}\cdot\bm{g}
                                - \nabla\cdot\bm{q} + \rho q_\star
  \\[8pt]
  \displaystyle\rho\,\frac{D}{Dt}\!\left(\frac{U^2}{2}\right)
      = \nabla\cdot(\bm{U}\cdot\ten{\sigma})
        + \rho\bm{U}\cdot\bm{g}
        + p\nabla\cdot\bm{U} - \ten{\tau}:\ten{D}
  \\[8pt]
  \displaystyle\rho\,\DD{e}
      = \underbrace{-p\,\nabla\cdot\bm{U}}_{\text{compression}}
        + \underbrace{\ten{\tau}:\ten{D}}_{\text{dissipation}}
        + \underbrace{-\nabla\cdot\bm{q} + \rho q_\star}_{\text{chaleur}}
  \end{cases}
\]
\end{importantbox}

La troisi\`eme ligne est la plus \'eclairante~: l'\'energie interne d'une particule fluide
augmente quand elle est comprim\'ee, quand la dissipation visqueuse la r\'echauffe, ou
quand de la chaleur lui est apport\'ee.

% ------------------------------------------
\section{Rappels de thermodynamique d'\'equilibre}
% ------------------------------------------

L'\'equation de l'\'energie interne fait intervenir $e$, $p$, $T$ --- mais ces trois grandeurs ne sont pas ind\'ependantes. Un gaz \`a l'\'equilibre est enti\`erement d\'ecrit par deux variables seulement (par exemple $\rho$ et $T$), et toutes les autres s'en d\'eduisent via des \emph{\'equations d'\'etat}. Pour fermer notre syst\`eme, il faut donc relier ces variables entre elles~: c'est le r\^ole de la thermodynamique. On rappelle ici les r\'esultats essentiels, en les motivant physiquement.

\subsection{Variables d'\'etat}

Un fluide simple est enti\`erement d\'ecrit par \emph{deux} variables ind\'ependantes.
En m\'ecanique des fluides, on choisit g\'en\'eralement $\rho$ et $e$, d'o\`u l'on
d\'eduit toutes les autres~:

\begin{center}
\begin{tabular}{cll}
\toprule
Variable & Nom & Unit\'e \\
\midrule
$T$ & Temp\'erature absolue & K \\
$p$ & Pression thermodynamique & Pa \\
$s$ & Entropie sp\'ecifique & $\mathrm{J\,kg^{-1}\,K^{-1}}$ \\
$h \equiv e + p/\rho$ & Enthalpie sp\'ecifique & $\mathrm{J\,kg^{-1}}$ \\
\bottomrule
\end{tabular}
\end{center}

En \'ecoulement, chaque particule fluide est suppos\'ee \`a l'\'equilibre \`a son \'echelle
microscopique (hypoth\`ese d'\'equilibre local), de sorte que $\rho(\bm{x},t)$ et
$e(\bm{x},t)$ d\'eterminent localement $T$, $p$, $s$, $h$.

\subsection{Relation de Gibbs et second principe}

Pour un syst\`eme ferm\'e, le premier principe s'\'ecrit
$\delta q = \mathrm{d}e + p\,\mathrm{d}(1/\rho)$.
Le second principe (Clausius) introduit l'entropie via
$\delta q_{\mathrm{rev}} = T\,\mathrm{d}s$. En combinant~:

\begin{importantbox}{Relation de Gibbs}
\[
  T\,\mathrm{d}s = \mathrm{d}e + p\,\mathrm{d}\!\left(\frac{1}{\rho}\right)
\]
\end{importantbox}

Cette identit\'e fondamentale relie des variations de variables d'\'etat, ind\'ependamment
du chemin suivi. Pour un processus quelconque, l'entropie augmente en pr\'esence
d'irr\'eversibilit\'es (viscosit\'e, conduction, chocs).

Dit autrement~: si on comprime une particule fluide (d$(1/\rho) < 0$, le volume sp\'ecifique diminue), l'\'energie interne augmente --- sauf si de la chaleur s'\'echappe simultan\'ement ($T\,\mathrm{d}s < 0$ possible si le processus est non-adiabatique). C'est exactement ce qui se passe dans le compresseur du M88-2~: l'air est comprim\'e (son volume diminue) et sa temp\'erature monte de $\sim 288\,\mathrm{K}$ \`a $\sim 800\,\mathrm{K}$ sur 6 \'etages.

\begin{defbox}{Transformations particuli\`eres}
\begin{itemize}
  \item \textbf{Adiabatique~:} $\delta q = 0 \Rightarrow \mathrm{d}s \geq 0$
  \item \textbf{Isentropique (adiabatique + r\'eversible)~:} $\mathrm{d}s = 0$
\end{itemize}
\end{defbox}

\subsection{Chaleurs sp\'ecifiques}

La r\'eponse thermique du fluide d\'epend des conditions~:
\[
  c_v = \left.\pd{e}{T}\right|_\rho \quad \text{(volume constant)},
  \qquad
  c_p = \left.\pd{h}{T}\right|_p \quad \text{(pression constante)}.
\]
On a toujours $c_p > c_v$ car une partie de la chaleur fournie \`a pression constante
sert \`a faire le travail de dilatation.

\subsection{Gaz parfait}

Pour les gaz peu denses (air \`a pression atmosph\'erique, gaz de combustion dans le
M88-2), le mod\`ele du \textbf{gaz parfait} s'applique~:

\begin{importantbox}{Gaz parfait --- loi d'\'etat et cons\'equences}
\[
  p = \rho r T,
  \qquad r = \frac{R}{M_{\mathrm{mol}}} \quad
  (R \simeq 8{,}314\,\mathrm{J\,mol^{-1}\,K^{-1}})
\]
Relation de Mayer~: $c_p - c_v = r$. \quad
Rapport~: $\gamma = c_p/c_v > 1$. \\
Entropie~: $s = c_v \ln(p/\rho^\gamma) + \mathrm{cst}$. \quad
Isentropique~: $p/\rho^\gamma = \mathrm{cst}$.
\end{importantbox}

Pour l'air \`a conditions standard~: $r = 287\,\mathrm{J\,kg^{-1}\,K^{-1}}$,
$\gamma = 1{,}4$, $c_p = 1005\,\mathrm{J\,kg^{-1}\,K^{-1}}$,
$c_v = 718\,\mathrm{J\,kg^{-1}\,K^{-1}}$.

% ------------------------------------------
\section{\'Equations de la thermodynamique du mouvement}
% ------------------------------------------

On dispose maintenant de tous les ingr\'edients pour \'ecrire le syst\`eme complet
gouvernant un fluide compressible.

\subsection{Lois constitutives}

Deux lois ferment le syst\`eme~:

\begin{defbox}{Lois constitutives --- fluide newtonien conducteur de chaleur}
\[
  \ten{\tau} = 2\mu\,\ten{D} + \lambda(\nabla\cdot\bm{U})\,\ten{I},
  \qquad
  \bm{q} = -k\nabla T.
\]
Les coefficients $\mu$, $\lambda$, $k$ d\'ependent en g\'en\'eral de l'\'etat local~:
$\mu = \mu(T)$ (loi de Sutherland, vue au chapitre~3), etc.
\end{defbox}

\subsection{Syst\`eme complet de Navier-Stokes compressible}

En rassemblant la conservation de la masse, de la quantit\'e de mouvement et de
l'\'energie interne avec les lois constitutives et une \'equation d'\'etat~:

\begin{importantbox}{Navier-Stokes compressible}
\[
  \begin{cases}
    \displaystyle\DD{\rho} + \rho\,\nabla\cdot\bm{U} = 0
    & \text{(masse)} \\[6pt]
    \displaystyle\rho\,\DD{\bm{U}} = -\nabla p + \nabla\cdot\ten{\tau} + \rho\bm{g}
    & \text{(quantit\'e de mouvement)} \\[6pt]
    \displaystyle\rho\,\DD{e} = -p\,\nabla\cdot\bm{U} + \ten{\tau}:\ten{D}
                                - \nabla\cdot\bm{q} + \rho q_\star
    & \text{(\'energie interne)}
  \end{cases}
\]
compl\'et\'e par $p = p(\rho,T)$, $e = e(\rho,T)$, $\mu,\lambda,k = f(\rho,T)$.
\end{importantbox}

Ce syst\`eme compte $5$ \'equations scalaires pour $5$ inconnues primaires ($\rho$, trois
composantes de $\bm{U}$, $e$ ou $T$)~; $p$ est \'eliminable via la loi d'\'etat.

\subsection{Conditions aux limites thermiques}

\`A une paroi solide, en plus de la condition de non-glissement
$\bm{U} = \bm{U}_{\mathrm{paroi}}$, il faut prescrire la condition thermique~:
\begin{itemize}
  \item \textbf{Paroi isotherme~:} $T = T_{\mathrm{paroi}}$.
  \item \textbf{Paroi adiabatique~:} $\bm{q}\cdot\bm{n} = 0$,
        soit $\partial T/\partial n = 0$.
\end{itemize}

% ------------------------------------------
\section{Entropie et second principe}
% ------------------------------------------

\subsection{\'Equation de transport de l'entropie}

On a vu que la relation de Gibbs relie $\mathrm{d}s$, $\mathrm{d}e$ et $\mathrm{d}(1/\rho)$ pour un syst\`eme \`a l'\'equilibre. Mais une particule fluide en mouvement n'est pas \`a l'\'equilibre global~: elle subit de la dissipation visqueuse et de la conduction thermique, deux processus irr\'eversibles. La question est~: \`a quel taux son entropie augmente-t-elle~? Pour r\'epondre, on combine la relation de Gibbs avec les bilans de masse et d'\'energie interne d\'ej\`a \'etablis.

Appliquons la relation de Gibbs \`a une particule fluide~:
$T\,\DD{s} = \DD{e} + \frac{p}{\rho^2}\DD{\rho}$.
La conservation de la masse donne $\frac{1}{\rho}\DD{\rho} = -\nabla\cdot\bm{U}$,
et en substituant l'\'equation de l'\'energie interne~:

\begin{importantbox}{\'Equation de transport de l'entropie}
\[
  \rho\,\DD{s}
  = \underbrace{\frac{\ten{\tau}:\ten{D}}{T}
                - \frac{\bm{q}\cdot\nabla T}{T^2}}_{\text{production d'entropie }\geq 0}
    + \nabla\cdot\!\left(\frac{\bm{q}}{T}\right)
    + \frac{\rho q_\star}{T}
\]
\end{importantbox}

Le terme de production est la somme de deux contributions toujours positives~:
\begin{itemize}
  \item $\ten{\tau}:\ten{D}/T \geq 0$~: dissipation m\'ecanique par viscosit\'e.
  \item $-\bm{q}\cdot\nabla T/T^2 = k(\nabla T)^2/T^2 \geq 0$~: dissipation
        thermique par conduction.
\end{itemize}

\subsection{Contraintes impos\'ees par le second principe}

Pour un syst\`eme isol\'e, la production totale d'entropie doit \^etre positive. En
analysant cette condition pour tous les \'etats possibles, on d\'emontre~:
\[
  \mu \geq 0, \quad \lambda \geq -\tfrac{2}{3}\mu, \quad k \geq 0.
\]
Ces in\'egalit\'es ne sont donc pas des choix mais des \emph{cons\'equences} du second
principe de la thermodynamique.

\subsection{\'Equation de transport de la temp\'erature}

\`A partir de $s = s(p,T)$ et des relations de Maxwell, on d\'erive une \'equation de
transport pour $T$~:

\begin{importantbox}{\'Equation de la temp\'erature (cas g\'en\'eral)}
\[
  \rho c_p\,\DD{T}
  = \ten{\tau}:\ten{D}
    + \beta T\,\DD{p}
    - \nabla\cdot\bm{q}
    + \rho q_\star,
\]
o\`u $\beta$ est le coefficient de dilatation thermique ($\beta = 1/T$ pour un gaz parfait).
\end{importantbox}

Pour un liquide incompressible ($\nabla\cdot\bm{U}=0$, $\beta \simeq 0$)~:
$\rho c\,\DD{T} = k\nabla^2 T + \ten{\tau}:\ten{D} + \rho q_\star$,
qui est l'\'equation de la chaleur convect\'ee.

% ------------------------------------------
\section{Effets de compressibilit\'e et nombre de Mach}
% ------------------------------------------

On dispose d\'esormais du syst\`eme complet d'\'equations. Mais \`a quoi ressemblent concr\`etement les \'ecoulements compressibles~? Comment quantifier le moment o\`u l'hypoth\`ese incompressible ne tient plus~? Et que se passe-t-il quand un avion d\'epasse la vitesse du son~? Pour r\'epondre \`a ces questions, il faut d'abord comprendre ce qu'est la vitesse du son et pourquoi elle joue un r\^ole si particulier.

\subsection{Vitesse du son --- naissance du nombre de Mach}

Imaginons qu'on claque des mains dans l'air. L'information ``il y a eu une surpression ici'' se propage sous forme d'onde sonore. \`A quelle vitesse~? C'est une question fondamentale, car cette vitesse fixe la fronti\`ere entre les r\'egimes subsonique et supersonique.

Physiquement, la vitesse du son est la vitesse \`a laquelle les mol\'ecules transmettent une perturbation de pression de proche en proche. Elle d\'epend de la \emph{compressibilit\'e isentropique} du milieu~: plus un gaz est ``raide'' (difficile \`a comprimer), plus le son s'y propage vite. Formellement, une analyse de perturbation lin\'eaire des \'equations compressibles montre que les ondes acoustiques se propagent \`a~:

\begin{defbox}{Vitesse du son}
\[
  c^2 = \left.\pd{p}{\rho}\right|_s = \frac{\gamma p}{\rho} = \gamma r T
  \qquad \text{(gaz parfait)}
\]
\end{defbox}

Cette d\'eriv\'ee est prise \`a entropie constante car la propagation acoustique est
adiabatique r\'eversible. Pour l'air \`a $T = 288\,\mathrm{K}$~:
$c = \sqrt{1{,}4 \times 287 \times 288} \approx 340\,\mathrm{m/s}$.

Le \textbf{nombre de Mach} compare la vitesse locale du fluide \`a la vitesse du son~:
$M = U/c = U/\sqrt{\gamma r T}$.

\begin{center}
\begin{tikzpicture}[font=\small]
  \draw[->] (0,0) -- (9,0) node[right] {$M$};
  \fill[blue!15]  (0,0) rectangle (1.5,0.5);
  \fill[green!20] (1.5,0) rectangle (4.2,0.5);
  \fill[orange!30](4.2,0) rectangle (6.0,0.5);
  \fill[red!25]   (6.0,0) rectangle (8.5,0.5);
  \node at (0.75,0.25) {incomp.};
  \node at (2.85,0.25) {subsonique};
  \node at (5.10,0.25) {transsonique};
  \node at (7.25,0.25) {supersonique};
  \foreach \x/\lab in {1.5/0.3, 4.2/0.8, 6.0/1.2}{
    \draw (\x,-0.15)--(\x,0.65) node[above,font=\tiny] {$M=\lab$};
  }
  \draw[blue!70!black, thick, ->] (2.95,-0.5) -- (2.95,-0.05);
  \node[blue!70!black, below] at (2.95,-0.5) {Rafale $M_\infty\!\simeq\!0{,}59$};
\end{tikzpicture}
\end{center}

\subsection{Mod\`ele 1D compressible --- enthalpie totale conserv\'ee}

Consid\'erons un \'ecoulement 1D stationnaire, non visqueux, sans gravit\'e ni apport
thermique, dans un conduit de section variable (mod\`ele de la tuy\`ere du M88-2).
La conservation de l'enthalpie totale g\'en\'eralise Bernoulli au cas compressible~:

L'\textbf{enthalpie totale} est d\'efinie par~:
\[
  h_t = h + \frac{U^2}{2} = e + \frac{p}{\rho} + \frac{U^2}{2}.
\]
Le long d'une ligne de courant en \'ecoulement adiabatique non visqueux, $h_t$ est
conserv\'ee~: la pression statique et l'enthalpie interne \'evoluent ensemble quand la
vitesse change, mais leur somme reste constante.

\subsection{Tube de Pitot compressible --- temp\'erature et pression de stagnation}

Le tube de Pitot (chapitre~1) mesure la pression d'arr\^et en amenant le fluide au
repos isentropiquement. La conservation de $h_t$ entre l'\'ecoulement libre et le
point d'arr\^et ($U_0 = 0$) donne~:
\[
  c_p T_1 + \frac{U_1^2}{2} = c_p T_0
  \quad\implies\quad
  T_0 = T_1\!\left(1 + \frac{\gamma-1}{2}M_1^2\right).
\]
Pour le Rafale \`a $M = 0{,}59$, $T_\infty = 288\,\mathrm{K}$~:
$T_0 \simeq 288\,(1 + 0{,}2 \times 0{,}59^2) \simeq 308\,\mathrm{K}$. L'\'echauffement
a\'erodynamique est d\'ej\`a de $20\,\mathrm{K}$.

La relation isentropique $p/\rho^\gamma = \mathrm{cst}$ combin\'ee \`a la loi des gaz
parfaits donne la pression de stagnation~:

\begin{importantbox}{Pression de stagnation --- tube de Pitot compressible}
\[
  \frac{p_0}{p_1} = \left(1 + \frac{\gamma-1}{2}M_1^2\right)^{\!\gamma/(\gamma-1)}
\]
En d\'eveloppant au premier ordre en $M^2$, on retrouve la formule de Bernoulli
incompressible $p_0 - p_1 \simeq \tfrac{1}{2}\rho U^2$. La d\'eviation atteint 5\,\%
\`a $M \simeq 0{,}3$~: c'est le crit\`ere classique de validit\'e du mod\`ele incompressible.
\end{importantbox}

\subsection{\'Ecoulements supersoniques et onde de choc}

Quand $M > 1$, la physique change radicalement. Une perturbation ne peut plus se propager
en amont~: le fluide ne ``sait pas'' qu'il va rencontrer un obstacle. L'adaptation se fait
brutalement via une \textbf{onde de choc}, surface de discontinuit\'e \`a travers laquelle~:
\begin{itemize}
  \item la vitesse chute soudainement~;
  \item la pression, la temp\'erature et la masse volumique sautent~;
  \item l'entropie \emph{augmente}~: le choc est irr\'eversible.
\end{itemize}

\begin{center}
\begin{tikzpicture}[font=\small, scale=1.0]
  \fill[gray!20] (0,-0.4) -- (4,0) -- (0,0.4) -- cycle;
  \draw[thick] (0,-0.4) -- (4,0) -- (0,0.4);
  \draw[red!80!black, very thick] (0,-1.5) -- (0,1.5) node[above, red!80!black] {choc droit};
  \foreach \y in {-1.2,-0.7,-0.2,0.3,0.8,1.2}{
    \draw[->, blue!70] (-2.2,\y) -- (-0.2,\y);
  }
  \node[blue!70] at (-1.5,-1.7) {$M_1 > 1$, $p_1$, $T_1$};
  \foreach \y in {-1.2,-0.7,0.8,1.2}{
    \draw[->, orange!80!black] (0.2,\y) -- (1.5,\y);
  }
  \node[orange!80!black] at (2.5,-1.7) {$M_2 < 1$, $p_2 > p_1$, $T_2 > T_1$};
\end{tikzpicture}
\end{center}

Les ondes de choc sont visibles autour du Rafale lors du passage transsonique et \`a
l'entr\'ee du r\'eacteur M88-2 o\`u un choc est utilis\'e pour d\'ec\'el\'erer le flux
avant la chambre de combustion. Les relations de Rankine-Hugoniot relient les grandeurs
de part et d'autre du choc via les bilans de masse, de quantit\'e de mouvement et
d'enthalpie totale.

% ==========================================
\section*{R\'esum\'e du chapitre}
\addcontentsline{toc}{section}{R\'esum\'e du chapitre}
% ==========================================

\begin{importantbox}{Ce qu'il faut retenir}
\begin{enumerate}
  \item En \'ecoulement compressible ($\nabla\cdot\bm{U}\neq 0$), la dynamique
        et la thermique sont coupl\'ees~: il faut r\'esoudre simultan\'ement masse,
        quantit\'e de mouvement et \'energie.
  \item La loi de Newton compressible introduit la deuxi\`eme viscosit\'e $\lambda$~;
        la loi de Fourier $\bm{q}=-k\nabla T$ mod\'elise la conduction.
  \item La relation de Gibbs $T\,\mathrm{d}s = \mathrm{d}e + p\,\mathrm{d}(1/\rho)$
        est l'identit\'e thermodynamique fondamentale.
  \item La vitesse du son $c = \sqrt{\gamma r T}$ et le nombre de Mach $M = U/c$
        caract\'erisent l'importance des effets compressibles~; le crit\`ere $M < 0{,}3$
        valide le mod\`ele incompressible.
  \item L'enthalpie totale $h_t = h + U^2/2$ est conserv\'ee le long d'une ligne de
        courant en \'ecoulement adiabatique non visqueux.
  \item Le second principe impose $\mu \geq 0$, $\lambda \geq -2\mu/3$, $k \geq 0$.
\end{enumerate}
\end{importantbox}

\end{document}